% Leçon 28 — Expression orale : gestion du budget — Chinois Tle A — Cours signé OnBuch+
% Programme officiel MINESEC. Compiler avec : tectonic ZH28-oral-gestion-du-budget.tex
\newcommand{\DOCMATIERE}{Chinois}
\newcommand{\DOCNIVEAU}{Tle A}
\newcommand{\DOCMODULE}{Module 4 : Activités économiques et monde du travail}
\newcommand{\DOCLECON}{ZH28}
\newcommand{\DOCTITRE}{Leçon 28 — Expression orale : gestion du budget}
% OnBuch+ — préambule « cours signés » ENRICHI (compilé avec tectonic / XeLaTeX)
% Système visuel « Pop » d'OnBuch+ : fond crème (jamais blanc), bordures encre
% épaisses, ombres « dures » décalées, titres Archivo Black en capitales.
% Version MATHÉMATIQUES : graphes (pgfplots/tikz), boîtes pédagogiques
% (définition, propriété, méthode, attention, le savais-tu, exercice résolu,
% exercice, TP, bilan), page de garde et signature OnBuch+ ancrée en bas.
%
% Macros attendues dans main.tex AVANT \input{preamble} :
%   \DOCMATIERE  \DOCNIVEAU  \DOCMODULE  \DOCLECON (numéro)  \DOCTITRE
\documentclass[11pt]{article}

\usepackage{fontspec}
\usepackage[french]{babel} % français = langue principale ; le chinois via \zh{..} / textebox (xeCJK)
\usepackage{xeCJK}
\usepackage{pifont} % \ding{51} \ding{55} utilisés par les modèles
\usepackage[a4paper,margin=2cm,top=3.4cm,bottom=2.8cm,headheight=1.4cm,footskip=1.3cm]{geometry}
\usepackage{amsmath,amssymb,amsfonts}
\usepackage{mathtools} % \xmapsto, \coloneqq... souvent utilisés par les modèles
\usepackage{cancel} % simplifications barrées (bilans de forces)
% \abs{z} (module, valeur absolue) et \norm{u} (norme), utilisés par les modèles
\providecommand{\abs}{}\let\abs\relax\DeclarePairedDelimiter{\abs}{\lvert}{\rvert}
\providecommand{\norm}{}\let\norm\relax\DeclarePairedDelimiter{\norm}{\lVert}{\rVert}
\usepackage[table]{xcolor} % [table] : fournit aussi \rowcolors (alternance de lignes)
\usepackage{pagecolor}
\usepackage{tikz}
\usetikzlibrary{arrows.meta,positioning,calc,shapes.geometric,shapes.misc,shapes.arrows,decorations.pathmorphing,decorations.pathreplacing,decorations.markings,patterns,shadows,mindmap,trees,fit,backgrounds,matrix,intersections,angles,quotes}
\usepackage{pgfplots}
\usepackage[version=4]{mhchem} % équations nucléaires \ce{^{235}_{92}U}
\usepackage[european,siunitx]{circuitikz} % schémas électriques
\pgfplotsset{compat=1.18}
\usepgfplotslibrary{fillbetween,groupplots} % aires entre courbes (calcul intégral)
\usepackage{enumitem}
\usepackage{fancyhdr}
% [most] charge toutes les bibliothèques tcolorbox (listings, minted, raster,
% documentation...) et gonfle inutilement la mémoire TeX sur un document long
% (~300 boîtes) : on ne charge que ce qui est réellement utilisé.
\usepackage[skins,breakable]{tcolorbox}
\usepackage{graphicx}
\usepackage{colortbl}
\usepackage{booktabs}
\usepackage{tabularx}
\usepackage{diagbox} % case d'en-tête diagonale des tableaux à double entrée
% Colonnes extensibles centrée / gauche / droite (C, L, R), souvent utilisées par les modèles
\newcolumntype{C}{>{\centering\arraybackslash}X}
\newcolumntype{L}{>{\raggedright\arraybackslash}X}
\newcolumntype{R}{>{\raggedleft\arraybackslash}X}
\usepackage{array}
\usepackage{multicol}
\usepackage{multirow}
\usepackage{siunitx}
% parse-numbers=false : accepte aussi \SI{2,5 \times 10^{-3}}{...}
\sisetup{locale=FR,output-decimal-marker={,},group-minimum-digits=4,parse-numbers=false}
\DeclareSIUnit\ohm{\ensuremath{\symup{\Omega}}} % U+2126 absent de la police
\DeclareSIUnit\division{div} % réglages d'oscilloscope (ms/div, V/div)
\DeclareSIUnit\tr{tr} % tours (vitesse de rotation, ex. \SI{1500}{\tr\per\minute})
\DeclareSIUnit\molar{M} % concentration molaire (ex. \SI{3}{\molar}, \SI{1}{\micro\molar})
\DeclareSIUnit\an{an} % année (le modèle confond parfois avec \year, un compteur TeX) — ex. \SI{5}{\kilo\meter\per\an}
\DeclareSIUnit\FCFA{FCFA} % franc CFA (ex. \SI{500}{\FCFA\per\kilo\gram})
\DeclareSIUnit\degreeBrix{\textdegree Brix} % teneur en sucre d'un sirop/jus (réfractométrie)
\DeclareSIUnit\month{mois} % le modèle utilise parfois \month, un compteur TeX, comme unité de durée
\usepackage{qrcode}
% \degree : commande fréquemment utilisée par les modèles pour un angle ou une
% température en dehors de \SI{}{} (non fournie par les paquets chargés).
\providecommand{\degree}{^{\circ}}
% \qed (fin de démonstration), utilisé par les modèles sans amsthm
\providecommand{\qed}{\hfill$\square$}
% \barre : double barre des valeurs interdites dans un tableau de variations (tkz-tab)
\providecommand{\barre}{\ensuremath{\|}}
% environnement proof (amsthm non chargé) utilisé par les modèles
\ifdefined\proof\else\newenvironment{proof}[1][Démonstration]{\par\noindent\textit{#1.}\ }{\qed\par}\fi
\usepackage{lastpage}
\usepackage{hyperref}
\hypersetup{hidelinks}

% ============================================================
% Palette « Pop » OnBuch+ — mêmes hex que la classe P (plus_ui.dart)
% ============================================================
\definecolor{poporange}{HTML}{F59321}
\definecolor{popdark}{HTML}{E07A0C}
\definecolor{popcream}{HTML}{FFF6E8}
\definecolor{popink}{HTML}{1C1714}
\definecolor{poppurple}{HTML}{7A5AE0}
\definecolor{popgreen}{HTML}{2E9E6A}
\definecolor{poppink}{HTML}{E0457B}
\definecolor{popblue}{HTML}{2F7DE1}
\definecolor{popgold}{HTML}{C98A12}
\definecolor{popmuted}{HTML}{8A7F76}
\definecolor{popline}{HTML}{EADFD2}
\definecolor{popgray}{HTML}{8A7F76}
\colorlet{popgrey}{popgray}
% Alias tolérants (le modèle nomme parfois les couleurs différemment)
\colorlet{popwhite}{white}
\colorlet{popblack}{popink}
\colorlet{popred}{poppink}
\colorlet{popyellow}{popgold}
% Teintes claires (fonds de tableaux/encadrés) — le modèle nomme parfois
% les couleurs avec un suffixe L (ex. popblueL) sans les avoir définies.
\colorlet{poporangeL}{poporange!20}
\colorlet{popdarkL}{popdark!20}
\colorlet{poppurpleL}{poppurple!20}
\colorlet{popgreenL}{popgreen!20}
\colorlet{poppinkL}{poppink!20}
\colorlet{popblueL}{popblue!20}
\colorlet{popgoldL}{popgold!20}
\colorlet{popredL}{popred!20}
\colorlet{popgrayL}{popgray!20}
% Teintes claires pour fonds de tableaux / zones
\colorlet{popblueL}{popblue!10!white}
\colorlet{poporangeL}{poporange!14!white}
\colorlet{popgreenL}{popgreen!12!white}
\colorlet{poppurpleL}{poppurple!12!white}
\colorlet{poppinkL}{poppink!12!white}
\colorlet{popgoldL}{popgold!14!white}

\pagecolor{popcream}

% ============================================================
% Typographie
% ============================================================
\defaultfontfeatures{Ligatures=TeX}
\setmainfont{PlusJakartaSans-Medium.ttf}[Path=../../../fonts/,BoldFont=PlusJakartaSans-Bold.ttf]
% Symboles, API et emojis absents de la police principale : repli sur DejaVu Sans (caractères actifs)
\newfontface\symfont{DejaVuSans.ttf}[Path=../../../fonts/]
\makeatletter
\newcommand\symactive[1]{\catcode#1=13 \begingroup\lccode`\~=#1 \lowercase{\endgroup\def~}{{\symfont\char#1}}}
\newcommand\symblank[1]{\catcode#1=13 \begingroup\lccode`\~=#1 \lowercase{\endgroup\def~}{}}
\newcommand\symhyph[1]{\catcode#1=13 \begingroup\lccode`\~=#1 \lowercase{\endgroup\def~}{\mbox{-}}}
\newcount\symc
\newcommand\symrange[2]{\symc=#1 \@whilenum\symc<#2 \do{\expandafter\symactive\expandafter{\the\symc}\advance\symc by 1 }}
\newcommand\symblankrange[2]{\symc=#1 \@whilenum\symc<#2 \do{\expandafter\symblank\expandafter{\the\symc}\advance\symc by 1 }}
\makeatother
\symrange{"0250}{"02FF}\symactive{"2713}\symactive{"2714}\symactive{"2717}\symactive{"2718}\symactive{"2610}\symactive{"2611}\symactive{"2612}
\symrange{"2600}{"27BF}
\catcode"2705=13 \begingroup\lccode`\~="2705 \lowercase{\endgroup\def~}{{\symfont\char"2713}}\symblank{"26BD}\symblank{"26A1}
\symrange{"2460}{"257F}\symactive{"223C}\symactive{"2248}\symactive{"2260}
\symactive{"01F8}\symactive{"01F9}\symactive{"1E3E}\symactive{"1E3F}
\symhyph{"2011}\symhyph{"2E1A}\symblankrange{"1F300}{"1FAFF}\symblank{"FE0F}
% Caractères chinois : WenQuanYi Zen Hei (sous-ensemble GB2312 + pinyin), gras/italique simulés
\setCJKmainfont{WenQuanYiZenHei-subset.ttf}[Path=../../../fonts/,AutoFakeBold=3,AutoFakeSlant=0.2]
\xeCJKsetup{CJKspace=false}
% Pinyin : voyelles à caron (ǎ ǐ ǒ ǔ ǖ ǘ ǚ ǜ) absentes de la police principale -> caractères actifs
\newfontface\pyfont{WenQuanYiZenHei-subset.ttf}[Path=../../../fonts/]
\newcommand\pyactive[1]{\catcode#1=13 \begingroup\lccode`\~=#1 \lowercase{\endgroup\def~}{{\pyfont\char#1}}}
\pyactive{"01CD}\pyactive{"01CE}\pyactive{"01CF}\pyactive{"01D0}\pyactive{"01D1}\pyactive{"01D2}\pyactive{"01D3}\pyactive{"01D4}\pyactive{"01D5}\pyactive{"01D6}\pyactive{"01D7}\pyactive{"01D8}\pyactive{"01D9}\pyactive{"01DA}\pyactive{"01DB}\pyactive{"01DC}
% Maths en sans-serif (Fira Math) avec chiffres et lettres droites de Plus
% Jakarta, pour que les formules se fondent dans le texte.
\usepackage[math-style=ISO,bold-style=ISO]{unicode-math}
\setmathfont{FiraMath-Regular.otf}[Path=../../../fonts/]
\setmathfont{PlusJakartaSans-Medium.ttf}[Path=../../../fonts/,range={up/num,up/{Latin,latin},\mathup}]
\usepackage{icomma} % virgule décimale sans espace en mode math
% Caractères mathématiques Unicode absents des polices de texte (Plus Jakarta,
% Archivo Black) : rendus par leur équivalent mathématique, en texte comme en
% formule — sinon ils s'affichent en carré vide (ex. « Arithmétique dans ℤ »).
\usepackage{newunicodechar}
\newunicodechar{ℝ}{\ensuremath{\mathbb{R}}}
\newunicodechar{ℤ}{\ensuremath{\mathbb{Z}}}
\newunicodechar{ℂ}{\ensuremath{\mathbb{C}}}
\newunicodechar{ℕ}{\ensuremath{\mathbb{N}}}
\newunicodechar{ℚ}{\ensuremath{\mathbb{Q}}}
\newunicodechar{𝔻}{\ensuremath{\mathbb{D}}}
\newunicodechar{α}{\ensuremath{\alpha}}
\newunicodechar{β}{\ensuremath{\beta}}
\newunicodechar{γ}{\ensuremath{\gamma}}
\newunicodechar{δ}{\ensuremath{\delta}}
\newunicodechar{ε}{\ensuremath{\varepsilon}}
\newunicodechar{θ}{\ensuremath{\theta}}
\newunicodechar{λ}{\ensuremath{\lambda}}
\newunicodechar{μ}{\ensuremath{\mu}}
\newunicodechar{σ}{\ensuremath{\sigma}}
\newunicodechar{τ}{\ensuremath{\tau}}
\newunicodechar{φ}{\ensuremath{\varphi}}
\newunicodechar{ψ}{\ensuremath{\psi}}
\newunicodechar{ω}{\ensuremath{\omega}}
\newunicodechar{Σ}{\ensuremath{\Sigma}}
% majuscules produites par \MakeUppercase dans les titres (α-aminés -> Α)
\newunicodechar{Α}{\ensuremath{\alpha}}
\newunicodechar{Β}{\ensuremath{\beta}}
\newunicodechar{Γ}{\ensuremath{\Gamma}}
\newunicodechar{Δ}{\ensuremath{\Delta}}
\newunicodechar{Ε}{\ensuremath{\varepsilon}}
\newunicodechar{Θ}{\ensuremath{\Theta}}
\newunicodechar{Λ}{\ensuremath{\Lambda}}
\newunicodechar{Μ}{\ensuremath{\mu}}
\newunicodechar{Τ}{\ensuremath{\tau}}
\newunicodechar{Φ}{\ensuremath{\Phi}}
\newunicodechar{Ψ}{\ensuremath{\Psi}}
\newunicodechar{Ω}{\ensuremath{\Omega}}
\newunicodechar{↦}{\ensuremath{\mapsto}}
\newunicodechar{∈}{\ensuremath{\in}}
\newunicodechar{∉}{\ensuremath{\notin}}
\newunicodechar{∧}{\ensuremath{\wedge}}
\newunicodechar{⇔}{\ensuremath{\Leftrightarrow}}
\newunicodechar{⟺}{\ensuremath{\Longleftrightarrow}}
\newunicodechar{⇒}{\ensuremath{\Rightarrow}}
\newunicodechar{⟹}{\ensuremath{\Longrightarrow}}
\newunicodechar{∘}{\ensuremath{\circ}}
\newunicodechar{∪}{\ensuremath{\cup}}
\newunicodechar{∩}{\ensuremath{\cap}}
\newunicodechar{≡}{\ensuremath{\equiv}}
\newunicodechar{⊂}{\ensuremath{\subset}}
\newunicodechar{⊥}{\ensuremath{\perp}}
\newunicodechar{∃}{\ensuremath{\exists}}
\newunicodechar{∀}{\ensuremath{\forall}}
\newunicodechar{∛}{\ensuremath{\sqrt[3]{\,}}}
\newunicodechar{∜}{\ensuremath{\sqrt[4]{\,}}}
\newunicodechar{⃗}{\ensuremath{{}^{\to}}}
\newunicodechar{ⁿ}{\ifmmode^{n}\else\textsuperscript{\ensuremath{n}}\fi}
\newunicodechar{ˣ}{\ifmmode^{x}\else\textsuperscript{\ensuremath{x}}\fi}
\newunicodechar{ᵇ}{\ifmmode^{b}\else\textsuperscript{\ensuremath{b}}\fi}
\newunicodechar{ʳ}{\ifmmode^{r}\else\textsuperscript{\ensuremath{r}}\fi}
\newunicodechar{ᵘ}{\ifmmode^{u}\else\textsuperscript{\ensuremath{u}}\fi}
\newunicodechar{ʸ}{\ifmmode^{y}\else\textsuperscript{\ensuremath{y}}\fi}
\newunicodechar{ᵃ}{\ifmmode^{a}\else\textsuperscript{\ensuremath{a}}\fi}
\newunicodechar{ᵉ}{\ifmmode^{e}\else\textsuperscript{\ensuremath{e}}\fi}
\newunicodechar{ᵏ}{\ifmmode^{k}\else\textsuperscript{\ensuremath{k}}\fi}
\newunicodechar{ᵖ}{\ifmmode^{p}\else\textsuperscript{\ensuremath{p}}\fi}
\newunicodechar{ᴺ}{\ifmmode^{N}\else\textsuperscript{\ensuremath{N}}\fi}
\newunicodechar{ᶜ}{\ifmmode^{c}\else\textsuperscript{\ensuremath{c}}\fi}
\newunicodechar{ʰ}{\ifmmode^{h}\else\textsuperscript{\ensuremath{h}}\fi}
\newunicodechar{ᵠ}{\ifmmode^{q}\else\textsuperscript{\ensuremath{q}}\fi}
\newunicodechar{ₙ}{\ifmmode_{n}\else\textsubscript{\ensuremath{n}}\fi}
\newunicodechar{ₐ}{\ifmmode_{a}\else\textsubscript{\ensuremath{a}}\fi}
\newunicodechar{ᵢ}{\ifmmode_{i}\else\textsubscript{\ensuremath{i}}\fi}
\newunicodechar{ᵣ}{\ifmmode_{r}\else\textsubscript{\ensuremath{r}}\fi}
\newunicodechar{ₖ}{\ifmmode_{k}\else\textsubscript{\ensuremath{k}}\fi}
\newunicodechar{ᵦ}{\ifmmode_{\beta}\else\textsubscript{\ensuremath{\beta}}\fi}
\newunicodechar{ₓ}{\ifmmode_{x}\else\textsubscript{\ensuremath{x}}\fi}
\newfontfamily\popdisplay{ArchivoBlack-Regular.ttf}[Path=../../../fonts/]
\newfontfamily\poplogo{SpaceGrotesk-Bold.ttf}[Path=../../../fonts/]
\newfontfamily\popsemi{PlusJakartaSans-SemiBold.ttf}[Path=../../../fonts/]
\newfontfamily\popxbold{PlusJakartaSans-ExtraBold.ttf}[Path=../../../fonts/]
\newfontfamily\popxboldls{PlusJakartaSans-ExtraBold.ttf}[Path=../../../fonts/,LetterSpace=3.0]

\setlist[enumerate]{leftmargin=1.7em,itemsep=5pt,topsep=4pt}
\setlist[itemize]{leftmargin=1.7em,itemsep=4pt,topsep=4pt,label={\color{poporange}\textbullet}}
\setlength{\parindent}{0pt}
\setlength{\parskip}{5pt}
\renewcommand{\arraystretch}{1.25}

% pgfplots : style Pop par défaut
\pgfplotsset{
  every axis/.append style={axis line style={popink,line width=1pt},tick style={popink},
    grid style={popline,line width=0.5pt},label style={font=\small},tick label style={font=\footnotesize}},
  popcurve/.style={poporange,line width=1.8pt,no markers,smooth},
}
% Même style disponible dans un \draw TikZ simple (hors pgfplots)
\tikzset{popcurve/.style={poporange,line width=1.8pt}}
\tikzset{popgrid/.style={popline,very thin}}
\colorlet{popcurve}{poporange} % le nom du style est parfois utilisé comme couleur

% ============================================================
% Pill / squiggle
% ============================================================
\newcommand{\poppill}[3]{%
  \tikz[baseline=-0.55ex]\node[draw=popink,line width=1.2pt,rounded corners=2.6pt,%
    fill=#2,inner xsep=6.5pt,inner ysep=3pt]{%
    \popxboldls\fontsize{7}{7}\selectfont\color{#3}\MakeUppercase{#1}};%
}
\newcommand{\popsquiggle}[1]{%
  \tikz[baseline=-0.4ex]{
    \draw[#1,line width=2.6pt,line cap=round]
      plot [smooth] coordinates {(0,0.09) (0.34,0.01) (0.68,0.21) (1.02,0.01) (1.36,0.10)};
  }%
}
% Mot-clé surligné (marqueur orange sous le texte)
\newcommand{\cle}[1]{{\popxbold\color{popdark}#1}}

% ============================================================
% En-tête / pied
% ============================================================
\pagestyle{fancy}
\fancyhf{}
\renewcommand{\headrulewidth}{0pt}
\renewcommand{\footrulewidth}{0pt}
\lhead{\raisebox{-0.15ex}{\poplogo\fontsize{13}{13}\selectfont\color{popink}OnBuch\color{poporange}+}}
\rhead{\poppill{\DOCMATIERE\ \textbullet\ \DOCNIVEAU\ \textbullet\ Leçon \DOCLECON}{white}{popink}}
\lfoot{\popxbold\fontsize{7.6}{7.6}\selectfont\color{popmuted}\MakeUppercase{Cours signé OnBuch+}\ \textbullet\ {\popsemi\DOCTITRE}}
\rfoot{\tikz[baseline=-0.6ex]\node[rounded corners=8.5pt,fill=popink,minimum height=17pt,minimum width=17pt,inner xsep=5pt,inner sep=0]{\popxbold\fontsize{8}{8}\selectfont\color{popcream}\thepage\,/\,\pageref*{LastPage}};}

% ============================================================
% Titres
% ============================================================
\newcommand{\chaptitre}[2]{%
  \begin{tcolorbox}[enhanced,colback=poporange,colframe=popink,boxrule=2.6pt,arc=11pt,
      drop shadow=popink,left=15pt,right=15pt,top=12pt,bottom=11pt,before skip=3pt,after skip=18pt]
    \poppill{#2}{popink}{popcream}\par\vspace{9pt}
    {\raggedright\hyphenpenalty=10000\exhyphenpenalty=10000\popdisplay\fontsize{22}{24}\selectfont\color{popcream}\MakeUppercase{#1}\par}
  \end{tcolorbox}
}
% Section numérotée : pastille orange avec le numéro + titre Archivo
\newcounter{popsec}
\newcommand{\coursec}[1]{%
  \stepcounter{popsec}\setcounter{popsub}{0}%
  \par\vspace{16pt}\noindent
  \begin{minipage}{\linewidth}
  \tikz[baseline=-0.7ex]\node[draw=popink,line width=1.6pt,fill=poporange,rounded corners=4pt,minimum size=22pt,inner sep=0]{\popdisplay\fontsize{12}{12}\selectfont\color{popcream}\thepopsec};%
  \hspace{8pt}{\popdisplay\fontsize{15}{17}\selectfont\color{popink}\MakeUppercase{#1}}\par
  \vspace{4pt}\noindent{\color{poporange}\rule{36pt}{3.6pt}}
  \end{minipage}\par\nopagebreak\vspace{8pt}
}
\newcounter{popsub}
\newcommand{\courssub}[1]{%
  \stepcounter{popsub}%
  \par\vspace{9pt}\noindent{\popxbold\fontsize{11.5}{13}\selectfont\color{popink}{\color{poporange}\thepopsec.\thepopsub}\hspace{6pt}#1}\par\nopagebreak\vspace{4pt}
}

% ============================================================
% Boîtes pédagogiques — même recette : fond blanc, bordure épaisse colorée,
% ombre dure encre, badge plein. Titre optionnel [#1].
% ============================================================
\tcbset{popbase/.style={breakable,enhanced,colback=white,boxrule=2.3pt,arc=9pt,
  shadow={3pt}{-3pt}{0pt}{popink},left=14pt,right=14pt,top=11pt,bottom=11pt,before skip=11pt,after skip=15pt,
  fontupper=\popsemi\fontsize{10.6}{14.5}\selectfont\color{popink},
  fontlower=\popsemi\fontsize{10.6}{14.5}\selectfont\color{popink}}}
% Coche dessinée (le glyphe ✓ n'existe pas dans les polices du document)
\newcommand{\popcheck}[1]{\tikz[baseline=-0.5ex]\draw[#1,line width=1.6pt,line cap=round,line join=round](-0.13,0.0)--(-0.04,-0.09)--(0.13,0.1);}
\providecommand{\coche}{\popcheck{popgreen}}
\providecommand{\resultat}[1]{\par\smallskip\noindent\fcolorbox{popgreen}{popgreenL}{\parbox{\dimexpr\linewidth-2\fboxsep-2\fboxrule\relax}{\textbf{Résultat.} #1}}\par\smallskip}
\newcommand{\popboxhead}[3]{\poppill{#1}{#2}{white}\ifx\relax#3\relax\else\hspace{6pt}{\popxbold\fontsize{10.6}{12}\selectfont\color{popink}#3}\fi\par\vspace{7pt}}

\newtcolorbox{aretenir}[1][]{popbase,colframe=popgreen,before upper={\popboxhead{À retenir}{popgreen}{#1}}}
% Texte en chinois (document de lecture, dialogue, extrait) : cadre
% d'étude. Un titre optionnel (source, auteur) s'affiche dans l'en-tête.
\newtcolorbox{textebox}[1][]{popbase,colframe=popblue,before upper={\popboxhead{文本 · Texte}{popblue}{#1}},after upper={}}
% Marques ✓ / ✗ (forme correcte / fautive) souvent utilisées par les modèles
\providecommand{\cross}{\textcolor{poppink}{\ensuremath{\times}}}
\providecommand{\xmark}{\cross}\providecommand{\crossmark}{\cross}\providecommand{\cmark}{\ensuremath{\checkmark}}
% Ligne de réponse / justification à compléter (inventée par les modèles)
\providecommand{\justification}[1]{\par\noindent\textit{Justification~:} \dotfill\par}
% \textipa (package tipa non chargé) : la police ne couvre pas l'API -> simple passage
\providecommand{\textipa}[1]{#1}
% Soulignement (ulem non chargé) : \uline, \ul
\providecommand{\uline}[1]{\underline{#1}}\providecommand{\ul}[1]{\underline{#1}}
% \glossaire{\item ...} inventée par les modèles : liste de glossaire
\providecommand{\glossaire}[1]{\begin{itemize}#1\end{itemize}}
% \alert (beamer) inventée par les modèles : texte mis en évidence
\providecommand{\alert}[1]{\textbf{\textcolor{poppink}{#1}}}
% variantes ASCII d'ordinaux écrites en macro
\providecommand{\iere}{\textsuperscript{re}}\providecommand{\ieme}{\textsuperscript{e}}\providecommand{\ere}{\textsuperscript{re}}\providecommand{\eme}{\textsuperscript{e}}\providecommand{\er}{\textsuperscript{er}}\providecommand{\ie}{\textsuperscript{e}}
% Ordinaux français écrits en macro par les modèles : 1\ière, 3\ième
\newcommand{\ière}{\textsuperscript{re}}\newcommand{\ième}{\textsuperscript{e}}
% \exemple (inventée par les modèles) : étiquette « Exemple : »
\providecommand{\exemple}{\textbf{Exemple~:}\ }
% \square utilisable aussi hors mode math (cases à cocher)
\AtBeginDocument{\let\squareMath\square\renewcommand{\square}{\ensuremath{\squareMath}}}
% Mot ou phrase en chinois à l'intérieur d'un paragraphe français
\newcommand{\zh}[1]{\ifmmode\text{#1}\else{#1}\fi}
\let\eng\zh \let\esp\zh \let\all\zh % alias tolérés
% flèches d'intonation inventées par les modèles
\providecommand{\textsearrow}{}\renewcommand{\textsearrow}{\ensuremath{\searrow}}\providecommand{\textnearrow}{}\renewcommand{\textnearrow}{\ensuremath{\nearrow}}
% \fr{..} : mot français (inventé par les modèles)
\providecommand{\fr}[1]{\textit{#1}}
% zhbox : encadré de texte chinois inventé par les modèles
\newenvironment{zhbox}{\begin{quote}}{\end{quote}}
% \courssubsub : titre de niveau 3 inventé par les modèles
\providecommand{\courssubsub}[1]{\par\medskip\noindent\textbf{#1}\par\smallskip}
% \zhbf : texte chinois en gras inventé par les modèles
\providecommand{\zhbf}[1]{\textbf{#1}}
% \vs : « versus » inventé par les modèles
\providecommand{\vs}{\textit{vs}\ }
% \blank : case à compléter inventée par les modèles
\providecommand{\blank}[1][]{\underline{\hspace{1.6em}}}
\newtcolorbox{exemplebox}[1][]{popbase,colframe=poporange,before upper={\popboxhead{Exemple}{poporange}{#1}}}
\newtcolorbox{definition}[1][]{popbase,colframe=popblue,before upper={\popboxhead{Définition}{popblue}{#1}}}
\newtcolorbox{propriete}[1][]{popbase,colframe=poppurple,before upper={\popboxhead{Propriété}{poppurple}{#1}}}
\newtcolorbox{methode}[1][]{popbase,colframe=popgold,before upper={\popboxhead{Méthode}{popgold}{#1}}}
\newtcolorbox{attention}[1][]{popbase,colframe=poppink,before upper={\popboxhead{Attention}{poppink}{#1}}}
\newtcolorbox{experience}[1][]{popbase,colframe=popblue,colback=popblueL,before upper={\popboxhead{Expérience}{popblue}{#1}}}
% « Le savais-tu ? » : carte violette pleine (contexte, culture, Cameroun)
\newtcolorbox{savaistu}[1][]{popbase,colback=poppurple,colframe=popink,
  fontupper=\popsemi\fontsize{10.4}{14.2}\selectfont\color{white},
  before upper={\poppill{Le savais-tu\,?}{white}{poppurple}\ifx\relax#1\relax\else\hspace{6pt}{\popxbold\color{white}#1}\fi\par\vspace{7pt}}}
% Exercice résolu : énoncé en haut, \tcblower puis la solution sur fond crème
\newcounter{popexo}\newcounter{popexr}
\newtcolorbox{exoresolu}[1][]{popbase,colframe=poporange,colbacklower=popcream,
  segmentation style={popink,line width=1.2pt,dashed},
  before upper={\stepcounter{popexr}\popboxhead{Exercice résolu \thepopexr}{poporange}{#1}},
  before lower={\poppill{Solution}{popink}{popcream}\par\vspace{6pt}}}
% Exercice (non résolu en place) — la correction est dans la partie Corrigés
\newtcolorbox{exercice}[1][]{popbase,colframe=popink,
  before upper={\stepcounter{popexo}\popboxhead{Exercice \thepopexo}{popink}{#1}}}
% Corrigé
\newtcolorbox{corrige}[1][]{popbase,colframe=popgreen,colback=popgreenL,
  before upper={\popboxhead{Corrigé}{popgreen}{#1}}}
% Objectifs / prérequis / bilan
\newtcolorbox{objectifs}{popbase,colframe=popink,colback=poporangeL,
  before upper={\popboxhead{Objectifs de la leçon}{popink}{}}}
\newtcolorbox{prerequis}{popbase,colframe=popink,colback=popblueL,
  before upper={\popboxhead{Prérequis}{popblue}{}}}
\newtcolorbox{bilan}{popbase,colframe=popink,colback=white,boxrule=2.8pt,
  before upper={\popboxhead{Fiche bilan}{poporange}{L'essentiel à connaître pour l'examen}}}
% Figure encadrée avec légende
\newtcolorbox{popfigure}[1][]{enhanced,breakable=false,colback=white,colframe=popline,boxrule=1.4pt,arc=8pt,
  left=10pt,right=10pt,top=10pt,bottom=8pt,before skip=10pt,after skip=12pt,halign=center,
  lower separated=false,halign lower=center,
  fontlower=\popsemi\fontsize{9}{11.5}\selectfont\color{popmuted},
  #1}
\newcommand{\legende}[1]{\par\vspace{6pt}{\popsemi\fontsize{9}{11.5}\selectfont\color{popmuted}#1}}

% ============================================================
% Signature OnBuch+
% ============================================================
\newcommand{\poplogoinline}[1]{{\poplogo\fontsize{#1}{#1}\selectfont\color{popink}OnBuch\color{poporange}+}}
\newcommand{\signatureonbuch}{%
  \begin{tcolorbox}[enhanced,colback=popink,colframe=popink,boxrule=2.2pt,arc=10pt,
      drop shadow=poporange,left=14pt,right=14pt,top=10pt,bottom=10pt,before skip=0pt,after skip=0pt]
    \tikz[baseline=-0.7ex]\node[circle,fill=popgreen,draw=popcream,line width=1.3pt,minimum size=22pt,inner sep=0]{\popcheck{white}};%
    \hspace{9pt}\begin{minipage}[c]{0.62\linewidth}
      {\popxbold\fontsize{12}{13}\selectfont\color{popcream}Signé\ \textbullet\ {\poplogo OnBuch\color{poporange}+}}\par\vspace{2pt}
      {\raggedright\popsemi\fontsize{8}{10.5}\selectfont\color{popline}\MakeUppercase{Équipe pédagogique OnBuch+}\\\MakeUppercase{Conforme au programme officiel MINESEC}\par}
    \end{minipage}\hfill
    {\poplogo\fontsize{22}{22}\selectfont\color{popcream}OnBuch\color{poporange}+}
  \end{tcolorbox}%
}

% ============================================================
% Page de garde
% #1 titre · #2 matière · #3 niveau · #4 module · #5 n° leçon · #6 durée ·
% #7 description / objectifs (texte ou itemize)
% ============================================================
\newcommand{\pagegarde}[7]{%
  \thispagestyle{empty}
  \newgeometry{a4paper,margin=1.7cm,top=1.6cm,bottom=1.4cm}%
  \begin{tikzpicture}[remember picture,overlay]
    \fill[poporange!18] ([xshift=-3.2cm,yshift=-3.6cm]current page.north east) circle (4.6cm);
    \fill[poppurple!14] ([xshift=2.4cm,yshift=7.6cm]current page.south west) circle (3.8cm);
  \end{tikzpicture}%
  {\poplogo\fontsize{30}{30}\selectfont\color{popink}OnBuch\color{poporange}+}\hfill
  \raisebox{0.6ex}{\poppill{Cours signé \textbullet\ Premium}{popink}{popcream}}\par
  \vspace{1pt}\popsquiggle{poppurple}\par
  \vspace{8pt}
  \begin{tcolorbox}[enhanced,colback=poporange,colframe=popink,boxrule=2.8pt,arc=13pt,
      drop shadow=popink,left=18pt,right=18pt,top=12pt,bottom=13pt,after skip=10pt]
    \poppill{#2 \textbullet\ #3}{popink}{popcream}\hspace{5pt}\poppill{#4}{white}{popink}\par\vspace{8pt}
    {\popdisplay\fontsize{14}{14}\selectfont\color{popink}LEÇON #5}\par\vspace{4pt}
    {\raggedright\hyphenpenalty=10000\exhyphenpenalty=10000\popdisplay\fontsize{25}{27}\selectfont\color{popcream}\MakeUppercase{#1}\par}
    \vspace{8pt}
    {\popxbold\fontsize{9.5}{9.5}\selectfont\color{popink}\MakeUppercase{Durée conseillée : #6}}
  \end{tcolorbox}
  \begin{tcolorbox}[enhanced,colback=white,colframe=popink,boxrule=2.2pt,arc=10pt,
      drop shadow=popink,left=16pt,right=16pt,top=10pt,bottom=10pt,after skip=8pt]
    \poppill{Dans cette leçon}{poppurple}{white}\par\vspace{5pt}
    \popsemi\fontsize{9.3}{12.6}\selectfont\color{popink}#7
  \end{tcolorbox}
  \noindent\begin{minipage}[t]{0.487\linewidth}
    \begin{tcolorbox}[enhanced,colback=popgreen,colframe=popink,boxrule=2.2pt,arc=10pt,
        drop shadow=popink,left=11pt,right=11pt,top=10pt,bottom=10pt,height=3.1cm,valign=center]
      \tcbox[colback=white,colframe=popink,boxrule=1.3pt,arc=5pt,boxsep=0pt,nobeforeafter,
        left=3pt,right=3pt,top=3pt,bottom=3pt,valign=center]{\qrcode[height=1.9cm]{https://play.google.com/store/apps/details?id=com.luvvix.onbuch&hl=fr}}%
      \hspace{8pt}\begin{minipage}[c]{0.5\linewidth}
        {\popdisplay\fontsize{11}{12.5}\selectfont\color{white}\MakeUppercase{Télécharge OnBuch}}\par\vspace{4pt}
        {\raggedright\popsemi\fontsize{8.4}{11}\selectfont\color{white}Gratuit \textbullet\ Google Play\\Tuteur IA, annales, concours, résultats\par}
      \end{minipage}
    \end{tcolorbox}
  \end{minipage}\hfill
  \begin{minipage}[t]{0.487\linewidth}
    \begin{tcolorbox}[enhanced,colback=poppurple,colframe=popink,boxrule=2.2pt,arc=10pt,
        drop shadow=popink,left=11pt,right=11pt,top=10pt,bottom=10pt,height=3.1cm,valign=center]
      \tikz[baseline=-0.7ex]\node[circle,fill=white,draw=popink,line width=1.3pt,minimum size=1.6cm,inner sep=0]{\popdisplay\fontsize{24}{24}\selectfont\color{poppurple}+};%
      \hspace{8pt}\begin{minipage}[c]{0.6\linewidth}
        {\popdisplay\fontsize{11}{12.5}\selectfont\color{white}\MakeUppercase{Passe à OnBuch+}}\par\vspace{4pt}
        {\raggedright\popsemi\fontsize{8.4}{11}\selectfont\color{white}Tous les cours signés, Tuteur IA illimité, fiches PDF — onglet OnBuch+ de l'app\par}
      \end{minipage}
    \end{tcolorbox}
  \end{minipage}
  \vspace{8pt}
  \popcontenu
  \vfill
  \signatureonbuch
  \restoregeometry
  \newpage
}

% Rangée « Ce cours contient » (page de garde)
\newcommand{\poptile}[3]{% #1 couleur #2 gros texte #3 légende
  \begin{tcolorbox}[enhanced,colback=#1,colframe=popink,boxrule=2pt,arc=9pt,shadow={2.5pt}{-2.5pt}{0pt}{popink},
      left=6pt,right=6pt,top=9pt,bottom=9pt,halign=center,valign=center,height=1.75cm,width=\linewidth,nobeforeafter]
    {\popdisplay\fontsize{12.5}{13}\selectfont\color{white}\MakeUppercase{#2}}\par\vspace{3pt}
    {\centering\popsemi\fontsize{7.8}{9.5}\selectfont\color{white}#3\par}
  \end{tcolorbox}}
\newcommand{\popcontenu}{%
  \par\noindent{\popxboldls\fontsize{7.5}{7.5}\selectfont\color{popmuted}CE COURS SIGNÉ CONTIENT}\par\vspace{7pt}
  \noindent\begin{minipage}[t]{0.235\linewidth}\poptile{popblue}{Cours}{complet, illustré, pas à pas}\end{minipage}\hfill
  \begin{minipage}[t]{0.235\linewidth}\poptile{poporange}{Méthodes}{et exercices résolus}\end{minipage}\hfill
  \begin{minipage}[t]{0.235\linewidth}\poptile{poppink}{Exercices}{type examen + corrigés}\end{minipage}\hfill
  \begin{minipage}[t]{0.235\linewidth}\poptile{popgreen}{Fiche bilan}{l'essentiel à retenir}\end{minipage}\par}

% Sommaire
\newtcolorbox{sommaire}{enhanced,colback=white,colframe=popink,boxrule=2.2pt,arc=10pt,
  drop shadow=popink,left=18pt,right=18pt,top=13pt,bottom=13pt,before skip=6pt,after skip=20pt,
  before upper={\poppill{Sommaire}{poppurple}{white}\par\vspace{9pt}\popsemi\fontsize{10.6}{14.5}\selectfont\color{popink}}}

% Signature de fin de cours — poussée tout en bas de la dernière page
\newcommand{\findecours}{%
  \par\vfill\nopagebreak
  \begin{center}\popsquiggle{poporange}\end{center}\vspace{4pt}
  \signatureonbuch
}
\AtBeginDocument{\def\llbracket{\mathopen{[\mkern-3.5mu[}}\def\rrbracket{\mathclose{]\mkern-3.5mu]}}} % intervalles d'entiers (glyphe ⟦ absent de la police)
% Glyphes absents de Fira Math : redéfinis à partir de glyphes disponibles
\AtBeginDocument{%
  \def\setminus{\mathbin{\backslash}}\let\smallsetminus\setminus
  \def\vdots{\mathord{\vbox{\baselineskip3.2pt\lineskiplimit0pt\kern4pt\hbox{.}\hbox{.}\hbox{.}}}}%
  \def\ddots{\mathinner{\mkern1mu\raise6.5pt\hbox{.}\mkern2mu\raise3.6pt\hbox{.}\mkern2mu\raise0.7pt\hbox{.}\mkern1mu}}%
  \def\triangle{\mathord{\scalebox{0.9}{$\symup{\Delta}$}}}%
  \def\nabla{\mathord{\raisebox{\depth}{\scalebox{1}[-1]{$\symup{\Delta}$}}}}%
  \def\checkmark{\popcheck{popdark}}%
  \def\star{\mathbin{\ast}}\def\bigstar{\mathord{\ast}}%
  \def\diamond{\mathbin{\scalebox{0.6}{$\Diamond$}}}%
  \def\bigcup{\mathop{\mathchoice{\raisebox{-0.3em}{\scalebox{1.7}{$\cup$}}}{\scalebox{1.25}{$\cup$}}{\cup}{\cup}}\displaylimits}%
  \def\bigcap{\mathop{\mathchoice{\raisebox{-0.3em}{\scalebox{1.7}{$\cap$}}}{\scalebox{1.25}{$\cap$}}{\cap}{\cap}}\displaylimits}%
  \def\lesseqqgtr{\mathrel{\lessgtr}}\def\bigoplus{\mathop{\oplus}\displaylimits}%
}
\providecommand{\ssi}{si et seulement si} % abréviation fréquente
\AtBeginDocument{\providecommand{\wideparen}{\overparen}} % arc de cercle

\begin{document}

\pagegarde{Leçon 28 — Expression orale : gestion du budget}{Chinois}{Tle A}{Module 4 : Activités économiques et monde du travail}{ZH28}{2 h}{Cette leçon d'expression orale entraîne les élèves à présenter un exposé simple et à participer à un débat guidé sur la gestion du budget (argent de poche, dépenses, économies). Elle fournit les formules fonctionnelles, les dialogues modèles et le travail des tons nécessaires pour s'exprimer avec aisance à l'oral.
\vspace{4pt}
\begin{multicols}{2}\small
\begin{itemize}
\item À la fin de cette leçon, je sais utiliser le vocabulaire essentiel du budget (argent, dépenser, économiser, prix) à l'oral.
\item À la fin de cette leçon, je sais présenter oralement un exposé simple de 8 à 10 phrases sur mon budget en chinois.
\item À la fin de cette leçon, je sais donner mon avis avec les structures 我觉得..., 我认为... et 对我来说...
\item À la fin de cette leçon, je sais exprimer l'accord et le désaccord poliment dans un débat (我同意 / 我不太同意).
\item À la fin de cette leçon, je sais relancer la parole et poser des questions à mes camarades (你呢？你觉得怎么样？).
\item À la fin de cette leçon, je sais prononcer correctement les tons des mots clés de la leçon, en particulier les suites de 3e tons.
\end{itemize}\end{multicols}}

\begin{sommaire}
\begin{multicols}{2}
\begin{enumerate}
\item Vocabulaire du budget
\item Dialogues modèles
\item Formules fonctionnelles
\item Les tons à travailler
\item Jeux de rôle et débat guidé
\item Production orale guidée et évaluation
\item Activité d'intégration
\item Méthodes et exercices résolus
\item Exercices
\item Corrigés des exercices
\item Fiche bilan
\end{enumerate}
\end{multicols}
\end{sommaire}

\begin{objectifs}
\begin{itemize}
\item À la fin de cette leçon, je sais utiliser le vocabulaire essentiel du budget (argent, dépenser, économiser, prix) à l'oral.
\item À la fin de cette leçon, je sais présenter oralement un exposé simple de 8 à 10 phrases sur mon budget en chinois.
\item À la fin de cette leçon, je sais donner mon avis avec les structures 我觉得..., 我认为... et 对我来说...
\item À la fin de cette leçon, je sais exprimer l'accord et le désaccord poliment dans un débat (我同意 / 我不太同意).
\item À la fin de cette leçon, je sais relancer la parole et poser des questions à mes camarades (你呢？你觉得怎么样？).
\item À la fin de cette leçon, je sais prononcer correctement les tons des mots clés de la leçon, en particulier les suites de 3e tons.
\item À la fin de cette leçon, je sais jouer un dialogue de négociation budgétaire avec un partenaire.
\item À la fin de cette leçon, je sais comparer brièvement la gestion de l'argent de poche au Cameroun et en Chine.
\end{itemize}
\end{objectifs}

\begin{prerequis}
\begin{itemize}
\item Les nombres et les sommes d'argent (块/元, 毛, 钱) vus en Première
\item Les verbes modaux 想, 要, 应该 et 可以
\item La structure d'opinion simple 我觉得 + phrase
\item Le vocabulaire des achats et des prix (买, 卖, 贵, 便宜, 多少)
\item Les questions avec 吗, 呢, 怎么 et 为什么
\end{itemize}
\end{prerequis}

\newpage

\coursec{Vocabulaire du budget}

C'est la fin du mois à Yaoundé. Amina a reçu son argent de poche de 10 000 FCFA le premier du mois. Mais après les trajets en taxi, les beignets du matin et le crédit téléphonique, il ne lui reste plus rien dès le 20 ! Elle soupire : « Où est passé mon argent ? » Son correspondant chinois, Li Ming, lycéen à Pékin, lui répond par message : en Chine, les lycéens apprennent à gérer leur \zh{零花钱} (línghuāqián, argent de poche) avec un petit budget mensuel, et ils savent dire ce qui est cher, ce qui est bon marché, et ce qu'ils mettent de côté.

\textbf{Problématique :} comment présenter son budget, donner son avis sur les dépenses et débattre en chinois ? Tout commence par le vocabulaire : c'est ce que nous allons apprendre aujourd'hui.

\courssub{Observation : un texte pour découvrir les mots}

Avant toute liste de mots, lisons le message que Li Ming envoie à Amina. Soulignez mentalement tous les mots qui parlent d'argent.

\begin{textebox}[Le message de Li Ming]
我每个月有零花钱。我把钱花在交通和电话上，剩下的存起来。\\
Wǒ měi ge yuè yǒu línghuāqián. Wǒ bǎ qián huā zài jiāotōng hé diànhuà shang, shèngxià de cún qǐlai.\\
« Je reçois de l'argent de poche chaque mois. Je le dépense pour le transport et le téléphone, et je mets le reste de côté. »
\end{textebox}

Ce court texte contient déjà cinq mots-clés de la leçon : \zh{钱} (qián, l'argent), \zh{零花钱} (línghuāqián, l'argent de poche), \zh{花} (huā, dépenser), \zh{每个月} (měi ge yuè, chaque mois) et \zh{存} (cún, mettre de côté). Observons-les un par un.

\courssub{L'argent et l'argent de poche : 钱 et 零花钱}

\begin{definition}[钱 (qián) : l'argent]
\zh{钱} (qián), nom, signifie « l'argent » en général. Le caractère contient la clé du métal \zh{钅} à gauche (autrefois, la monnaie était en métal) : on l'écrit de gauche à droite, d'abord la clé, puis la partie droite. On dit \zh{钱} pour parler de l'argent en général, jamais pour un billet précis.
\end{definition}

\begin{exemplebox}[钱 en contexte]
\zh{我没有钱。} \emph{Wǒ méiyǒu qián.} « Je n'ai pas d'argent. »\\
\zh{这本书多少钱？} \emph{Zhè běn shū duōshao qián ?} « Combien coûte ce livre ? » (litt. « ce livre, combien d'argent ? »)
\end{exemplebox}

\begin{definition}[零花钱 (línghuāqián) : l'argent de poche]
\zh{零花钱} (línghuāqián) se décompose ainsi : \zh{零} (líng, zéro, petit), \zh{花} (huā, dépenser), \zh{钱} (qián, argent) : littéralement « l'argent des petites dépenses ». C'est exactement notre « argent de poche », celui que les parents donnent aux enfants et aux lycéens, à Yaoundé comme à Pékin.
\end{definition}

\begin{exemplebox}[零花钱 en contexte]
\zh{妈妈每个月给我零花钱。} \emph{Māma měi ge yuè gěi wǒ línghuāqián.} « Maman me donne de l'argent de poche chaque mois. »\\
\zh{你的零花钱够吗？} \emph{Nǐ de línghuāqián gòu ma ?} « Ton argent de poche suffit-il ? »
\end{exemplebox}

\courssub{Le budget et les dépenses : 预算 et 花费}

\begin{definition}[预算 (yùsuàn) : le budget]
\zh{预算} (yùsuàn), nom, désigne la somme d'argent \emph{prévue} pour des dépenses sur une période donnée (une semaine, un mois, une année). \zh{预} (yù) signifie « à l'avance » et \zh{算} (suàn) « calculer » : le budget, c'est donc littéralement « le calcul fait à l'avance ». On dit \zh{做预算} (zuò yùsuàn, faire un budget).
\end{definition}

\begin{exemplebox}[预算 en contexte]
\zh{我每个月做一个小预算。} \emph{Wǒ měi ge yuè zuò yí ge xiǎo yùsuàn.} « Je fais un petit budget chaque mois. »\\
\zh{这个月的预算是一万。} \emph{Zhè ge yuè de yùsuàn shì yíwàn.} « Le budget de ce mois est de dix mille. »
\end{exemplebox}

\begin{definition}[花费 (huāfèi) : les dépenses]
\zh{花费} (huāfèi), nom, désigne les dépenses, l'argent dépensé. Il est formé de \zh{花} (huā, dépenser) et \zh{费} (fèi, frais). Attention : c'est un nom, pas un verbe.
\end{definition}

\begin{exemplebox}[花费 en contexte]
\zh{交通的花费太多了。} \emph{Jiāotōng de huāfèi tài duō le.} « Les dépenses de transport sont trop importantes. »
\end{exemplebox}

\courssub{Les verbes de l'argent : 花, 省, 存, 浪费, 需要}

\begin{propriete}[花 (huā) : dépenser]
\zh{花} (huā) est un verbe qui signifie « dépenser » (de l'argent, mais aussi du temps). Structure : \textbf{Sujet + 花 + argent/temps (+ complément)}.\\
\zh{我花了一百块。} \emph{Wǒ huā le yìbǎi kuài.} « J'ai dépensé cent yuans. »\\
\zh{她花了很多时间学习。} \emph{Tā huā le hěn duō shíjiān xuéxí.} « Elle a passé beaucoup de temps à étudier. »
\end{propriete}

\begin{propriete}[把钱花在...上 : dépenser son argent pour...]
Structure très utile à l'oral : \textbf{把 + 钱 + 花在 + nom + 上}. La particule \zh{把} (bǎ) place l'objet (\zh{钱}) avant le verbe, et \zh{在...上} (zài...shang) introduit le domaine de la dépense.\\
\zh{我把钱花在交通上。} \emph{Wǒ bǎ qián huā zài jiāotōng shang.} « Je dépense mon argent pour le transport. »\\
\zh{Amina 把钱花在电话上。} \emph{Amina bǎ qián huā zài diànhuà shang.} « Amina dépense son argent pour le téléphone. »
\end{propriete}

\begin{propriete}[省 (shěng) et 存 (cún) : économiser et mettre de côté]
\zh{省} (shěng), verbe, signifie « économiser », c'est-à-dire dépenser moins : \zh{省钱} (shěng qián, économiser de l'argent). \zh{存} (cún), verbe, signifie « mettre de côté, déposer » (à la banque, dans une tirelire) : \zh{存钱} (cún qián, mettre de l'argent de côté). La nuance : on \zh{省} en dépensant moins, on \zh{存} en gardant ce qui reste.\\
\zh{我想省点钱。} \emph{Wǒ xiǎng shěng diǎn qián.} « Je veux économiser un peu d'argent. »\\
\zh{剩下的钱，我存起来。} \emph{Shèngxià de qián, wǒ cún qǐlai.} « L'argent qui reste, je le mets de côté. »
\end{propriete}

\begin{propriete}[浪费 (làngfèi) et 需要 (xūyào)]
\zh{浪费} (làngfèi), verbe et adjectif, signifie « gaspiller » : \zh{别浪费钱！} \emph{Bié làngfèi qián !} « Ne gaspille pas l'argent ! » \zh{需要} (xūyào), verbe, signifie « avoir besoin de » ; il se place devant un nom ou un verbe : \zh{我需要买一本书。} \emph{Wǒ xūyào mǎi yì běn shū.} « J'ai besoin d'acheter un livre. »
\end{propriete}

\courssub{Cher ou bon marché : 贵, 便宜, 价格}

\begin{exemplebox}[Deux réactions au marché de Mokolo]
\zh{这本书太贵了！} \emph{Zhè běn shū tài guì le !} « Ce livre est trop cher ! »\\
\zh{那个市场的东西很便宜。} \emph{Nà ge shìchǎng de dōngxi hěn piányi.} « Les produits de ce marché sont très bon marché. »
\end{exemplebox}

\begin{definition}[贵 (guì), 便宜 (piányi), 价格 (jiàgé)]
\zh{贵} (guì), adjectif : « cher ». Son contraire est \zh{便宜} (piányi), adjectif : « bon marché ». \zh{价格} (jiàgé), nom : « le prix ». On demande le prix avec \zh{价格是多少？} (jiàgé shì duōshao ?) ou plus simplement \zh{多少钱？} (duōshao qián ?).
\end{definition}

\begin{attention}[La prononciation de 便宜]
\zh{便宜} se prononce \textbf{piányi} (2\ieme{} ton + ton neutre) et non \emph{biànyí}. C'est l'erreur la plus fréquente des francophones, car le caractère \zh{便} se lit \emph{biàn} dans \zh{方便} (fāngbiàn, pratique). Retenez : pour les prix, c'est toujours \textbf{piányi}. Autre piège : ne dites pas \zh{价格很贵} à la place de \zh{东西很贵} si vous parlez de l'objet ; avec \zh{价格}, on dit \zh{价格很高} (le prix est élevé).
\end{attention}

\courssub{Le temps et les quantités : 每个月 et 一半}

\begin{definition}[每个月 (měi ge yuè) et 一半 (yíbàn)]
\zh{每个月} (měi ge yuè) signifie « chaque mois » : \zh{每} (měi, chaque) + classificateur \zh{个} + \zh{月} (yuè, mois). \zh{一半} (yíbàn) signifie « la moitié » : \zh{我把一半的钱存起来。} \emph{Wǒ bǎ yíbàn de qián cún qǐlai.} « Je mets la moitié de l'argent de côté. »
\end{definition}

\courssub{Carte mentale du lexique}

\begin{popfigure}
\begin{tikzpicture}[mindmap, grow cyclic, every node/.style={concept}, root concept/.append style={concept color=popblue, font=\large\bfseries}, level 1/.append style={level distance=3.4cm, sibling angle=90}, level 2/.append style={level distance=2.2cm, sibling angle=45, concept color=popgreenL}]
\node[root concept, align=center]{钱\\ qián\\ l'argent}
  child[concept color=poporangeL]{ node[align=center]{gagner\\ 得到 dédào} }
  child[concept color=poppinkL]{ node[align=center]{dépenser\\ 花 huā}
      child { node[align=center]{浪费 làngfèi\\ gaspiller} }
  }
  child[concept color=popgreenL]{ node[align=center]{économiser\\ 省 shěng / 存 cún} }
  child[concept color=popgoldL]{ node[align=center]{prix\\ 贵 guì / 便宜 piányi}
      child { node[align=center]{价格 jiàgé\\ le prix} }
  };
\end{tikzpicture}
\legende{Carte mentale du vocabulaire de l'argent autour du mot central 钱 (qián).}
\end{popfigure}

\courssub{Tableau récapitulatif du vocabulaire}

\begin{center}
\begin{tabularx}{\linewidth}{|X|X|X|X|}
\hline
\rowcolor{popblueL} Caractères & Pinyin & Nature & Traduction \\
\hline
钱 & qián & nom & l'argent \\
\hline
零花钱 & línghuāqián & nom & l'argent de poche \\
\hline
预算 & yùsuàn & nom & le budget \\
\hline
花费 & huāfèi & nom & les dépenses \\
\hline
花 & huā & verbe & dépenser \\
\hline
省 & shěng & verbe & économiser \\
\hline
存 & cún & verbe & mettre de côté \\
\hline
贵 & guì & adjectif & cher \\
\hline
便宜 & piányi & adjectif & bon marché \\
\hline
价格 & jiàgé & nom & le prix \\
\hline
需要 & xūyào & verbe & avoir besoin de \\
\hline
浪费 & làngfèi & verbe / adj. & gaspiller \\
\hline
每个月 & měi ge yuè & expression & chaque mois \\
\hline
一半 & yíbàn & nom & la moitié \\
\hline
\end{tabularx}
\end{center}

\courssub{Comparaison chinois / français}

\begin{center}
\begin{tabularx}{\linewidth}{|X|X|X|}
\hline
\rowcolor{popblueL} Notion & Français & Chinois \\
\hline
Article & « de l'argent » (article obligatoire) & 钱 seul, sans article : \zh{我没有钱。} \\
\hline
Verbe « dépenser » & « dépenser \emph{pour} le transport » & 花在...上 : \zh{花在交通上} (la préposition encadre le nom) \\
\hline
« Cher » & « ce livre \emph{est} cher » (verbe être) & pas de verbe être : \zh{这本书很贵。} \\
\hline
« Chaque mois » & « chaque » + nom directement & 每 + classificateur + nom : \zh{每个月} \\
\hline
\end{tabularx}
\end{center}

\courssub{Exercice résolu}

\begin{exoresolu}[Complétez avec le mot qui convient]
Complétez les phrases avec : 预算, 零花钱, 便宜, 存, 浪费.\\
1. 妈妈给我五千块 \_\_\_\_。\\
2. 这个东西很 \_\_\_\_，只要一百块。\\
3. 我把剩下的钱 \_\_\_\_ 起来。\\
4. 别 \_\_\_\_ 钱！\\
5. 这个月你的 \_\_\_\_ 是多少？
\tcblower
\textbf{Solution détaillée :}\\
1. \zh{零花钱} : la maman donne de l'argent de poche (5 000 FCFA). \emph{Māma gěi wǒ wǔqiān kuài línghuāqián.} « Maman me donne 5 000 d'argent de poche. »\\
2. \zh{便宜} : « seulement cent » indique un prix bas. \emph{Zhè ge dōngxi hěn piányi.} « Cet article est bon marché. »\\
3. \zh{存} : on met de côté ce qui reste. \emph{Wǒ bǎ shèngxià de qián cún qǐlai.}\\
4. \zh{浪费} : après \zh{别} (ne... pas), il faut un verbe d'action négative à éviter : gaspiller. \emph{Bié làngfèi qián !}\\
5. \zh{预算} : on demande le montant prévu pour le mois. \emph{Zhè ge yuè nǐ de yùsuàn shì duōshao ?} « Quel est ton budget ce mois-ci ? »
\end{exoresolu}

\begin{attention}[Pièges fréquents des francophones]
\begin{itemize}
\item Ne traduisez pas « cher » par \zh{贵} suivi de \zh{的} devant un nom : on dit \zh{很贵的书} (un livre cher), mais jamais \zh{贵的书} sans \zh{很} à l'oral courant.
\item \zh{花} ne s'utilise pas seul pour « acheter » : « j'achète un livre » se dit \zh{我买书}, pas \zh{我花书}.
\item \zh{一半} se place avant \zh{的} + nom : \zh{一半的钱} (la moitié de l'argent), et non après le nom comme en français.
\end{itemize}
\end{attention}

\begin{savaistu}[Les enveloppes rouges du Nouvel An]
En Chine, lors de la fête du Printemps (le Nouvel An chinois), les enfants reçoivent des \zh{红包} (hóngbāo), des enveloppes rouges contenant de l'argent, offertes par les parents et les grands-parents. Le rouge est la couleur du bonheur et de la chance. On encourage les enfants à ne pas tout dépenser : une partie du \zh{红包} est mise de côté (\zh{存起来}), première leçon de gestion de budget ! Au Cameroun, les cadeaux d'argent des fêtes de fin d'année jouent un rôle comparable : c'est souvent là qu'Amina reçoit son premier « capital » de l'année.
\end{savaistu}

\begin{exercice}[À vous de jouer]
Traduisez en chinois : 1. « Je reçois de l'argent de poche chaque mois. » 2. « Ce téléphone est trop cher ! » 3. « Je dépense mon argent pour le transport. » 4. « Je mets la moitié de l'argent de côté. » Puis, à l'oral, présentez en trois phrases votre propre budget mensuel en utilisant au moins quatre mots du tableau.
\end{exercice}

\begin{corrige}[Exercice « À vous de jouer »]
1. \zh{我每个月有零花钱。} \emph{Wǒ měi ge yuè yǒu línghuāqián.}\\
2. \zh{这个电话太贵了！} \emph{Zhè ge diànhuà tài guì le !}\\
3. \zh{我把钱花在交通上。} \emph{Wǒ bǎ qián huā zài jiāotōng shang.}\\
4. \zh{我把一半的钱存起来。} \emph{Wǒ bǎ yíbàn de qián cún qǐlai.}\\
Pour la production orale, modèle possible : \zh{我每个月有一万块零花钱。我把钱花在交通和电话上。剩下的钱，我存起来。} Vérifiez la prononciation de \zh{便宜} (piányi) et les tons de \zh{预算} (yùsuàn, 4\ieme{} + 4\ieme{} ton).
\end{corrige}

\coursec{Dialogues modèles}

Dans la section précédente, nous avons construit le vocabulaire du budget : \zh{零花钱} (línghuāqián, argent de poche), \zh{存钱} (cún qián, économiser), \zh{省钱} (shěng qián, faire des économies), \zh{花钱} (huā qián, dépenser). Mais un vocabulaire qui reste dans un tableau ne sert à rien : il faut le faire \emph{vivre} dans la bouche. C'est exactement le rôle des dialogues modèles. Un dialogue modèle est une petite conversation authentique, courte, que l'on mémorise, que l'on joue, puis que l'on transforme. C'est la méthode la plus efficace pour passer du « je connais les mots » au « je sais parler ». Dans cette section, nous allons observer trois dialogues, comprendre leurs structures, puis apprendre à les faire nôtres.

\courssub{Dialogue 1 : parler de son argent de poche}

\begin{textebox}[Dialogue 1 — L'argent de poche]
A ：你每个月有多少零花钱？\\
\emph{Nǐ měige yuè yǒu duōshao línghuāqián?}\\
« Combien d'argent de poche as-tu chaque mois ? »\\[4pt]
B ：大概两百块。你呢？\\
\emph{Dàgài liǎng bǎi kuài. Nǐ ne?}\\
« À peu près deux cents yuans. Et toi ? »\\[4pt]
A ：我也有两百块。我把钱花在吃饭和交通上。\\
\emph{Wǒ yě yǒu liǎng bǎi kuài. Wǒ bǎ qián huā zài chīfàn hé jiāotōng shàng.}\\
« Moi aussi, j'ai deux cents yuans. Je dépense mon argent pour la nourriture et le transport. »
\end{textebox}

\begin{propriete}[Structure : Demander une quantité avec 多少]
Pour demander une quantité indéterminée ou importante, on utilise \cle{Sujet + 有 + 多少 + Nom ?}.
\zh{你每个月有多少零花钱？} \emph{Nǐ měige yuè yǒu duōshao línghuāqián?}
Le mot \zh{多少} (duōshao) ne prend \textbf{pas} de classificateur après lui (contrairement à \zh{几} jǐ). On dit \zh{多少钱}, \zh{多少人}, \zh{多少书}.
\end{propriete}

\begin{propriete}[Répondre par une approximation : 大概 + chiffre]
\zh{大概} (dàgài, « à peu près, environ ») se place \emph{avant} le groupe numéral.
\zh{大概两百块。} \emph{Dàgài liǎng bǎi kuài.}
\zh{块} (kuài) est le classificateur familier de la monnaie (équivalent oral de \zh{元} yuán).
\end{propriete}

\begin{propriete}[Relancer la conversation : 你呢？]
La particule \zh{呢} (ne) placée après un pronom ou un nom permet de retourner la question à l'interlocuteur sans la répéter : \zh{你呢？} (Et toi ?), \zh{他呢？} (Et lui ?). C'est un outil de fluidité conversationnelle indispensable.
\end{propriete}

\begin{propriete}[La structure 把 + objet + 花在...上]
Pour dire « dépenser quelque chose \emph{pour} quelque chose » en insistant sur la destination de l'objet, le chinois utilise la construction en \cle{把} :
\begin{center}
Sujet + 把 + Objet (argent) + 花 + 在 + Poste de dépense + 上
\end{center}
\zh{我把钱花在吃饭和交通上。} \emph{Wǒ bǎ qián huā zài chīfàn hé jiāotōng shàng.} « Je dépense mon argent pour la nourriture et le transport. »\\
La préposition \zh{在} (zài) introduit le domaine, et \zh{上} (shàng) le ferme : \zh{在...上}. Sans \zh{把}, on dit simplement \zh{我花钱吃饭} (je dépense de l'argent pour manger), mais la structure en \zh{把} met l'accent sur \emph{la gestion de l'argent}.
\end{propriete}

\begin{attention}[两 (liǎng) ou 二 (èr) ? — Erreur classique]
Devant un classificateur ou une mesure, « deux » se dit obligatoirement \cle{两} (liǎng), jamais \zh{二} (èr).
\begin{itemize}
\item \textbf{Correct :} \zh{两百块} (liǎng bǎi kuài), \zh{两个人} (liǎng ge rén), \zh{两本书} (liǎng běn shū).
\item \textbf{Usage de 二 :} numéros de série, mathématiques, ordinaux : \zh{二月} (février), \zh{二楼} (2e étage), \zh{第二} (deuxième), \zh{一二三} (compter).
\end{itemize}
Dire \zh{二块} est une erreur typique des francophones car le français n'a qu'un seul mot « deux ».
\end{attention}

\courssub{Dialogue 2 : donner son avis et nuancer}

\begin{textebox}[Dialogue 2 — Faut-il économiser ?]
A ：我觉得学生应该学会省钱。\\
\emph{Wǒ juéde xuésheng yīnggāi xuéhuì shěng qián.}\\
« Je pense que les élèves devraient apprendre à économiser. »\\[4pt]
B ：我同意，但是不能什么都不买。\\
\emph{Wǒ tóngyì, dànshì bù néng shénme dōu bù mǎi.}\\
« Je suis d'accord, mais on ne peut pas ne rien acheter du tout. »
\end{textebox}

\begin{propriete}[Exprimer son opinion : 我觉得 + phrase]
\zh{我觉得} (wǒ juéde, « je pense que / je trouve que ») introduit une opinion personnelle. Il est suivi d'une phrase complète.
\zh{我觉得学生应该学会省钱。}
Les verbes modaux nuancent l'opinion :
\begin{itemize}
\item \zh{应该} (yīnggāi) : devoir moral, conseil (« devrait »).
\item \zh{能} (néng) : capacité, possibilité objective (« pouvoir »).
\item \zh{会} (huì) : savoir-faire acquis (« savoir »).
\end{itemize}
Ici \zh{学会} (xuéhuì) = « apprendre à / arriver à savoir faire ».
\end{propriete}

\begin{propriete}[Désaccord poli : 我同意，但是...]
La réponse de B est le modèle standard du débat courtois en chinois :
1. Valider : \zh{我同意} (Wǒ tóngyì, « Je suis d'accord »).
2. Nuancer : \zh{但是} (dànshì, « mais ») + objection argumentée.
Ne jamais opposer un « non » sec (\zh{不对} ou \zh{不}) en ouverture.
\end{propriete}

\begin{propriete}[Négation totale : 什么都不... / 谁都不...]
\zh{什么都不买} (shénme dōu bù mǎi) = « ne rien acheter du tout ».
Structure : \cle{Mot-interrogatif + 都 + Négation + Verbe}.
Le mot interrogatif perd sa valeur de question pour devenir un indéfini totalisant (« quoi que ce soit », « qui que ce soit »). \zh{都} (dōu) porte la totalisation.
\begin{itemize}
\item \zh{我什么都不知道。} (Wǒ shénme dōu bù zhīdào.) « Je ne sais rien du tout. »
\item \zh{他哪儿都不去。} (Tā nǎr dōu bù qù.) « Il ne va nulle part. »
\end{itemize}
\end{propriete}

\courssub{Dialogue 3 : relancer et conclure}

\begin{textebox}[Dialogue 3 — Demander l'avis de l'autre]
A ：你觉得怎么样？\\
\emph{Nǐ juéde zěnmeyàng?}\\
« Qu'en penses-tu ? »\\[4pt]
B ：对我来说，存钱很重要。\\
\emph{Duì wǒ láishuō, cún qián hěn zhòngyào.}\\
« Pour moi, mettre de l'argent de côté est très important. »
\end{textebox}

\begin{propriete}[Deux formules fonctionnelles à automatiser]
\begin{enumerate}
\item \zh{你觉得怎么样？} (Nǐ juéde zěnmeyàng?) — « Qu'en penses-tu ? » (litt. « Tu penses comment ? »).
  On peut préciser le sujet : \zh{你觉得省钱怎么样？} (Que penses-tu d'économiser ?).
\item \cle{对...来说} (duì...láishuō) — « Du point de vue de..., pour... ».
  Structure : \zh{对} + Personne/Groupe + \zh{来说} + Commentaire.
  \zh{对学生来说，时间很重要。} (Duì xuésheng láishuō, shíjiān hěn zhòngyào.) « Pour les élèves, le temps est important. »
\end{enumerate}
\end{propriete}

\begin{exemplebox}[Exemples traduits à réutiliser à l'oral]
\zh{我把钱花在书上。} \emph{Wǒ bǎ qián huā zài shū shàng.} « Je dépense mon argent pour les livres. »\\[4pt]
\zh{你应该存一点儿钱。} \emph{Nǐ yīnggāi cún yìdiǎnr qián.} « Tu devrais mettre un peu d'argent de côté. »\\[4pt]
\zh{对我来说，交通太贵了。} \emph{Duì wǒ láishuō, jiāotōng tài guì le.} « Pour moi, le transport est trop cher. »\\[4pt]
\zh{我同意你的看法。} \emph{Wǒ tóngyì nǐ de kànfǎ.} « Je suis d'accord avec ton point de vue. »
\end{exemplebox}

\courssub{Le dialogue complet à lire}

Voici maintenant le dialogue complet, qui enchaîne les trois mouvements : \textbf{présenter} sa situation, \textbf{donner son avis}, \textbf{relancer} et répondre. Lisez-le d'abord en silence, puis à voix haute.

\begin{textebox}[Dialogue complet — Mon budget]
A ：你每个月有多少零花钱？\\
\emph{Nǐ měige yuè yǒu duōshao línghuāqián?}\\
« Combien d'argent de poche as-tu chaque mois ? »\\[4pt]
B ：大概两百块。你呢？\\
\emph{Dàgài liǎng bǎi kuài. Nǐ ne?}\\
« À peu près deux cents yuans. Et toi ? »\\[4pt]
A ：我也有两百块。我把钱花在吃饭和交通上。\\
\emph{Wǒ yě yǒu liǎng bǎi kuài. Wǒ bǎ qián huā zài chīfàn hé jiāotōng shàng.}\\
« Moi aussi, j'ai deux cents yuans. Je dépense mon argent pour la nourriture et le transport. »\\[4pt]
B ：我觉得你应该存一点儿。\\
\emph{Wǒ juéde nǐ yīnggāi cún yìdiǎnr.}\\
« Je pense que tu devrais mettre un peu d'argent de côté. »\\[4pt]
A ：我同意！\\
\emph{Wǒ tóngyì!}\\
« Je suis d'accord ! »
\end{textebox}

\begin{center}
\begin{tabularx}{\linewidth}{|X|X|X|X|}
\hline
\rowcolor{popblueL} Caractères & Pinyin & Nature & Traduction \\
\hline
零花钱 & línghuāqián & nom & argent de poche \\
\hline
大概 & dàgài & adverbe & à peu près, environ \\
\hline
块 & kuài & classificateur & yuan (familier) \\
\hline
花 & huā & verbe & dépenser \\
\hline
交通 & jiāotōng & nom & transport \\
\hline
觉得 & juéde & verbe & penser, trouver \\
\hline
应该 & yīnggāi & verbe modal & devoir (conseil) \\
\hline
存 & cún & verbe & économiser, mettre de côté \\
\hline
一点儿 & yìdiǎnr & quantificateur & un peu \\
\hline
同意 & tóngyì & verbe & être d'accord \\
\hline
\end{tabularx}
\end{center}

\courssub{La logique d'un échange : présenter, opiner, relancer}

Tout dialogue d'opinion en chinois suit une chorégraphie en trois temps. Le schéma suivant la résume ; gardez-le en tête pendant les jeux de rôle.

\begin{popfigure}
\begin{tikzpicture}[node distance=1.6cm]
\node[draw=popblue, fill=popblueL, rounded corners, thick, align=center, minimum width=3.4cm, minimum height=1.4cm] (pres) {\textbf{Présenter}\\ \zh{我有两百块。}\\ \emph{Wǒ yǒu...}};
\node[draw=poporange, fill=poporangeL, rounded corners, thick, align=center, minimum width=3.4cm, minimum height=1.4cm, right=of pres] (opin) {\textbf{Opiner}\\ \zh{我觉得应该省钱。}\\ \emph{Wǒ juéde...}};
\node[draw=popgreen, fill=popgreenL, rounded corners, thick, align=center, minimum width=3.4cm, minimum height=1.4cm, right=of opin] (rel) {\textbf{Relancer}\\ \zh{你呢？你觉得呢？}\\ \emph{Nǐ ne?}};
\draw[-{Stealth}, very thick, popdark] (pres) -- (opin);
\draw[-{Stealth}, very thick, popdark] (opin) -- (rel);
\draw[-{Stealth}, very thick, poppurple, dashed] (rel.south) to[bend left=25] node[below, align=center, font=\small]{l'autre répond à son tour} (pres.south);
\end{tikzpicture}
\legende{Les trois mouvements d'un dialogue d'opinion : présenter sa situation, donner son avis, relancer l'interlocuteur — et la conversation repart.}
\end{popfigure}

\begin{methode}[Mémoriser un dialogue en cinq étapes]
\begin{enumerate}
\item \textbf{Lire} : lisez le dialogue en silence, repérez les mots connus et les structures.
\item \textbf{Répéter en chœur} : la classe répète après le professeur, phrase par phrase, en soignant les tons.
\item \textbf{Jouer avec le texte} : en binômes, jouez le dialogue en lisant, avec l'intonation et les gestes.
\item \textbf{Jouer sans le texte} : fermez le manuel ; ne gardez que le schéma présenter $\rightarrow$ opiner $\rightarrow$ relancer.
\item \textbf{Substituer} : remplacez les sommes (\zh{两百块} $\rightarrow$ \zh{三百块}, \zh{五百块}) et les postes de dépenses (\zh{吃饭} $\rightarrow$ \zh{买书}, \zh{打电话}, \zh{买衣服}).
\end{enumerate}
\end{methode}

\begin{exoresolu}[Substitution guidée]
Énoncé : Remplacez dans la phrase \zh{我把钱花在吃饭和交通上} les deux postes de dépenses par « les livres » (\zh{书}) et « le téléphone » (\zh{电话}), puis traduisez.
\tcblower
Solution : La structure \zh{把} + objet + \zh{花在} + A + \zh{和} + B + \zh{上} reste identique ; seuls les noms de postes changent. On obtient :\\
\zh{我把钱花在书和电话上。}\\
\emph{Wǒ bǎ qián huā zài shū hé diànhuà shàng.}\\
« Je dépense mon argent pour les livres et le téléphone. »\\
\emph{Attention :} \zh{上} reste toujours en fin de cadre, après le dernier poste cité.
\end{exoresolu}

\begin{attention}[Pièges de prononciation dans ces dialogues]
\begin{itemize}
\item \zh{多少} (duōshao) : 1er ton + ton léger (pas de marque sur \zh{shao}) ; ne dites pas « duōshǎo » avec deux tons pleins.
\item \zh{觉得} (juéde) : 2e ton + ton léger sur \zh{得} ; un ton plein alourdit la phrase.
\item \zh{一点儿} (yìdiǎnr) : \zh{一} (yī) subit la sandhi et se prononce 4e ton (yì) devant le 3e ton de \zh{点} (diǎn). Le \zh{儿} (r) se fusionne (érhuà).
\item \zh{同意} (tóngyì) : 2e ton puis 4e ton ; les francophones ont tendance à aplatir le 2e ton (montant).
\item \zh{大概} (dàgài) : deux 4e tons (chute nette deux fois).
\end{itemize}
\end{attention}

\begin{experience}[Lecture en chœur puis en binômes]
\textbf{Étape 1 (Chorale) :} Le professeur lit le dialogue complet, la classe répète phrase par phrase. Ensuite, lecture par rôles : la moitié gauche de la classe joue A, la moitié droite joue B, puis on échange.
\textbf{Étape 2 (Binômes) :} Jouez le dialogue trois fois de suite :
\begin{enumerate}
\item Avec le texte sous les yeux.
\item Sans le texte, en s'appuyant sur le schéma « Présenter $\rightarrow$ Opiner $\rightarrow$ Relancer ».
\item En substitution libre : sommes (\zh{一百块}, \zh{三百块}, \zh{五百块}) et postes (\zh{吃饭}, \zh{交通}, \zh{买书}, \zh{买衣服}, \zh{打电话}, \zh{上网}).
\end{enumerate}
Terminez chaque échange par une relance : \zh{你呢？} ou \zh{你觉得怎么样？}
\end{experience}

\begin{exercice}[À vous de parler]
En binômes, construisez votre propre dialogue en suivant le schéma imposé :
\begin{enumerate}
\item Question sur l'argent de poche (\zh{多少}).
\item Réponse avec \zh{大概} + somme + relance \zh{你呢？}.
\item Présentation de vos dépenses avec la structure en \zh{把} (\zh{花在...上}).
\item Avis avec \zh{我觉得...应该...}.
\item Accord ou désaccord poli (\zh{我同意，但是...} ou \zh{我不同意，因为...}).
\end{enumerate}
Présentez votre dialogue devant la classe (1 à 2 minutes par binôme).
\end{exercice}

\begin{corrige}[Exercice — À vous de parler]
\textbf{Modèle de production attendue (exemple) :}\\[4pt]
A ：你每个月有多少零花钱？ \emph{(Nǐ měige yuè yǒu duōshao línghuāqián?)}\\[4pt]
B ：大概三百块。你呢？ \emph{(Dàgài sān bǎi kuài. Nǐ ne?)}\\[4pt]
A ：我有两百块。我把钱花在买书和交通上。 \emph{(Wǒ yǒu liǎng bǎi kuài. Wǒ bǎ qián huā zài mǎi shū hé jiāotōng shàng.)}\\[4pt]
B ：我觉得你应该存一点儿钱。 \emph{(Wǒ juéde nǐ yīnggāi cún yìdiǎnr qián.)}\\[4pt]
A ：我同意，但是我也想买衣服。 \emph{(Wǒ tóngyì, dànshì wǒ yě xiǎng mǎi yīfu.)}\\[8pt]
\textbf{Critères d'évaluation (grille orale) :}
\begin{itemize}
\item \textbf{Prononciation / Tons} (exactitude des tons, sandhi de \zh{一}, tons légers) : 5 pts
\item \textbf{Structures} (\zh{把} correct, \zh{觉得/应该}, \zh{对...来说} si utilisé, \zh{两} vs \zh{二}) : 5 pts
\item \textbf{Interaction} (relance \zh{你呢？} / \zh{你觉得怎么样？}, fluidité, réactivité) : 5 pts
\item \textbf{Vocabulaire} (richesse des postes de dépenses, connecteurs \zh{但是/因为}) : 5 pts
\end{itemize}
Total : 20 pts.
\end{corrige}

\begin{aretenir}[L'essentiel des dialogues modèles]
\begin{itemize}
\item Demander : \zh{你每个月有多少零花钱？} — Répondre : \zh{大概两百块。}
\item Relancer : \zh{你呢？} / \zh{你觉得怎么样？}
\item Dépenser (focus objet) : \zh{我把钱花在...上。} (structure en \zh{把})
\item Opiner : \zh{我觉得...应该...} ; Nuancer : \zh{我同意，但是...}
\item Négation totale : \zh{什么都不...} / \zh{谁都不...}
\item \textbf{Deux + classificateur = \zh{两} (liǎng), jamais \zh{二} (èr).}
\end{itemize}
\end{aretenir}

\coursec{Vocabulaire du budget}

\begin{definition}[Mots-clés du thème « Gestion du budget »]
Le lexique suivant est indispensable pour comprendre, parler et écrire sur le budget personnel ou familial. Il correspond au niveau HSK 3--4 (niveau Terminale). Mémorisez les caractères, le pinyin (tons inclus), la nature grammaticale et le sens.
\end{definition}

\begin{center}
\begin{tabularx}{\linewidth}{|X|X|X|X|}
\hline
\rowcolor{popblueL} Caractères & Pinyin & Nature & Traduction \\
\hline
预算 & yùsuàn & n. & budget \\
收入 & shōurù & n. & revenus, rentrée d'argent \\
支出 & zhīchū & n. & dépenses, sortie d'argent \\
工资 & gōngzī & n. & salaire \\
生活费 & shēnghuófèi & n. & frais de vie, argent pour la vie courante \\
零花钱 & línghuāqián & n. & argent de poche \\
节约 & jiéyuē & v./adj. & économiser / économe \\
浪费 & làngfèi & v./adj. & gaspiller / gaspilleur \\
存钱 & cúnqián & v. & mettre de côté, épargner \\
储蓄 & chǔxù & n. & épargne \\
管理 & guǎnlǐ & v. & gérer, administrer \\
负责 & fùzé & v. & être responsable de, s'occuper de \\
家庭 & jiātíng & n. & famille, foyer \\
成员 & chéngyuán & n. & membre (de la famille, d'un groupe) \\
困难 & kùnnan & adj. & difficile \\
容易 & róngyì & adj. & facile \\
重要 & zhòngyào & adj. & important \\
建议 & jiànyì & v./n. & suggérer / suggestion \\
计划 & jìhuà & v./n. & planifier / plan \\
记账 & jìzhàng & v. & tenir les comptes, noter les dépenses \\
\hline
\end{tabularx}
\end{center}

\begin{exemplebox}[Mise en contexte : phrases types]
\zh{爸爸每个月的工资是家里的主要收入。} \emph{Bàba měi ge yuè de gōngzī shì jiā lǐ de zhǔyào shōurù.} « Le salaire mensuel de papa est le revenu principal de la maison. »\\
\zh{我们要学会管理支出，不能浪费钱。} \emph{Wǒmen yào xuéhuì guǎnlǐ zhīchū, bùnéng làngfèi qián.} « Nous devons apprendre à gérer les dépenses, ne pas gaspiller l'argent. »\\
\zh{妈妈建议我每个月存一点钱。} \emph{Māma jiànyì wǒ měi ge yuè cún yìdiǎn qián.} « Maman me suggère d'épargner un peu d'argent chaque mois. »\\
\zh{记账是一个好习惯，对预算很有帮助。} \emph{Jìzhàng shì yí ge hǎo xíguàn, duì yùsuàn hěn yǒu bāngzhù.} « Tenir les comptes est une bonne habitude, très utile pour le budget. »
\end{exemplebox}

\begin{attention}[Pièges lexicaux fréquents]
\begin{itemize}
\item \zh{收入} (rentrées) vs \zh{支出} (sorties) : ne confondez pas les radicaux \zh{收} (recevoir) et \zh{支} (payer/dépenser).
\item \zh{零花钱} (argent de poche) : \zh{零} signifie « fragmenté, petit » ; \zh{花} « dépenser » ; \zh{钱} « argent ». Littéralement : « argent pour petites dépenses ».
\item \zh{节约} (économiser, vertueux) vs \zh{省钱} (coûter moins cher / faire des économies). \zh{省钱} peut être sujet : \zh{这个办法省钱} (Cette méthode fait économiser de l'argent). \zh{节约} qualifie une personne ou une action : \zh{他很节约} (Il est économe).
\item \zh{存钱} (verbe d'action : mettre à la banque) vs \zh{储蓄} (nom : l'épargne, le capital accumulé).
\end{itemize}
\end{attention}

\coursec{Dialogues modèles}

\begin{textebox}[Dialogue 1 : Deux camarades parlent de l'argent de poche (informel)]
A : 嘿，你每个月有多少零花钱？\\
B : 两百元。你呢？\\
A : 我有一百五。我觉得不太够用，想找个兼职。\\
B : 兼职挺好的，但是别影响学习。我建议你先记账，看看钱花哪儿了。\\
A : 好主意！我今晚就开始。\\
\emph{A : Hēi, nǐ měi ge yuè yǒu duōshao línghuāqián?}\\
\emph{B : Liǎng bǎi yuán. Nǐ ne?}\\
\emph{A : Wǒ yǒu yì bǎi wǔ. Wǒ juéde bú tài gòuyòng, xiǎng zhǎo ge jiānzhí.}\\
\emph{B : Jiānzhí tǐng hǎo de, dànshì bié yǐngxiǎng xuéxí. Wǒ jiànyì nǐ xiān jìzhàng, kànkan qián huā nǎr le.}\\
\emph{A : Hǎo zhǔyi! Wǒ jīnwǎn jiù kāishǐ.}\\
« A : Hé, combien as-tu d'argent de poche par mois ?\\
B : 200 yuans. Et toi ?\\
A : J'en ai 150. Je trouve que ce n'est pas assez, je veux trouver un job étudiant.\\
B : Le job c'est bien, mais ne laisse pas ça affecter tes études. Je te conseille de d'abord tenir tes comptes, voir où part l'argent.\\
A : Bonne idée ! Je commence ce soir. »
\end{textebox}

\begin{textebox}[Dialogue 2 : Exposé formel en classe — Le budget de ma famille]
大家好！\textbf{我先介绍一下}我们家的预算。我们家有五口人：爸爸、妈妈、两个弟弟和我。爸爸每个月发工资，妈妈负责管理家里的钱。\textbf{然后}，\textbf{我觉得}管理预算不太容易，因为东西越来越贵。\textbf{对我来说}，省钱比花钱难。有的人说学生不需要管钱，\textbf{可是}我\textbf{不太同意}：学生也应该学习节约。\textbf{最后}，\textbf{所以，我认为}每个人都应该有一个简单的预算。\textbf{我的介绍完了}，谢谢大家！\\
\emph{Dàjiā hǎo! Wǒ xiān jièshào yíxià wǒmen jiā de yùsuàn. Wǒmen jiā yǒu wǔ kǒu rén: bàba, māma, liǎng ge dìdi hé wǒ. Bàba měi ge yuè fā gōngzī, māma fùzé guǎnlǐ jiā lǐ de qián. Ránhòu, wǒ juéde guǎnlǐ yùsuàn bú tài róngyì, yīnwèi dōngxi yuè lái yuè guì. Duì wǒ láishuō, shěng qián bǐ huā qián nán. Yǒu de rén shuō xuésheng bù xūyào guǎn qián, kěshì wǒ bú tài tóngyì: xuésheng yě yīnggāi xuéxí jiéyuē. Zuìhòu, suǒyǐ, wǒ rènwéi měi ge rén dōu yīnggāi yǒu yí ge jiǎndān de yùsuàn. Wǒ de jièshào wán le, xièxie dàjiā!}\\
« Bonjour à tous ! Je vais d'abord présenter le budget de notre famille. Nous sommes cinq : papa, maman, mes deux petits frères et moi. Papa reçoit son salaire chaque mois, maman gère l'argent de la maison. Ensuite, je trouve que gérer un budget n'est pas très facile, car les choses deviennent de plus en plus chères. Pour moi, économiser est plus difficile que dépenser. Certains disent que les élèves n'ont pas besoin de gérer l'argent, mais je ne suis pas tout à fait d'accord : les élèves aussi devraient apprendre à économiser. Enfin, donc, je pense que chacun devrait avoir un budget simple. Mon exposé est terminé, merci à tous ! »
\end{textebox}

\coursec{Formules fonctionnelles}

Dans la section précédente, nous avons écouté et lu des dialogues modèles sur la gestion du budget familial. Vous avez remarqué que les locuteurs n'improvisent jamais totalement : ils s'appuient sur des \cle{formules toutes faites} pour ouvrir la parole, donner leur avis, réagir et conclure. Ces formules, appelées \cle{formules fonctionnelles}, sont les outils de base de tout exposé et de tout débat en chinois. Les connaître par cœur, c'est gagner en fluidité, en confiance et en points à l'examen. Cette section les présente une par une, avec leur structure, leur emploi et leurs pièges.

\courssub{Observer avant de mémoriser : un mini-exposé modèle}

Lisons d'abord ce court exposé d'un élève de Terminale sur le budget de sa famille. Repérez les formules en gras : elles seront expliquées ensuite.

\textbf{Questions d'observation :}
\begin{enumerate}
\item Quelle formule ouvre l'exposé ? Quelle formule le ferme ?
\item Combien de formules servent à donner un avis personnel ? Lesquelles ?
\item Quels mots relient les étapes du discours ?
\end{enumerate}

\courssub{Présenter : ouvrir la parole}

\begin{definition}[La formule d'ouverture standard]
\zh{我先介绍一下...} (\emph{wǒ xiān jièshào yíxià...}) signifie « je vais d'abord présenter... ». C'est la formule standard pour commencer un exposé.
\end{definition}

Décomposons la structure : \zh{我} (je) + \zh{先} (d'abord) + \zh{介绍} (présenter) + \zh{一下} (un peu, marqueur d'aspect délimitatif qui adoucit le verbe). Le mot \zh{一下} est essentiel : il rend le ton modeste et naturel, comme le « rapidement » ou le « un peu » du français oral.

\begin{exemplebox}[Varier l'ouverture selon le contexte]
\zh{我先介绍一下我的家庭预算。} \emph{Wǒ xiān jièshào yíxià wǒ de jiātíng yùsuàn.} « Je vais d'abord présenter le budget de ma famille. »\\
\zh{今天我想谈谈学生的零花钱。} \emph{Jīntiān wǒ xiǎng tántan xuésheng de línghuāqián.} « Aujourd'hui, je voudrais parler de l'argent de poche des élèves. »\\
\zh{我的题目是：怎么管理生活费？} \emph{Wǒ de tímù shì: zěnme guǎnlǐ shēnghuófèi?} « Mon sujet est : comment gérer ses frais de vie ? »
\end{exemplebox}

\courssub{Donner son avis : 我觉得, 我认为, 对我来说}

\begin{propriete}[Exprimer une opinion personnelle : structures]
\zh{我觉得} (\emph{wǒ juéde}, « je trouve que / j'ai l'impression que ») et \zh{我认为} (\emph{wǒ rènwéi}, « je pense que / j'estime que ») se construisent tous les deux avec une \textbf{phrase complète} (sujet + prédicat) après eux :\\
\cle{Sujet + 觉得 / 认为 + [Sujet + Verbe / Adjectif + Complément]}.\\
\zh{我觉得} est courant à l'oral, subjectif, basé sur le ressenti. \zh{我认为} est un peu plus \textbf{formel}, basé sur le raisonnement ; il convient parfaitement à l'exposé à l'examen.\\
La formule \zh{对我来说} (\emph{duì wǒ láishuō}, « pour moi, en ce qui me concerne ») se place en \textbf{tête de phrase} (position de thème), suivie d'une virgule, avant le commentaire.
\end{propriete}

\begin{exemplebox}[Trois façons de donner son avis]
\zh{我觉得管理预算很重要。} \emph{Wǒ juéde guǎnlǐ yùsuàn hěn zhòngyào.} « Je trouve que gérer un budget est très important. »\\
\zh{我认为学生应该学习节约。} \emph{Wǒ rènwéi xuésheng yīnggāi xuéxí jiéyuē.} « Je pense que les élèves devraient apprendre à économiser. »\\
\zh{对我来说，省钱比花钱难。} \emph{Duì wǒ láishuō, shěng qián bǐ huā qián nán.} « Pour moi, économiser est plus difficile que dépenser. »
\end{exemplebox}

\begin{attention}[Piège majeur : pas de verbe 是 après 觉得 et 认为 devant un adjectif]
En français on dit « je pense que c'\emph{est} important », et le francophone traduit souvent mot à mot : *\zh{我觉得是重要}. \textbf{C'est faux.} En chinois, l'adjectif seul joue le rôle de prédicat (verbe d'état) : \zh{我觉得管理预算很重要}.\\
On ne met \zh{是} que devant un \textbf{nom} (nominalisation) : \zh{我认为这是一个好办法。} (\emph{Wǒ rènwéi zhè shì yí ge hǎo bànfǎ.}) « Je pense que c'est une bonne méthode. »
\end{attention}

\courssub{Exprimer l'accord et le désaccord}

\begin{definition}[Accord et désaccord : nuances de politesse]
Accord : \zh{我同意} (\emph{wǒ tóngyì}, « je suis d'accord ») ; \zh{你说得对} (\emph{nǐ shuō de duì}, « tu as raison / ce que tu dis est juste »).\\
Désaccord poli : \zh{我不太同意} (\emph{wǒ bú tài tóngyì}, « je ne suis pas tout à fait d'accord »), souvent suivi de \zh{可是} (\emph{kěshì}) ou \zh{不过} (\emph{búguò}, « cependant / mais »).
\end{definition}

En chinois comme en français, un désaccord brutal est mal perçu. La culture chinoise valorise l'harmonie du groupe (\zh{和谐}) : on adoucit toujours avec \zh{不太} (« pas tout à fait ») plutôt que \zh{不同意} sec, et on enchaîne avec une concession introduite par \zh{可是} ou \zh{不过}.

\begin{exemplebox}[Réagir dans un débat]
\zh{我同意你的看法。} \emph{Wǒ tóngyì nǐ de kànfǎ.} « Je suis d'accord avec ton point de vue. »\\
\zh{你说得对，钱不容易管理。} \emph{Nǐ shuō de duì, qián bù róngyì guǎnlǐ.} « Tu as raison, l'argent n'est pas facile à gérer. »\\
\zh{我不太同意你的看法。} \emph{Wǒ bú tài tóngyì nǐ de kànfǎ.} « Je ne suis pas tout à fait d'accord avec ton point de vue. »\\
\zh{你说得有道理，不过我觉得节约也很重要。} \emph{Nǐ shuō de yǒu dàolǐ, búguò wǒ juéde jiéyuē yě hěn zhòngyào.} « Ce que tu dis est sensé, cependant je trouve qu'économiser est aussi très important. »
\end{exemplebox}

\begin{attention}[Erreur classique : \zh{我是同意} / \zh{我不是同意}]
Le francophone pense « je \emph{suis} d'accord » et ajoute le verbe \zh{是}. \textbf{Erreur !} En chinois, \zh{同意} est un verbe à part entière (verbe transitif) : on dit simplement \zh{我同意}, sans \zh{是}. De même : \zh{我不同意} (négation directe avec \zh{不}) et non *\zh{我不是同意}.
\end{attention}

\courssub{Relancer la parole et conclure}

Un débat vit grâce aux questions qui redonnent la parole aux autres. Trois formules à maîtriser absolument :

\begin{exemplebox}[Relancer un camarade]
\zh{你呢？} \emph{Nǐ ne?} « Et toi ? » (la plus courte, utilisée après avoir donné son propre avis)\\
\zh{你觉得怎么样？} \emph{Nǐ juéde zěnmeyàng?} « Qu'en penses-tu ? / Comment trouves-tu ça ? »\\
\zh{你有什么想法？} \emph{Nǐ yǒu shénme xiǎngfǎ?} « Quelles idées as-tu ? / Qu'en penses-tu ? » (plus ouvert)
\end{exemplebox}

Pour conclure, deux formules complémentaires : \zh{所以，我认为...} (\emph{suǒyǐ, wǒ rènwéi...}, « donc, je pense que... ») annonce la conclusion logique, et \zh{我的介绍完了} (\emph{wǒ de jièshào wán le}, « mon exposé est terminé », \zh{完了} = action achevée) ferme poliment la prise de parole, souvent suivie de \zh{谢谢大家} (« merci à tous »).

\begin{methode}[Structurer un mini-exposé en quatre étapes (armature universelle)]
\begin{enumerate}
\item \textbf{Ouvrir} avec \zh{先} (d'abord) : \zh{我先介绍一下...} — annoncer le sujet clairement.
\item \textbf{Développer} avec \zh{然后} (ensuite) : donner un ou deux arguments avec \zh{我觉得} ou \zh{对我来说}.
\item \textbf{Terminer le développement} avec \zh{最后} (enfin) : dernier argument, exemple personnel ou concession.
\item \textbf{Conclure} avec \zh{所以} (donc) : \zh{所以，我认为...} puis \zh{我的介绍完了，谢谢大家！}
\end{enumerate}
Cette armature \zh{先 → 然后 → 最后 → 所以} fonctionne pour tout exposé de deux minutes au baccalauréat.
\end{methode}

\courssub{Tableau récapitulatif des formules fonctionnelles}

\begin{center}
\begin{tabularx}{\linewidth}{|X|X|X|}
\hline
\rowcolor{popblueL} Fonction & Formule chinoise + pinyin & Équivalent français \\
\hline
Ouvrir & 我先介绍一下... \emph{wǒ xiān jièshào yíxià...} & Je vais d'abord présenter... \\
\hline
Donner son avis (oral) & 我觉得... \emph{wǒ juéde...} & Je trouve que... \\
\hline
Donner son avis (formel) & 我认为... \emph{wǒ rènwéi...} & Je pense que... / J'estime que... \\
\hline
Personnaliser (thème) & 对我来说，... \emph{duì wǒ láishuō,...} & Pour moi, ... / En ce qui me concerne... \\
\hline
Accord & 我同意。/ 你说得对。 \emph{wǒ tóngyì / nǐ shuō de duì} & Je suis d'accord. / Tu as raison. \\
\hline
Désaccord poli & 我不太同意... 可是/不过... \emph{wǒ bú tài tóngyì... kěshì/búguò...} & Je ne suis pas tout à fait d'accord... cependant... \\
\hline
Relancer & 你呢？/ 你觉得怎么样？/ 你有什么想法？ \emph{nǐ ne / nǐ juéde zěnmeyàng / nǐ yǒu shénme xiǎngfǎ} & Et toi ? / Qu'en penses-tu ? / Quelles idées as-tu ? \\
\hline
Conclure & 所以，我认为... / 我的介绍完了。 \emph{suǒyǐ, wǒ rènwéi... / wǒ de jièshào wán le} & Donc, je pense que... / Mon exposé est terminé. \\
\hline
\end{tabularx}
\end{center}

\begin{popfigure}
\begin{tikzpicture}[
  fct/.style={rectangle, rounded corners, draw=popblue, fill=popblueL, align=center, font=\small, minimum width=3.2cm, minimum height=0.9cm},
  frm/.style={rectangle, rounded corners, draw=poporange, fill=poporangeL, align=center, font=\small, minimum width=5.6cm, minimum height=0.9cm},
  arr/.style={-{Stealth[length=2mm]}, thick, popdark}
]
\node[fct] (f1) at (0,0) {Présenter};
\node[frm, align=center] (r1) at (6.2,0){我先介绍一下...\\ \emph{wǒ xiān jièshào yíxià...}};
\node[fct] (f2) at (0,-1.6) {Donner son avis};
\node[frm, align=center] (r2) at (6.2,-1.6){我觉得 / 我认为 / 对我来说\\ \emph{wǒ juéde / wǒ rènwéi / duì wǒ láishuō}};
\node[fct] (f3) at (0,-3.2) {Accord / désaccord};
\node[frm, align=center] (r3) at (6.2,-3.2){我同意 / 我不太同意，不过...\\ \emph{wǒ tóngyì / wǒ bú tài tóngyì, búguò...}};
\node[fct] (f4) at (0,-4.8) {Relancer};
\node[frm, align=center] (r4) at (6.2,-4.8){你呢？你觉得怎么样？\\ \emph{nǐ ne? nǐ juéde zěnmeyàng?}};
\node[fct] (f5) at (0,-6.4) {Conclure};
\node[frm, align=center] (r5) at (6.2,-6.4){所以，我认为... 我的介绍完了。\\ \emph{suǒyǐ, wǒ rènwéi... wǒ de jièshào wán le}};
\draw[arr] (f1) -- (r1);
\draw[arr] (f2) -- (r2);
\draw[arr] (f3) -- (r3);
\draw[arr] (f4) -- (r4);
\draw[arr] (f5) -- (r5);
\end{tikzpicture}
\legende{Carte mentale des formules fonctionnelles : chaque fonction de communication correspond à une ou plusieurs formules chinoises.}
\end{popfigure}

\courssub{Comparaison avec le français : points de transfert et pièges}

\begin{center}
\begin{tabularx}{\linewidth}{|X|X|X|}
\hline
\rowcolor{popgreenL} Situation & En français & En chinois \\
\hline
Donner un avis (adjectif) & « je pense que c'\emph{est} important » (verbe être) & 我觉得很重要。(\textbf{pas de 是} devant l'adjectif) \\
\hline
Accord & « je \emph{suis} d'accord » (être + adjectif) & 我同意。(\textbf{verbe seul}, sans 是) \\
\hline
Désaccord direct & « je ne suis pas d'accord » (franc, direct) & 我不太同意。(\textbf{adouci par 不太}) \\
\hline
Relance simple & « et toi ? » & 你呢？(\textbf{mot + particule 呢}) \\
\hline
\end{tabularx}
\end{center}

\begin{exoresolu}[Compléter un échange de débat]
Énoncé : Complétez ce mini-débat sur l'argent de poche avec les formules fonctionnelles appropriées (caractères + pinyin).\\
A : \_\_\_\_\_\_\_\_\_\_\_\_ (j'ouvre : « je vais d'abord présenter ») 学生的零花钱问题。\\
B : \_\_\_\_\_\_\_\_\_\_\_\_ (avis oral : « je trouve que ») 零花钱很有用。\\
A : \_\_\_\_\_\_\_\_\_\_\_\_ (désaccord poli : « je ne suis pas tout à fait d'accord »), \_\_\_\_\_\_\_\_\_\_\_\_ (concession : « cependant ») 学生应该先学习。\\
B : \_\_\_\_\_\_\_\_\_\_\_\_ (accord partiel : « tu as raison »), 可是节约也很重要。\\
A : \_\_\_\_\_\_\_\_\_\_\_\_ (conclusion : « donc, je pense que ») 每个人都应该有预算。
\tcblower
Solution détaillée :\\
A : \zh{我先介绍一下}学生的零花钱问题。 \emph{Wǒ xiān jièshào yíxià xuésheng de línghuāqián wèntí.} — Formule d'ouverture avec \zh{一下} pour adoucir.\\
B : \zh{我觉得}零花钱很有用。 \emph{Wǒ juéde línghuāqián hěn yǒuyòng.} — Avis oral, sans \zh{是} devant l'adjectif \zh{有用}.\\
A : \zh{我不太同意}，\zh{不过}学生应该先学习。 \emph{Wǒ bú tài tóngyì, búguò xuésheng yīnggāi xiān xuéxí.} — Désaccord poli + concession avec \zh{不过}.\\
B : \zh{你说得对}，可是节约也很重要。 \emph{Nǐ shuō de duì, kěshì jiéyuē yě hěn zhòngyào.} — Accord partiel (\zh{你说得对}) + nuance.\\
A : \zh{所以，我认为}每个人都应该有预算。 \emph{Suǒyǐ, wǒ rènwéi měi ge rén dōu yīnggāi yǒu yùsuàn.} — Conclusion logique avec \zh{所以} + \zh{我认为} (registre formel).
\end{exoresolu}

\coursec{Les tons à travailler}

\begin{propriete}[Rappels de phonétique indispensables pour l'oral]
La maîtrise des tons est notée à part entière au baccalauréat. Trois phénomènes sont critiques dans cette leçon :
\end{propriete}

\begin{definition}[Sandhi du 4e ton : \zh{不} (bù)]
\zh{不} se prononce en 2e ton (\textbf{bú}) devant un 4e ton.
\begin{itemize}
\item \zh{不太} \rightarrow \emph{bú tài} (pas \emph{bù tài})
\item \zh{不同意} \rightarrow \emph{bú tóngyì}
\item \zh{不好} \rightarrow \emph{bú hǎo}
\item \zh{不是} \rightarrow \emph{bú shì}
\end{itemize}
Devant les autres tons (1, 2, 3), \zh{不} reste 4e ton : \zh{不贵} (\emph{bù guì}), \zh{不忙} (\emph{bù máng}), \zh{不敢} (\emph{bù gǎn}).
\end{definition}

\begin{definition}[Sandhi du 1er ton : \zh{一} (yī)]
\zh{一} change de ton selon le ton du mot suivant :
\begin{itemize}
\item Devant 4e ton $\rightarrow$ 2e ton : \zh{一下} \rightarrow \emph{yíxià}, \zh{一定} \rightarrow \emph{yídìng}.
\item Devant 1, 2, 3e tons $\rightarrow$ 4e ton : \zh{一个} \rightarrow \emph{yí ge} (souvent neutre \emph{yí ge} à l'oral), \zh{一起} \rightarrow \emph{yìqǐ}, \zh{一年} \rightarrow \emph{yìnián}.
\item Isolé, à la fin de phrase, ou dans un nombre ordinal $\rightarrow$ 1er ton : \zh{第一} (\emph{dì yī}), \zh{一百} (\emph{yì bǎi} — nombre).
\end{itemize}
\end{definition}

\begin{definition}[Ton neutre (轻声) sur les particules grammaticales]
Les particules de structure et de modalité perdent leur ton propre et deviennent légères, courtes, sans contour de ton marqué :
\zh{一下} (\emph{yíxi\`a}), \zh{大家} (\emph{dàji\`a}), \zh{什么} (\emph{shénme}), \zh{呢} (\emph{ne}), \zh{了} (\emph{le}), \zh{的} (\emph{de}), \zh{得} (\emph{de}), \zh{地} (\emph{de}).
\end{definition}

\begin{definition}[Intonation de phrase : Déclaratif vs Interrogatif]
\begin{itemize}
\item \textbf{Phrase déclarative / Exposé :} Courbe globale descendante (high $\rightarrow$ low). Ex: \zh{我的介绍完了。} \emph{Wǒ de jièshào wán le.}
\item \textbf{Question \zh{吗} / \zh{呢} (oui/non) :} Montée finale nette sur la particule. Ex: \zh{你呢？} \emph{Nǐ ne?}
\item \textbf{Question \zh{什么} / \zh{怎么} / \zh{怎么样} (QQ) :} Courbe descendante sur le mot interrogatif, puis basse. Ex: \zh{你有什么想法？} \emph{Nǐ yǒu shénme xiǎngfǎ?}
\end{itemize}
\end{definition}

\begin{experience}[Entraînement ciblé : lecture à voix haute (5 min)]
En binôme, lisez le dialogue 2 (exposé modèle) en vous concentrant \textbf{uniquement} sur :
\begin{enumerate}
\item Les sandhi de \zh{不} (\zh{不太}, \zh{不同意}) et \zh{一} (\zh{一下}, \zh{一个}).
\item Les tons neutres (\zh{一下}, \zh{大家}, \zh{什么}, \zh{呢}, \zh{了}).
\item L'intonation finale descendante sur les affirmations (\zh{完了}, \zh{大家}) et montante sur \zh{你呢？}.
\end{enumerate}
Corrigez-vous mutuellement. Enregistrez-vous si possible.
\end{experience}

\coursec{Jeux de rôle et débat guidé}

\begin{methode}[Organisation d'un débat en classe (20 min)]
\begin{enumerate}
\item \textbf{Préparation (5 min) :} En binôme, tirez un rôle (A ou B) et un sujet. Notez 3 arguments avec les formules fonctionnelles.
\item \textbf{Échange (5 min) :} Jouez le dialogue. L'enseignant circule, note les tons et les formules utilisées.
\item \textbf{Rotation (5 min) :} Changez de partenaire, inversez les rôles (A devient B), même sujet.
\item \textbf{Synthèse (5 min) :} Un binôme volontaire passe au tableau. La classe identifie les formules entendues (\zh{先}, \zh{然后}, \zh{我觉得}, \zh{不太同意}, \zh{所以}...).
\end{enumerate}
\end{methode}

\begin{experience}[Jeu de rôle : Cartes « Pour / Contre »]
\textbf{Sujet :} « Faut-il donner de l'argent de poche aux lycéens ? » (\zh{高中生要不要零花钱？})

\textbf{Carte A (Pour) :} Tu penses que l'argent de poche apprend la gestion. Utilise : \zh{我先介绍一下}, \zh{我觉得}, \zh{对我来说}, \zh{你呢？}, \zh{所以我认为}.
\textbf{Carte B (Contre / Nuancé) :} Tu penses que c'est risqué sans éducation financière. Utilise : \zh{我不太同意}, \zh{不过}, \zh{你说得对}, \zh{你有什么想法？}, \zh{我的介绍完了}.

\textbf{Contraintes :} Minimum 6 répliques au total. Chaque élève doit utiliser \textbf{au moins 4 formules différentes} du tableau récapitulatif.
\end{experience}

\begin{experience}[Débat guidé : Motion formelle]
\textbf{Motion :} « La gestion du budget devrait être une matière obligatoire au lycée. » (\zh{理财课应该成为高中必修课。})

\textbf{Déroulement :}
\begin{enumerate}
\item Vote initial à main levée (Pour / Contre / Abstention).
\item Groupes de 4 : 2 « Pour », 2 « Contre ». Préparent 3 arguments chacun (5 min).
\item Débat en « ping-pong » : Un argument Pour $\rightarrow$ Un argument Contre $\rightarrow$ Réplique Pour...
\item Vote final. Comparaison.
\end{enumerate}
\textbf{Formules de débat avancé (bonus) :}
\zh{一方面... 另一方面...} (D'un côté... de l'autre...), \zh{举个例子} (Par exemple), \zh{据我所知} (Pour autant que je sache), \zh{总而言之} (En résumé).
\end{experience}

\coursec{Production orale guidée et évaluation}

\begin{methode}[Préparation de l'exposé noté (2 minutes) — Étape par étape]
\begin{enumerate}
\item \textbf{Choix du sujet (30 s) :} « Mon budget de lycéen » / « Comment je gère mon argent de poche » / « Le budget de ma famille à Yaoundé/Douala/Buea ».
\item \textbf{Planification (1 min) :} Remplissez le squelette :\\
\zh{1. 我先介绍一下... (Sujet)}\\
\zh{2. 然后，我觉得/对我来说... (Argument 1)}\\
\zh{3. 最后，... (Argument 2 / Exemple personnel)}\\
\zh{4. 所以，我认为... 我的介绍完了。 (Conclusion)}
\item \textbf{Répétition silencieuse (30 s) :} Vérifiez les tons (\zh{bú tài}, \zh{yíxià}, particules neutres) et l'enchaînement \zh{先 $\rightarrow$ 然后 $\rightarrow$ 最后 $\rightarrow$ 所以}.
\item \textbf{Passage à l'oral (2 min) :} Debout, regard vers la classe/jury, pas de lecture, notes autorisées (mots-clés seulement).
\end{enumerate}
\end{methode}

\begin{definition}[Grille d'auto-évaluation / évaluation par les pairs (sur 20)]
Utilisez cette grille lors des entraînements pour cibler vos progrès.
\end{definition}

\begin{center}
\begin{tabularx}{\linewidth}{|X|c|c|c|c|}
\hline
\rowcolor{poppurpleL} Critère & 5 pts (Excellent) & 3-4 pts (Acquis) & 1-2 pts (Fragile) & 0 pt (Non acquis) \\
\hline
\textbf{Prononciation / Tons} & Tons justes, sandhi (\zh{bù}/\zh{yī}) maîtrisés, intonation naturelle & Quelques erreurs de tons isolées, sandhi hésitant & Erreurs fréquentes sur tons de base, pas de sandhi & Incompréhensible, tons plats \\
\hline
\textbf{Vocabulaire budget} & 8+ mots du lexique utilisés correctement & 5-7 mots utilisés & $<$ 5 mots, répétitions, anglicismes & Vocabulaire hors sujet \\
\hline
\textbf{Structures / Grammaire} & Formules fonctionnelles variées (4+), ordre SOV correct, pas de \zh{是} superflu & 2-3 formules, petites erreurs d'ordre (compléments) & 1 formule, erreurs \zh{是} + adj, \zh{我是同意} & Phrases télégraphiques, pas de structure \\
\hline
\textbf{Cohérence / Enchaînement} & Armature \zh{先-然后-最后-所以} fluide, connecteurs logiques & Armature présente mais transitions raides & Sauts logiques, pas de connecteurs & Défilé de phrases sans lien \\
\hline
\textbf{Interaction (si débat)} & Relance naturelle (\zh{你呢？}), réagit aux arguments (accord/désaccord poli) & Relance présente, réaction basique & Attend son tour, ne réagit pas & Monologue, ignore l'autre \\
\hline
\end{tabularx}
\end{center}

\begin{exercice}[Production orale finale : « Mon budget idéal »]
Préparez un exposé de \textbf{2 minutes maximum} sur le thème : \zh{我的理想预算} (Mon budget idéal).
\textbf{Contraintes obligatoires :}
\begin{enumerate}
\item Utiliser \textbf{au moins 6 mots} du tableau de vocabulaire (Section 1).
\item Utiliser \textbf{l'armature complète} : \zh{先} — \zh{然后} — \zh{最后} — \zh{所以}.
\item Intégrer \textbf{au moins 3 formules fonctionnelles} différentes (ouvrir, avis, conclure).
\item Respecter les \textbf{sandhi} (\zh{不}, \zh{一}) et les \textbf{tons neutres}.
\item Contexte camerounais bienvenu : mentionner \zh{喀麦隆}, \zh{雅温得}, \zh{杜阿拉}, \zh{非洲法郎} (CFA), \zh{学费}, \zh{吃饭}...
\end{enumerate}
\textit{Conseil :} Rédigez d'abord le squelette en caractères/pinyin sur une fiche bristol (mots-clés seulement), puis entraînez-vous 3 fois à voix haute avant de passer.
\end{exercice}

\begin{corrige}[Exercice : Modèle de réponse type (exemple pour l'enseignant)]
\zh{大家好！我先介绍一下我的理想预算。我是喀麦隆的高中生，在雅温得上学。然后，我觉得收入很重要。我的收入是父母给的生活费，每个月五万非洲法郎。对我来说，记账非常有用，我每天用手机记账。最后，所以我认为存钱比花钱重要。我想每个月存一万法郎，买书或者备用。我的介绍完了，谢谢大家！}\\
\emph{Dàjiā hǎo! Wǒ xiān jièshào yíxià wǒ de lǐxiǎng yùsuàn. Wǒ shì Kāmàilóng de gāozhōngshēng, zài Yǎwēndé shàngxué. Ránhòu, wǒ juéde shōurù hěn zhòngyào. Wǒ de shōurù shì fùmǔ gěi de shēnghuófèi, měi ge yuè wǔ wàn Fēizhōu fǎláng. Duì wǒ láishuō, jìzhàng fēicháng yǒuyòng, wǒ měi tiān yòng shǒujī jìzhàng. Zuìhòu, suǒyǐ wǒ rènwéi cúnqián bǐ huā qián zhòngyào. Wǒ xiǎng měi ge yuè cún yí wàn fǎláng, mǎi shū huòzhě bèiyòng. Wǒ de jièshào wán le, xièxie dàjiā!}
\end{corrige}

\begin{savaistu}[Les connecteurs, vos meilleurs amis à l'oral du baccalauréat]
À l'épreuve orale, l'examinateur n'écoute pas seulement vos tons et votre vocabulaire : il évalue votre capacité à \textbf{organiser un discours}. Les connecteurs \zh{先} (d'abord), \zh{然后} (ensuite), \zh{但是} (mais) et \zh{所以} (donc) sont des marqueurs d'organisation très valorisés. Un exposé simple mais bien articulé avec ces quatre mots obtient toujours une meilleure note qu'un discours riche mais désordonné. Au Cameroun comme en Chine, un bon orateur est d'abord un orateur clair !
\end{savaistu}

\begin{aretenir}[L'essentiel de la leçon 28 — Gestion du budget (Expression orale)]
\begin{itemize}
\item \textbf{Lexique clé :} \zh{预算, 收入, 支出, 工资, 生活费, 零花钱, 节约, 浪费, 存钱, 管理, 负责, 记账}.
\item \textbf{Formules fonctionnelles :} Ouvrir (\zh{我先介绍一下}), Avis (\zh{我觉得}/\zh{我认为}/\zh{对我来说}), Accord (\zh{我同意}), Désaccord poli (\zh{我不太同意...不过}), Relance (\zh{你呢？}/\zh{你有什么想法？}), Conclure (\zh{所以我认为... 我的介绍完了}).
\item \textbf{Grammaire piège :} Pas de \zh{是} devant adjectif après \zh{觉得}/\zh{认为} ; pas de \zh{是} devant \zh{同意}.
\item \textbf{Phonétique :} Sandhi \zh{不} (bú devant 4e ton), \zh{一} (yí/yì selon contexte), tons neutres (particules), intonation finale (descendante affirmation, montante \zh{ne/ma}).
\item \textbf{Structure d'exposé :} \zh{先 $\rightarrow$ 然后 $\rightarrow$ 最后 $\rightarrow$ 所以}.
\item \textbf{Culture :} Politesse du désaccord (\zh{不太}), harmonie du groupe, importance de l'épargne en Chine (\zh{储蓄率} élevée) comme au Cameroun (tontines, \zh{会}).
\end{itemize}
\end{aretenir}

\coursec{Les tons à travailler}

Dans la section précédente, nous avons mémorisé les formules fonctionnelles du débat : \zh{我觉得…} (Wǒ juéde…, « je pense que… »), \zh{我同意} (Wǒ tóngyì, « je suis d'accord »), \zh{你说得对} (Nǐ shuō de duì, « tu as raison »). Mais une formule parfaitement mémorisée ne sert à rien si elle est prononcée avec de mauvais tons : votre interlocuteur chinois ne comprendra tout simplement pas. Cette section est donc consacrée à l'entraînement phonétique intensif sur les mots clés de notre thème : la gestion du budget. L'objectif est triple : réviser les quatre tons et le ton neutre, maîtriser le sandhi du troisième ton, et placer correctement l'intonation des questions.

\courssub{Révision des quatre tons et du ton neutre}

Commençons par observer quatre mots essentiels de notre leçon. Prononcez-les à voix haute en suivant la courbe mélodique indiquée :

\begin{exemplebox}[Quatre mots, quatre tons]
\zh{花钱} (huā qián, « dépenser de l'argent ») : \cle{huā} est au premier ton, haut et plat ; \cle{qián} est au deuxième ton, montant.\\
\zh{省钱} (shěng qián, « économiser de l'argent ») : \cle{shěng} est au troisième ton, qui descend puis remonte.\\
\zh{太贵了！} (Tài guì le !, « C'est trop cher ! ») : \cle{guì} est au quatrième ton, qui tombe brusquement.
\end{exemplebox}

\begin{popfigure}
\begin{tikzpicture}[scale=1.0]
% Grille de fond (niveaux 1 à 5)
\draw[popline, dashed] (0,1) -- (11,1);
\draw[popline, dashed] (0,2) -- (11,2);
\draw[popline, dashed] (0,3) -- (11,3);
\draw[popline, dashed] (0,4) -- (11,4);
% Ton 1 : haut et plat (niveau 5)
\draw[popblue, very thick] (0.5,4) -- (2.5,4);
\node[popblue, align=center] at (1.5,4.6) {\textbf{1er ton}\\ 花 huā};
\node[popblue, below] at (1.5,3.8) {\footnotesize haut et plat (55)};
% Ton 2 : montant (niveau 3 à 5)
\draw[popgreen, very thick] (3.5,2) -- (5.5,4);
\node[popgreen, align=center] at (4.5,4.6) {\textbf{2e ton}\\ 钱 qián};
\node[popgreen, below] at (4.5,1.8) {\footnotesize montant (35)};
% Ton 3 : bas, descente-remontée (niveau 2-1-4)
\draw[poporange, very thick] (6.5,3) .. controls (7.1,1) and (7.9,1) .. (8.5,3.5);
\node[poporange, align=center] at (7.5,4.6) {\textbf{3e ton}\\ 省 shěng};
\node[poporange, below] at (7.5,0.8) {\footnotesize descend-remonte (214)};
% Ton 4 : chute brutale (niveau 5 à 1)
\draw[poppurple, very thick] (9.5,4) -- (11.5,1);
\node[poppurple, align=center] at (10.5,4.6) {\textbf{4e ton}\\ 贵 guì};
\node[poppurple, below] at (10.5,0.8) {\footnotesize tombe brusquement (51)};
\end{tikzpicture}
\legende{Les quatre courbes tonales appliquées aux mots clés du budget : huā (1er), qián (2e), shěng (3e), guì (4e). Niveaux de 1 (bas) à 5 (haut).}
\end{popfigure}

\begin{aretenir}[Les cinq hauteurs de la syllabe chinoise]
\begin{itemize}
\item \textbf{1er ton} (ā) : haut, plat, constant (55) — comme un médecin qui vous demande de dire « aaa » longuement.
\item \textbf{2e ton} (á) : montant (35) — comme le « hein ? » interrogatif et surpris du français.
\item \textbf{3e ton} (ǎ) : bas, avec une descente puis une remontée (214) — à l'oral rapide, on n'entend souvent que la partie basse (demi-3e ton).
\item \textbf{4e ton} (à) : descendant et sec (51) — comme un « non ! » catégorique et bref.
\item \textbf{Ton neutre} (a) : court, léger, sans accent fixe, sa hauteur dépend du ton précédent — comme dans \zh{觉得} juéde ou \zh{我们} wǒmen.
\end{itemize}
\end{aretenir}

\begin{center}
\begin{tabularx}{\linewidth}{|X|X|X|X|}
\hline
\rowcolor{popblueL} Caractères & Pinyin & Ton & Traduction \\
\hline
花 & huā & 1er & dépenser ; fleur \\
\hline
钱 & qián & 2e & argent \\
\hline
省 & shěng & 3e & économiser \\
\hline
贵 & guì & 4e & cher \\
\hline
买 & mǎi & 3e & acheter \\
\hline
卖 & mài & 4e & vendre \\
\hline
零花钱 & línghuāqián & 2e-1er-2e & argent de poche \\
\hline
觉得 & juéde & 2e + neutre & penser, trouver \\
\hline
\end{tabularx}
\end{center}

\courssub{Le sandhi du troisième ton (变调 biàndiào)}

Observez attentivement ces deux phrases tirées de nos dialogues :

\begin{exemplebox}[Observation]
\zh{我想省钱。} — \emph{Wǒ xiǎng shěng qián.} — « Je veux économiser de l'argent. »\\
\zh{我很好。} — \emph{Wǒ hěn hǎo.} — « Je vais très bien. »
\end{exemplebox}

À l'écrit, \zh{我} est wǒ (3e ton) et \zh{想} est xiǎng (3e ton). Mais à l'oral, aucun Chinois ne prononce deux courbes complètes « descente-remontée » à la suite : ce serait trop lent et trop lourd. La langue applique automatiquement une transformation.

\begin{propriete}[Sandhi du troisième ton]
Quand \textbf{deux syllabes au troisième ton se suivent}, la première se prononce comme un \textbf{deuxième ton} (montant). L'écriture pinyin ne change pas, seule la prononciation change.\\
Schéma : 3e + 3e $\to$ \textbf{2e} + 3e\\
\zh{我想} wǒ xiǎng se prononce \emph{wó xiǎng}.\\
\zh{很好} hěn hǎo se prononce \emph{hén hǎo}.\\
\zh{你好} nǐ hǎo se prononce \emph{ní hǎo}.\\
\zh{我省钱} wǒ shěng qián se prononce \emph{wó shěng qián}.
\end{propriete}

\begin{attention}[Le pinyin écrit ne change jamais]
On écrit toujours \zh{我想} wǒ xiǎng avec deux marques de 3e ton, même si on prononce wó xiǎng. Dans un examen écrit, ne « corrigez » jamais le pinyin : la règle du sandhi est une règle de \cle{prononciation}, pas d'orthographe. De même, si trois 3e tons se suivent, on applique la règle par groupes de sens (mots) : \zh{我也想买} (wǒ yě xiǎng mǎi) se groupe souvent en [wǒ yě] [xiǎng mǎi] $\to$ \emph{wó yé xiáng mǎi} (2e+2e+2e+3e).
\end{attention}

\courssub{Suites difficiles à travailler}

Deux enchaînements de nos dialogues demandent un entraînement particulier.

\textbf{1. 我觉得 (wǒ juéde) : 3e + 2e + neutre.} Ici, pas de sandhi (le mot suivant n'est pas au 3e ton). La difficulté est ailleurs : le 3e ton de \zh{我} devant un autre ton (sauf 3e) se réduit à sa \textbf{moitié basse} (on dit « demi-3e ton ») : on descend légèrement sans remonter. Puis \zh{觉} monte (2e ton) et \zh{得} est léger et court (ton neutre). Répétez : \emph{wǒ( bas) jué(montant) de(léger)}.

\textbf{2. 零花钱 (línghuāqián) : 2e + 1er + 2e.} Trois tons pleins à enchaîner sans pause : montant, haut-plat, montant. Découpez d'abord en deux : \zh{零花} (línghuā) puis \zh{钱} (qián), puis soudez : \emph{línghuāqián}.

\begin{exemplebox}[Phrases d'entraînement]
\zh{我觉得零花钱太少了。} — \emph{Wǒ juéde línghuāqián tài shǎo le.} — « Je trouve que l'argent de poche est trop peu. »\\
\zh{你的零花钱够吗？} — \emph{Nǐ de línghuāqián gòu ma ?} — « Ton argent de poche suffit-il ? »\\
\zh{我觉得这本书不贵。} — \emph{Wǒ juéde zhè běn shū bù guì.} — « Je trouve que ce livre n'est pas cher. »
\end{exemplebox}

\courssub{L'intonation des questions : 吗 et 呢}

Au-delà des tons de chaque syllabe, la phrase entière possède une mélodie globale, comme en français. Deux cas sont essentiels pour notre débat.

\begin{propriete}[Intonation des questions]
\begin{itemize}
\item Question avec \cle{吗} (ma) : l'intonation est \textbf{montante} à la fin, comme le « est-ce que… ? » du français. La particule ma elle-même est au ton neutre, légère, mais portée par la montée de la phrase. Exemple : \zh{贵吗？} \emph{Guì ma ?} (« C'est cher ? ») — la voix monte sur ma.
\item Question avec \cle{呢} (ne) : l'intonation est \textbf{retombante} à la fin. ne sert à relancer (« et toi ? ») ou à demander où est quelque chose. Exemple : \zh{我呢？} \emph{Wǒ ne ?} (« Et moi ? ») — la voix redescend sur ne.
\end{itemize}
\end{propriete}

\begin{exemplebox}[Comparaison à voix haute]
\zh{你想买吗？} — \emph{Nǐ xiǎng mǎi ma ?} (intonation montante) — « Tu veux acheter ? »\\
\zh{你想买什么呢？} — \emph{Nǐ xiǎng mǎi shénme ne ?} (intonation retombante) — « Qu'est-ce que tu veux acheter ? »\\
\zh{我的钱够，你的呢？} — \emph{Wǒ de qián gòu, nǐ de ne ?} — « Mon argent suffit, et le tien ? »
\end{exemplebox}

\courssub{Paires minimales : mǎi / mài et shěng / shēng}

Une \cle{paire minimale} est un couple de mots qui ne diffèrent que par un seul trait phonétique — ici, uniquement le ton. C'est le test le plus redoutable pour un francophone.

\begin{attention}[mǎi ou mài : le sens s'inverse complètement !]
Confondre \zh{买} mǎi (acheter, 3e ton) et \zh{卖} mài (vendre, 4e ton) inverse totalement le sens de la phrase ! Si vous dites \zh{我想卖一本书} en voulant dire « je veux acheter un livre », votre interlocuteur comprendra que vous voulez \textbf{vendre} un livre. Au marché de Mokolo ou de Douala, l'erreur peut coûter cher ! Moyen mnémotechnique : \zh{买} mǎi (3e ton, la voix « creuse ») = on \emph{creuse} dans sa poche pour payer ; \zh{卖} mài (4e ton, la voix tombe) = l'argent \emph{tombe} dans la caisse. Remarquez aussi l'écriture : \zh{卖} possède un trait supplémentaire en haut (\zh{十}), ce que l'on « ajoute » quand on vend.
\end{attention}

\begin{exemplebox}[Les deux phrases en regard]
\zh{我想买一本书。} — \emph{Wǒ xiǎng mǎi yì běn shū.} — « Je veux acheter un livre. »\\
\zh{我想卖我的旧书。} — \emph{Wǒ xiǎng mài wǒ de jiù shū.} — « Je veux vendre mes vieux livres. »
\end{exemplebox}

Deuxième paire : \zh{省} shěng (économiser, 3e ton) et \zh{生} shēng (naître, 1er ton). Comparer : \zh{我省钱。} \emph{Wǒ shěng qián.} (« J'économise de l'argent. ») et \zh{我生在雅温得。} \emph{Wǒ shēng zài Yǎwēndé.} (« Je suis né à Yaoundé. »).

\begin{methode}[Vérifier ses tons en trois étapes]
\begin{enumerate}
\item \textbf{Lent et exagéré} : dites le mot seul, très lentement, en dessinant la courbe du ton avec la main dans l'air (plat, montée, creux, chute).
\item \textbf{Vitesse normale} : répétez le mot cinq fois à vitesse de conversation, sans la main, en gardant la même mélodie.
\item \textbf{Dans la phrase} : replacez le mot dans la phrase complète du dialogue et vérifiez que le ton survit au débit rapide. Si le ton se perd, revenez à l'étape 1.
\end{enumerate}
\end{methode}

\begin{exoresolu}[Identifier et corriger les tons]
Énoncé : un élève prononce les phrases suivantes avec les tons indiqués. Trouvez l'erreur de ton dans chaque phrase et donnez la prononciation correcte.\\
1. \zh{我买书} prononcé \emph{wǒ mài shū} pour dire « j'achète un livre ».\\
2. \zh{太贵了} prononcé \emph{tài guí le}.\\
3. \zh{我想省钱} prononcé \emph{wǒ(3e plein) xiǎng(3e plein) shěng qián}, avec deux 3e tons complets.
\tcblower
Solution détaillée :\\
1. L'élève a prononcé mài (4e ton, vendre) au lieu de \textbf{mǎi} (3e ton, acheter). Correction : \emph{wǒ mǎi shū}. Le sens était inversé !\\
2. \zh{贵} est au 4e ton (guì), qui tombe, et non au 2e ton montant (guí). Correction : \emph{tài guì le}.\\
3. Deux 3e tons consécutifs : par le sandhi, \zh{我想} se prononce \emph{wó xiǎng} (2e + 3e). La suite correcte est donc \emph{wó xiǎng shěng qián}. L'écriture reste wǒ xiǎng.
\end{exoresolu}

\courssub{Travail rythmique : frapper les tons}

Pour ancrer les tons dans le corps et non seulement dans la tête, nous utilisons un travail rythmique collectif : on \cle{scande} les phrases du dialogue en frappant dans les mains ou sur la table, un coup par syllabe, en exagérant la hauteur de la voix.

\begin{experience}[Scander le dialogue en frappant les tons]
En groupes de quatre, prenez la phrase : \zh{我想省钱，不想乱花钱。} (\emph{Wǒ xiǎng shěng qián, bù xiǎng luàn huā qián.} — « Je veux économiser, je ne veux pas dépenser n'importe comment. »)
\begin{enumerate}
\item Découpez la phrase en syllabes et marquez le ton de chacune au tableau : wǒ(3) xiǎng(3) shěng(3) qián(2) — bù(4) xiǎng(3) luàn(4) huā(1) qián(2).
\item Appliquez les transformations : \zh{我想} $\to$ wó xiǎng (sandhi 3e+3e) ; \zh{不} bù (4e ton) reste 4e ton car suivi de \zh{想} (3e ton) — \textbf{rappel : bù devient 2e ton seulement devant un 4e ton}.
\item Scandez en frappant : un coup fort sur chaque syllabe, la voix suivant la courbe. Recommencez trois fois en accélérant.
\item Un élève récite seul ; les autres lèvent la main dès qu'un ton est faux.
\end{enumerate}
\end{experience}

\begin{exercice}[À vous de jouer]
1. Marquez les tons de : \zh{我觉得这本书太贵了。} puis prononcez-la trois fois en appliquant le sandhi si nécessaire.\\
2. Prononcez les paires minimales à voix haute : mǎi/mài, shěng/shēng, puis faites deviner à votre voisin lequel des deux mots vous avez dit.\\
3. Posez la question \zh{你的零花钱够吗？} avec l'intonation montante, puis relancez avec \zh{你的呢？} avec l'intonation retombante.
\end{exercice}

\begin{corrige}[Exercice]
1. Wǒ(3/demi-3e) jué(2) de(neutre) zhè(4) běn(3) shū(1) tài(4) guì(4) le(neutre). Pas de sandhi 3e+3e ici (\zh{我} est suivi d'un 2e ton : demi-3e ton bas).\\
2. mǎi = 3e ton (creux), mài = 4e ton (chute) ; shěng = 3e ton, shēng = 1er ton (haut et plat).\\
3. Sur \zh{吗} la voix monte ; sur \zh{呢} la voix redescend doucement.
\end{corrige}

\begin{aretenir}[L'essentiel phonétique de la leçon]
\begin{itemize}
\item Quatre tons + ton neutre : huā (1er), qián (2e), shěng (3e), guì (4e).
\item Sandhi : 3e + 3e $\to$ 2e + 3e (wǒ xiǎng $\to$ wó xiǎng), prononciation seule, jamais l'écriture.
\item Demi-3e ton : 3e ton suivi d'un ton 1, 2, 4 ou neutre $\to$ seule la partie basse.
\item \zh{不} bù : 4e ton $\to$ 2e ton (bú) seulement devant un 4e ton.
\item mǎi (acheter) / mài (vendre) : un ton différent, un sens opposé.
\item Question avec 吗 : intonation montante ; avec 呢 : intonation retombante.
\item Méthode : lent et exagéré, puis vitesse normale, puis dans la phrase.
\end{itemize}
\end{aretenir}

\coursec{Vocabulaire du budget : revenus, dépenses et gestion}

Avant de savoir en parler, il faut connaître les mots. Voici le lexique de base pour gérer un budget, classer ses dépenses et exprimer ses choix financiers. Ce vocabulaire est indispensable pour les jeux de rôle et le débat qui suivront.

\begin{center}
\begin{tabularx}{\linewidth}{|X|X|X|X|}
\hline
\rowcolor{popblueL} Caractères & Pinyin & Nature & Traduction \\
\hline
零花钱 & línghuāqián & n. & argent de poche \\
\hline
工资 & gōngzī & n. & salaire \\
\hline
红包 & hóngbāo & n. & enveloppe rouge (cadeau d'argent) \\
\hline
收入 & shōurù & n. & revenu, rentrée d'argent \\
\hline
支出 & zhīchū & n. & dépense, sortie d'argent \\
\hline
预算 & yùsuàn & n. & budget, prévision budgétaire \\
\hline
分配 & fēnpèi & v. & répartir, distribuer \\
\hline
吃饭 & chīfàn & v. & manger ; poste « nourriture » \\
\hline
交通 & jiāotōng & n. & transport \\
\hline
娱乐 & yúlè & n. & loisirs, divertissement \\
\hline
学习用品 & xuéxí yòngpǐn & n. & fournitures scolaires \\
\hline
零食 & língshí & n. & snacks, grignotines \\
\hline
花 & huā & v. & dépenser (argent, temps) \\
\hline
省 & shěng & v. & économiser, épargner \\
\hline
省着花 & shěngzhe huā & expr. & dépenser avec économie, faire attention \\
\hline
记账 & jìzhàng & v. & tenir ses comptes, noter ses dépenses \\
\hline
存钱 & cún qián & v. & mettre de l'argent de côté, épargner \\
\hline
紧张 & jǐnzhāng & adj. & serré (budget), tendu \\
\hline
富裕 & fùyù & adj. & à l'aise (financièrement), confortable \\
\hline
乱花 & luàn huā & v. & dépenser n'importe comment, gaspiller \\
\hline
\end{tabularx}
\end{center}

\begin{propriete}[Classificateurs courants pour l'argent et les objets]
En chinois, tout nom comptable nécessite un classificateur entre le nombre/démonstratif et le nom.
\begin{center}
\begin{tabularx}{\linewidth}{|X|X|X|}
\hline
\rowcolor{popblueL} Classificateur & Pinyin & Emploi principal \\
\hline
元 & yuán & Unité principale du yuan (Chine) ; aussi \zh{块} (kuài) à l'oral \\
\hline
块 & kuài & Yuan (oral), équivalent de « balles » ou « euros » familier \\
\hline
角 & jiǎo & Dixième de yuan (10 fen = 1 jiao) \\
\hline
分 & fēn & Centième de yuan \\
\hline
笔 & bǐ & Une dépense, une entrée de compte (\zh{一笔钱}) \\
\hline
份 & fèn & Une part, une portion (\zh{一份工资}) \\
\hline
张 & zhāng & Billet, ticket, carte (\zh{一张纸}, \zh{一张卡}) \\
\hline
个 & gè & Classificateur générique (\zh{一个红包}) \\
\hline
\end{tabularx}
\end{center}
\textbf{Exemples :} \zh{十块钱} (shí kuài qián, « 10 yuan / 10 balles ») ; \zh{三笔支出} (sān bǐ zhīchū, « trois dépenses ») ; \zh{一份预算} (yī fèn yùsuàn, « un budget »).

\end{propriete}

\begin{attention}[Faux amis et pièges lexicaux]
\begin{itemize}
\item \zh{省} (shěng) = « économiser » (verbe) OU « province » (nom). Ne pas confondre avec \zh{性} (xìng, nature/sexe) ni \zh{胜} (shèng, victoire).
\item \zh{花} (huā) = « dépenser » (argent) OU « fleur » (nom). \zh{花钱} (huā qián) = dépenser de l'argent.
\item \zh{零花钱} (línghuāqián) : \zh{零} (líng) =零碎, \zh{花} (huā) =dépenser, \zh{钱} (qián) =argent. Litt. « argent pour menues dépenses ».
\item \zh{记账} (jìzhàng) : \zh{记} (jì) = noter, \zh{账} (zhàng) = compte. Ne pas écrire \zh{帐} (zhàng, tente/toile) qui est un caractère proche mais erroné ici.
\end{itemize}
\end{attention}

\coursec{Dialogues modèles : situations de la vie courante}

Deux dialogues types pour entendre le vocabulaire en contexte et travailler les tons. Écoutez (ou lisez à haute voix), observez les structures, puis répétez.

\courssub{Dialogue 1 : « Fin de mois difficile » (deux amis)}

\begin{textebox}[Dialogue 1 — Fin de mois]
\zh{A：哎，这个月怎么过呀？}\par
\emph{Āi, zhège yuè zěnme guò ya ?} — « Dis, comment on fait pour ce mois-ci ? »\par
\zh{B：别提了，零花钱早就花光了。}\par
\emph{Bié tí le, línghuāqián zǎo jiù huā guāng le.} — « N'en parle pas, l'argent de poche est parti depuis longtemps. »\par
\zh{A：你买什么了？}\par
\emph{Nǐ mǎi shénme le ?} — « Tu as acheté quoi ? »\par
\zh{B：游戏皮肤、奶茶、还请同学吃饭…… 乱花了一通。}\par
\emph{Yóuxì pífū, nǎichá, hái qǐng tóngxué chīfàn… luàn huā le yì tōng.} — « Skins de jeux, bubble tea, j'ai payé resto aux copains… j'ai dilapidé n'importe comment. »\par
\zh{A：下个月你得记账，省着花。应该先买学习用品，不应该天天点外卖。}\par
\emph{Xià ge yuè nǐ děi jìzhàng, shěngzhe huā. Yīnggāi xiān mǎi xuéxí yòngpǐn, bù yīnggāi tiāntiān diǎn wàimài.} — « Le mois prochain, tu tiens tes comptes et tu fais attention. Tu dois d'abord acheter les fournitures, pas commander à manger tous les jours. »\par
\zh{B：说得对。我一定改。}\par
\emph{Shuō de duì. Wǒ yīdìng gǎi.} — « T'as raison. Je change, promis. »
\end{textebox}

\courssub{Dialogue 2 : « Conseil de famille » (quatre personnes)}

\begin{textebox}[Dialogue 2 — Conseil de famille]
\zh{爸爸：这个月预算十万法郎。我们要分配一下。}\par
\emph{Bàba : Zhège yuè yùsuàn shí wàn fǎláng. Wǒmen yào fēnpèi yí xià.} — « Papa : Budget du mois 100 000 francs. Il faut répartir. »\par
\zh{妈妈：吃饭最重要，给四万。}\par
\emph{Māma : Chīfàn zuì zhòngyào, gěi sì wàn.} — « Maman : Manger c'est prioritaire, on met 40 000. »\par
\zh{儿子：交通也要两万，我每天坐车上学。}\par
\emph{Érzi : Jiāotōng yě yào liǎng wàn, wǒ měitiān zuò chē shàngxué.} — « Fils : Transport aussi 20 000, je prends le bus tous les jours. »\par
\zh{女儿：娱乐一万够了，学习用品一万。剩下两万存着。}\par
\emph{Nǚ'ér : Yúlè yí wàn gòu le, xuéxí yòngpǐn yí wàn. Shèng xià liǎng wàn cún zhe.} — « Fille : Loisirs 10 000 suffisent, fournitures 10 000. Le reste 20 000 on épargne. »\par
\zh{爸爸：好，就这样定了。记得记账！}\par
\emph{Bàba : Hǎo, jiù zhème dìng le. Jìde jìzhàng !} — « Papa : Bon, c'est décidé. N'oubliez pas de noter les dépenses ! »
\end{textebox}

\begin{methode}[Travailler un dialogue en trois temps]
\begin{enumerate}
\item \textbf{Compréhension globale :} Qui parle ? De quoi ? Quel est le ton (familier, sérieux) ?
\item \textbf{Analyse linguistique :} Repérez les structures cibles (\zh{应该}, \zh{因为...所以...}, \zh{得}, \zh{了}), le vocabulaire nouveau, les classificateurs.
\item \textbf{Appropriation orale :} Lisez à deux, inversez les rôles, changez les détails (montants, postes de dépenses), jouez sans notes.
\end{enumerate}
\end{methode}

\coursec{Formules fonctionnelles pour débattre et conseiller}

Ces « briques » linguistiques sont les outils de votre expression orale. Mémorisez-les par cœur ; elles doivent sortir automatiquement.

\courssub{Exprimer son avis / son opinion}

\begin{center}
\begin{tabularx}{\linewidth}{|X|X|X|}
\hline
\rowcolor{popblueL} Structure & Exemple (caractères + pinyin) & Traduction \\
\hline
我觉得 + phrase & \zh{我觉得应该给零花钱。} \emph{Wǒ juéde yīnggāi gěi línghuāqián.} & Je pense qu'il faut donner de l'argent de poche. \\
\hline
我认为 + phrase & \zh{我认为记账很重要。} \emph{Wǒ rènwéi jìzhàng hěn zhòngyào.} & Je considère que tenir ses comptes est important. \\
\hline
依我看来 + phrase & \zh{依我看来，省钱是好习惯。} \emph{Yī wǒ kànlái, shěng qián shì hǎo xíguàn.} & À mon avis, économiser est une bonne habitude. \\
\hline
\end{tabularx}
\end{center}

\courssub{Donner un conseil / une recommandation}

\begin{center}
\begin{tabularx}{\linewidth}{|X|X|X|}
\hline
\rowcolor{popblueL} Structure & Exemple & Traduction \\
\hline
你应该 + V & \zh{你应该先存钱。} \emph{Nǐ yīnggāi xiān cún qián.} & Tu devrais d'abord épargner. \\
\hline
你不应该 + V & \zh{你不应该乱花钱。} \emph{Nǐ bù yīnggāi luàn huā qián.} & Tu ne devrais pas gaspiller. \\
\hline
建议你 + V & \zh{建议你每周记一次账。} \emph{Jiànyì nǐ měi zhōu jì yí cì zhàng.} & Je te conseille de tenir tes comptes chaque semaine. \\
\hline
最好 + V & \zh{最好带现金，别只用手机。} \emph{Zuìhǎo dài xiànjīn, bié zhǐ yòng shǒujī.} & Mieux vaut emporter du cash, pas seulement le téléphone. \\
\hline
\end{tabularx}
\end{center}

\courssub{Être d'accord / Pas d'accord / Nuancer}

\begin{center}
\begin{tabularx}{\linewidth}{|X|X|X|}
\hline
\rowcolor{popblueL} Structure & Exemple & Traduction \\
\hline
我同意 / 你说得对 & \zh{我同意你的看法。} \emph{Wǒ tóngyì nǐ de kànfǎ.} & Je suis d'accord / Tu as raison. \\
\hline
我不太同意 & \zh{我不太同意，因为...} \emph{Wǒ bù tài tóngyì, yīnwèi...} & Je ne suis pas tout à fait d'accord, car... \\
\hline
有道理，可是... & \zh{你的话有道理，可是...} \emph{Nǐ de huà yǒu dàoli, kěshì...} & Ton argument tient la route, mais... \\
\hline
不一定吧 & \zh{不一定吧，看情况。} \emph{Bù yídìng ba, kàn qíngkuàng.} & Pas sûr, ça dépend. \\
\hline
\end{tabularx}
\end{center}

\courssub{Relancer / Passer la parole / Conclure}

\begin{center}
\begin{tabularx}{\linewidth}{|X|X|X|}
\hline
\rowcolor{popblueL} Structure & Exemple & Traduction \\
\hline
你呢？ / 你怎么看？ & \zh{你呢？你怎么看？} \emph{Nǐ ne? Nǐ zěnme kàn?} & Et toi ? Qu'en penses-tu ? \\
\hline
还有别的想法吗？ & \zh{还有别的想法吗？} \emph{Hái yǒu biéde xiǎngfǎ ma?} & D'autres idées ? \\
\hline
所以，我们认为... & \zh{所以，我们认为应该给零花钱。} \emph{Suǒyǐ, wǒmen rènwéi yīnggāi gěi línghuāqián.} & Donc, nous pensons qu'il faut donner de l'argent de poche. \\
\hline
总之... & \zh{总之，记账是必须的。} \emph{Zǒngzhī, jìzhàng shì bìxū de.} & En résumé, tenir ses comptes est indispensable. \\
\hline
\end{tabularx}
\end{center}

\coursec{Les tons à travailler : points de vigilance}

La prononciation juste est la condition de la compréhension. Cette leçon met l'accent sur quatre difficultés tonales fréquentes chez les francophones.

\begin{propriete}[Rappel : les quatre tons et le ton neutre]
\begin{itemize}
\item \textbf{1er ton (ˉ) :} haut et plat. \zh{花} (huā), \zh{分} (fēn), \zh{预} (yù $\to$ yú en composition ? Non, yù 4e ton. \zh{预算} yùsuàn).
\item \textbf{2e ton (ˊ) :} montant. \zh{钱} (qián), \zh{零} (líng), \zh{来} (lái).
\item \textbf{3e ton (ˇ) :} bas-descendant-remontant (souvent bas-descendant seul). \zh{省} (shěng), \zh{记} (jì), \zh{给} (gěi), \zh{很} (hěn).
\item \textbf{4e ton (ˋ) :} descendant fort. \zh{应该} (yīnggāi - 1-1), \zh{不} (bù $\to$ bú devant 4e ton), \zh{花} (huà - verbe « transformer » rare, ici huā 1er ton), \zh{分配} (fēnpèi - 1-4).
\item \textbf{Ton neutre :} léger, court. \zh{了} (le), \zh{吧} (ba), \zh{呢} (ne), \zh{的} (de), \zh{着} (zhe dans \zh{省着花} shěngzhe huā).
\end{itemize}
\end{propriete}

\begin{attention}[Pièges tonaux de la leçon]
\begin{enumerate}
\item \textbf{\zh{不} (bù) + 4e ton = \zh{不} devient 2e ton (bú).} \rightarrow \zh{不应该} se prononce \emph{bú yīnggāi} (pas \emph{bù yīnggāi}). Idem : \zh{不是} (bú shì), \zh{不要} (bú yào).
\item \textbf{Deux 3e tons de suite : le premier devient 2e ton.} \zh{省着} (shěng zhe $\to$ shéng zhe). \zh{记账} (jì zhàng $\to$ jí zhàng). \zh{很好} (hěn hǎo $\to$ hén hǎo).
\item \textbf{\zh{一} (yī) :} 1er ton seul/chiffre ; 4e ton (yì) devant 1,2,3e tons ; 2e ton (yí) devant 4e ton. \zh{一万} (yí wàn), \zh{一笔} (yì bǐ), \zh{一样} (yí yàng).
\item \textbf{\zh{零花钱} :} \emph{líng huā qián} (2-1-2). Attention à ne pas dire \emph{lǐng} (3e ton, « commander »).
\item \textbf{\zh{预算} :} \emph{yù suàn} (4-4). \zh{算} est 4e ton ici (calculer), pas 3e ton.
\end{enumerate}
\end{attention}

\begin{experience}[Entraînement ciblé sur les tons (5 min)]
En binôme, lisez la liste ci-dessous. L'écouteur corrige les tons (surtout \zh{不应该}, \zh{省着花}, \zh{记账}, \zh{一万}, \zh{零花钱}). Changez de rôle.
\begin{enumerate}
\item \zh{你不应该乱花钱。} (bú yīnggāi luàn huā qián)
\item \zh{我一定省着花。} (wǒ yīdìng shéngzhe huā)
\item \zh{请给我一万法郎。} (qǐng gěi wǒ yí wàn fǎláng)
\item \zh{他每天记账。} (tā měitiān jí zhàng)
\item \zh{零花钱不多。} (línghuāqián bù duō)
\item \zh{因为预算紧张，所以我们省着花。} (yīnwèi yùsuàn jǐnzhāng, suǒyǐ wǒmen shéngzhe huā)
\end{enumerate}
\end{experience}

\coursec{Jeux de rôle et débat guidé}

Maintenant que le vocabulaire, les dialogues, les formules et les tons sont en place, passons à la pratique interactive. Cette section propose deux jeux de rôle et un débat structuré pour mobiliser toutes les compétences en situation de communication quasi-réelle.

\courssub{Jeu de rôle 1 : « À la fin du mois »}

\textbf{Situation.} C'est le 28 du mois. L'élève A a dépensé tout son argent de poche. L'élève B lui donne des conseils avisés. La structure grammaticale centrale est \cle{应该} (\emph{yīnggāi}) et sa négation \cle{不应该} (\emph{bú yīnggāi}).

\begin{propriete}[La structure 应该 / 不应该 + Verbe]
\textbf{Forme :} Sujet + (不)应该 + Verbe (+ Complément).\par
\textbf{Emploi :} Exprimer un conseil, une obligation morale, une recommandation. Équivalent de « tu devrais / tu ne devrais pas ».\par
\textbf{Négation :} \zh{不} se place \emph{avant} \zh{应该} $\to$ \zh{不应该} (prononcé \emph{bú yīnggāi}). \textbf{Jamais} \zh{应该不}.\par
\textbf{Nuance :} \zh{应该} est plus fort que \zh{建议你} (conseil) mais moins contraignant que \zh{得} (děi, obligation forte/nécessité).
\end{propriete}

\begin{exemplebox}[Dialogue modèle à imiter et transformer]
\zh{A：我这个月的零花钱都花完了！}\par
\emph{Wǒ zhège yuè de línghuāqián dōu huā wán le!} — « J'ai dépensé tout mon argent de poche du mois ! »\par
\zh{B：你不应该买那么多零食。你应该记账。}\par
\emph{Nǐ bú yīnggāi mǎi nàme duō língshí. Nǐ yīnggāi jìzhàng.} — « Tu ne devrais pas acheter autant de snacks. Tu devrais tenir tes comptes. »\par
\zh{A：你说得对。下个月我一定省着花。}\par
\emph{Nǐ shuō de duì. Xià ge yuè wǒ yīdìng shéngzhe huā.} — « Tu as raison. Le mois prochain, je ferai attention, c'est sûr. »
\end{exemplebox}

\begin{experience}[À vous de jouer ! — Jeu de rôle 1]
Par groupes de deux. Jouez la scène \textbf{trois fois} en changeant le motif de la ruine :
\begin{enumerate}
\item Argent parti en \zh{游戏皮肤} (skins de jeux vidéo).
\item Argent parti en \zh{话费} (crédit téléphonique / forfait).
\item Argent parti en \zh{聚餐} (repas entre amis / sorties).
\end{enumerate}
\textbf{Contrainte :} Le rôle B doit donner \textbf{au moins trois conseils} distincts utilisant \zh{你应该...} et \zh{你不应该...}. Puis inversez les rôles.
\textbf{Pistes :} \zh{你应该先买学习用品。} (\emph{Nǐ yīnggāi xiān mǎi xuéxí yòngpǐn.}) ; \zh{你不应该每天打车。} (\emph{Nǐ bú yīnggāi měitiān dǎchē.}) ; \zh{你应该少喝奶茶。} (\emph{Nǐ yīnggāi shǎo hē nǎichá.}) ; \zh{你不应该借钱给别人。} (\emph{Nǐ bú yīnggāi jiè qián gěi biérén.}).
\end{experience}

\courssub{Jeu de rôle 2 : « Le conseil de famille » (Yaoundé)}

\textbf{Situation.} Une famille de Yaoundé dispose d'un budget fictif de \zh{十万法郎} (100 000 FCFA) pour le mois. En groupes de quatre (Père \zh{爸爸}, Mère \zh{妈妈}, Fils \zh{儿子}, Fille \zh{女儿}), vous devez répartir cette somme entre quatre postes obligatoires : \zh{吃饭} (nourriture), \zh{上学/交通} (école/transport), \zh{学习用品} (fournitures), \zh{娱乐} (loisirs). Chaque membre défend son poste !

\begin{exemplebox}[Phrases utiles pour le conseil de famille]
\zh{我觉得吃饭最重要，应该给吃饭四万。}\par
\emph{Wǒ juéde chīfàn zuì zhòngyào, yīnggāi gěi chīfàn sì wàn.} — « Je pense que la nourriture est prioritaire, il faut allouer 40 000. »\par
\zh{可是交通也不能太少，我们每天要坐车上学。}\par
\emph{Kěshì jiāotōng yě bù néng tài shǎo, wǒmen měitiān yào zuò chē shàngxué.} — « Mais le transport ne peut pas être trop faible, on prend le bus tous les jours. »\par
\zh{娱乐一万够了，剩下的存起来。}\par
\emph{Yúlè yí wàn gòu le, shèng xià de cún qǐlai.} — « 10 000 pour les loisirs suffisent, le reste on l'épargne. »\par
\zh{好，那就这么定吧！记得记账。}\par
\emph{Hǎo, nà jiù zhème dìng ba! Jìde jìzhàng.} — « Bon, c'est décidé comme ça ! N'oubliez pas de noter les dépenses. »
\end{exemplebox}

\begin{experience}[Conseil de famille — Déroulement]
\begin{enumerate}
\item \textbf{Préparation (2 min) :} Chaque membre note secrètement sa proposition de répartition (ex: Manger 40k, Transport 20k, Fournitures 20k, Loisirs 10k, Épargne 10k).
\item \textbf{Négociation (5 min) :} Tour de table. Chacun annonce sa proposition avec \zh{我觉得...应该给...} . Les autres réagissent avec \zh{可是...} / \zh{我同意...} / \zh{不太同意...} .
\item \textbf{Décision (1 min) :} Consensus ou vote. Le père/mère tranche : \zh{那就这么定了。} .
\item \textbf{Restitution orale (1 min/groupe) :} Un rapporteur présente le budget final à la classe en chinois.
\end{enumerate}
\end{experience}

\courssub{Débat guidé : « Faut-il donner de l'argent de poche aux lycéens ? »}

C'est l'activité sommative de la leçon. Elle structure la pensée argumentative en chinois.

\begin{textebox}[Trame du débat : Pour ou contre l'argent de poche ?]
\zh{主持人：今天我们讨论：应该不应该给中学生零花钱？}\par
\emph{Zhǔchírén: Jīntiān wǒmen tǎolùn: yīnggāi bu yīnggāi gěi zhōngxuéshēng línghuāqián?}\par
« L'animateur : Aujourd'hui, nous discutons : faut-il donner de l'argent de poche aux collégiens et lycéens ? »\par
\zh{甲方：我觉得学生应该有零花钱，因为他们可以学习管理钱。}\par
\emph{Jiǎfāng: Wǒ juéde xuésheng yīnggāi yǒu línghuāqián, yīnwèi tāmen kěyǐ xuéxí guǎnlǐ qián.}\par
« Camp Pour : Je pense que les élèves devraient avoir de l'argent de poche, car ils peuvent ainsi apprendre à gérer l'argent. »\par
\zh{乙方：我不太同意，因为有些学生有了钱就乱花。}\par
\emph{Yǐfāng: Wǒ bù tài tóngyì, yīnwèi yǒuxiē xuésheng yǒu le qián jiù luàn huā.}\par
« Camp Contre : Je ne suis pas tout à fait d'accord, car certains élèves, dès qu'ils ont de l'argent, le dépensent n'importe comment. »\par
\zh{甲方：你的话有道理，可是不给钱他们怎么学会管理？}\par
\emph{Jiǎfāng: Nǐ de huà yǒu dàoli, kěshì bù gěi qián tāmen zěnme xuéhuì guǎnlǐ?}\par
« Camp Pour : Ton argument tient la route, mais sans argent, comment apprendre à gérer ? »\par
\zh{乙方：可以给，但要少给，还要教他们记账。}\par
\emph{Yǐfāng: Kěyǐ gěi, dàn yào shǎo gěi, hái yào jiāo tāmen jìzhàng.}\par
« Camp Contre : On peut en donner, mais peu, et il faut leur apprendre à tenir les comptes. »\par
\zh{报告员：所以，我们认为应该给零花钱，可是不能太多，还要学会记账。}\par
\emph{Bàogàoyuán: Suǒyǐ, wǒmen rènwéi yīnggāi gěi línghuāqián, kěshì bù néng tài duō, hái yào xuéhuì jìzhàng.}\par
« Le rapporteur : Donc, nous pensons qu'il faut donner de l'argent de poche, mais pas trop, et qu'il faut apprendre à tenir ses comptes. »
\end{textebox}

\begin{propriete}[La structure 因为...所以... (Yīnwèi... Suǒyǐ...)]
\textbf{Forme :} \zh{因为} + Cause ，\zh{所以} + Conséquence.\par
\textbf{Règle d'or :} En chinois, les deux marqueurs \textbf{coexistent fréquemment dans la même phrase complexe}. C'est correct et même recommandé pour marquer la logique cause-effet. En français, « Parce que... donc... » est une lourdeur/faute ; on choisit l'un des deux.\par
\textbf{Variantes :} On peut n'utiliser que \zh{因为} (en début de phrase, la conséquence est implicite) ou que \zh{所以} (la cause est contextuelle). Mais on ne sépare jamais les deux propositions par un point final si on garde les deux mots.
\end{propriete}

\begin{exemplebox}[因为...所以... en action : variations]
\zh{因为钱不多，所以我们应该省着花。}\par
\emph{Yīnwèi qián bù duō, suǒyǐ wǒmen yīnggāi shéngzhe huā.} — « Comme l'argent est limité, nous devons faire attention. »\par
\zh{因为他每天记账，所以他知道钱花在哪儿。}\par
\emph{Yīnwèi tā měitiān jí zhàng, suǒyǐ tā zhīdào qián huā zài nǎr.} — « Parce qu'il note ses comptes chaque jour, il sait où part son argent. »\par
\zh{因为交通费太贵了，所以我走路去学校。}\par
\emph{Yīnwèi jiāotōngfèi tài guì le, suǒyǐ wǒ zǒulù qù xuéxiào.} — « Comme le transport coûte cher, j'y vais à pied. »\par
\zh{零花钱能教我们理财，所以它是必要的。} (Seul \zh{所以} explicite, cause en sujet)\par
\emph{Línghuāqián néng jiāo wǒmen lǐcái, suǒyǐ tā shì bìxū de.} — « L'argent de poche nous apprend la finance, donc c'est nécessaire. »
\end{exemplebox}

\begin{attention}[Piège majeur pour francophones : « Parce que... donc »]
\begin{itemize}
\item \textbf{Français :} \emph{Parce que} je suis occupé, \emph{donc} je ne suis pas venu $\to$ \textbf{FAUTE}. On dit : « \emph{Comme} je suis occupé, je ne suis pas venu » OU « Je suis occupé, \emph{donc} je ne suis pas venu ».
\item \textbf{Chinois :} \zh{因为我很忙，所以没去。} (\emph{Yīnwèi wǒ hěn máng, suǒyǐ méi qù.}) $\to$ \textbf{CORRECT et NATUREL}.
\item \textbf{Traduction Version (Ch $\to$ Fr) :} Supprimez l'un des deux connecteurs.
\item \textbf{Traduction Thème (Fr $\to$ Ch) :} Gardez les deux (\zh{因为}... \zh{所以}...).
\end{itemize}
\end{attention}

\courssub{Les rôles du débat et la méthode pour le réussir}

Un débat efficace en chinois repose sur une distribution claire des rôles et une procédure ritualisée.

\begin{center}
\begin{tabularx}{\linewidth}{|X|X|X|}
\hline
\rowcolor{popblueL} Rôle & Formule clé (Caractères + Pinyin) & Fonction \\
\hline
Animateur \zh{主持人} & \zh{你呢？你怎么看？} \emph{Nǐ ne? Nǐ zěnme kàn?} & Lancer le sujet, distribuer la parole, relancer si silence \\
\hline
Orateur « Pour » \zh{甲方} & \zh{我觉得...因为...} \emph{Wǒ juéde... yīnwèi...} & Présenter les arguments favorables (min. 2) \\
\hline
Orateur « Contre » \zh{乙方} & \zh{我不太同意，因为...} \emph{Wǒ bù tài tóngyì, yīnwèi...} & Présenter les arguments défavorables (min. 2) \\
\hline
Rapporteur \zh{报告员} & \zh{所以，我们认为...} \emph{Suǒyǐ, wǒmen rènwéi...} & Synthétiser les deux camps, formuler la conclusion commune \\
\hline
Observateurs (x2) & Grille d'évaluation (voir plus bas) & Noter tons, formules, fluidité, interaction \\
\hline
\end{tabularx}
\end{center}

\begin{methode}[Réussir un débat en chinois en cinq étapes]
\begin{enumerate}
\item \textbf{Écouter activement} l'autre jusqu'au bout, sans couper. Repérez sa structure (\zh{我觉得...}, \zh{因为...}).
\item \textbf{Valider avant de contrer} (politesse chinoise) : \zh{你说得对，可是...} (\emph{Nǐ shuō de duì, kěshì...}) ou \zh{有道理，但是...} (\emph{Yǒu dàoli, dànshì...}).
\item \textbf{Argumenter} avec sa propre structure : \zh{我不太同意，因为...} / \zh{我觉得...因为...所以...}.
\item \textbf{Relancer} si le débat marque le pas : L'animateur utilise \zh{你呢？} , \zh{还有别的想法吗？} , \zh{请说说你的看法。} (\emph{Qǐng shuōshuō nǐ de kànfǎ.}).
\item \textbf{Conclure collectivement} : Le rapporteur résume \zh{甲方认为... 乙方认为... 所以，我们认为...} .
\end{enumerate}
\end{methode}

\begin{popfigure}
\begin{tikzpicture}[
  node distance=0.55cm and 1.1cm,
  etape/.style={rectangle, rounded corners=3pt, draw=popblue, fill=popblueL, thick, align=center, minimum height=0.9cm, font=\small, text width=2.8cm},
  fleche/.style={-{Stealth[length=2.5mm]}, thick, popdark},
  retour/.style={-{Stealth[length=2.5mm]}, thick, dashed, poppurple}]
\node[etape, align=center] (ouv){Ouverture\\ \zh{今天我们讨论...}};
\node[etape, below=of ouv, fill=poppurpleL, draw=poppurple, align=center] (pour){Arguments POUR\\ \zh{我觉得...因为...}};
\node[etape, below=of pour, fill=poporangeL, draw=poporange, align=center] (contre){Arguments CONTRE\\ \zh{我不太同意，因为...}};
\node[etape, below=of contre, fill=popgoldL, draw=popgold, align=center] (relance){Relances / Échanges\\ \zh{你呢？\quad 你说得对，可是...}};
\node[etape, below=of relance, fill=popgreenL, draw=popgreen, align=center] (concl){Conclusion Rapporteur\\ \zh{所以，我们认为...}};
\draw[fleche] (ouv) -- (pour);
\draw[fleche] (pour) -- (contre);
\draw[fleche] (contre) -- (relance);
\draw[fleche] (relance) -- (concl);
\draw[retour] (relance.west) .. controls +(-1.8,0.5) and +(-1.8,-0.5) .. (pour.west);
\draw[retour] (relance.west) .. controls +(-1.8,1.2) and +(-1.8,-1.2) .. (contre.west);
\end{tikzpicture}
\legende{Organigramme du débat guidé. Les flèches pointillées montrent que les relances peuvent faire rebondir sur les arguments précédents.}
\end{popfigure}

\courssub{Grille d'observation et d'évaluation (Observateurs / Enseignant)}

Pendant le débat, deux élèves observateurs remplissent cette grille pour \textbf{chaque orateur} (Animateur, Pour, Contre, Rapporteur). Elle sert d'outil d'auto-évaluation et d'évaluation sommative.

\begin{center}
\begin{tabularx}{\linewidth}{|X|X|X|X|}
\hline
\rowcolor{popblueL} Critère & \textbf{很好} (Excellent) & \textbf{好} (Bien) & \textbf{加油} (À améliorer) \\
\hline
Prononciation des tons & Tons justes partout (y compris sandhi \zh{不应该}, \zh{省着花}) & 1-2 erreurs de ton non gênantes & Tons souvent confondus (3e/4e, neutre) \\
\hline
Formules fonctionnelles & Utilise 3 formules cibles ou plus (\zh{因为...所以...}, \zh{应该}, \zh{你呢?}, \zh{你说得对}) & Utilise 1 à 2 formules cibles & Aucune formule cible, français ou chinois très approximatif \\
\hline
Fluidité & Parle sans longue pause, débit naturel & Quelques hésitations / auto-corrections & Lit ses notes, longs silences, haché \\
\hline
Interaction / Relance & Répond à l'autre, relance (\zh{你呢?}), utilise \zh{你说得对，可是...} & Répond quand on l'invite, peu d'initiative & Reste silencieux, ne réagit pas aux autres \\
\hline
Contenu / Argumentation & Arguments pertinents, exemples concrets (Cameroun/Chine), structure claire & Arguments simples, un exemple, structure visible & Arguments absents ou hors sujet \\
\hline
\end{tabularx}
\end{center}

\begin{exoresolu}[Préparer une prise de parole type « Orateur Pour »]
\textbf{Énoncé.} Vous êtes l'orateur « Pour » (\zh{甲方}). Rédigez et prononcez une intervention structurée en trois phrases utilisant : 1. \zh{我觉得} (avis), 2. \zh{因为...所以...} (argumentation), 3. Un exemple concret (\zh{比如} / \zh{例如}).\par
\tcblower
\textbf{Solution détaillée.}\par
1. \textbf{Annoncer l'avis :} \zh{我觉得应该给高中生零花钱。} (\emph{Wǒ juéde yīnggāi gěi gāozhōngshēng línghuāqián.}) — « Je pense qu'il faut donner de l'argent de poche aux lycéens. »\par
2. \textbf{Justifier (Cause $\to$ Conséquence) :} \zh{因为零花钱能教我们理财，所以它是一种财商教育。} (\emph{Yīnwèi línghuāqián néng jiāo wǒmen lǐcái, suǒyǐ tā shì yì zhǒng cáishāng jiàoyù.}) — « Parce que l'argent de poche nous apprend la finance, c'est une éducation financière. »\par
3. \textbf{Exemple concret :} \zh{比如，我每个月记一次账，现在我不乱花钱了，还存了一千块。} (\emph{Bǐrú, wǒ měi ge yuè jì yí cì zhàng, xiànzài wǒ bù luàn huā qián le, hái cún le yì qiān kuài.}) — « Par exemple, je tiens mes comptes chaque mois, maintenant je ne gaspille plus, j'ai même économisé 1000 yuan. »\par
\textbf{Méthode universelle :} \cle{Avis} $\to$ \cle{Raison (因为...所以...)} $\to$ \cle{Exemple (比如...)}. Cette structure en trois temps convient à toute prise de parole argumentée (oral d'examen, débat, exposé).
\end{exoresolu}

\begin{attention}[Erreurs fréquentes à l'oral — Check-list finale]
\begin{itemize}
\item \textbf{Accord \zh{同意} :} Ne dites pas \zh{我同意你} (Chinglish). Dites : \zh{我同意你的看法} (\emph{Wǒ tóngyì nǐ de kànfǎ}), \zh{我赞成} (\emph{Wǒ zànchéng}), ou simplement \zh{你说得对} / \zh{有道理}.
\item \textbf{Homophones \zh{钱} vs \zh{前} :} Même pinyin \emph{qián} (2e ton). \zh{钱} = argent (radical \zh{钅} or). \zh{前} = devant/avant (radical \zh{⻌} marche). À l'écrit : vigilance absolue.
\item \textbf{\zh{觉得} + [Phrase] :} Après \zh{觉得}, il faut une proposition complète (Sujet + Verbe...), pas un nom seul. \checkmark \zh{我觉得零花钱很重要}. \cross \zh{我觉得零花钱}.
\item \textbf{\zh{省着花} (shěngzhe huā) :} \zh{着} (zhe) ton neutre, marque la manière/état. Ne pas dire \zh{省花} (incorrect).
\item \textbf{Classificateur \zh{笔} pour dépenses :} \zh{一笔钱}, \zh{三笔支出}. Pas \zh{个} pour une dépense comptable.
\end{itemize}
\end{attention}

\begin{savaistu}[Argent de poche, Mobile Money et WeChat Pay : Cameroun / Chine]
\textbf{Cameroun :} L'essor du \textbf{Mobile Money} (MTN MoMo, Orange Money) a transformé la gestion de l'argent de poche. Beaucoup de parents envoient les \zh{零花钱} directement sur le téléphone de l'élève. Chaque transaction (retrait, achat crédit, transfert) laisse une trace numérique : c'est un \zh{记账} (tenue de comptes) automatique ! Cela apprend la traçabilité.
\textbf{Chine :} Les lycéens chinois vivent dans une société \textbf{sans cash} (\zh{无现金社会}). \zh{微信支付} (Wēixìn zhīfù, WeChat Pay) et \zh{支付宝} (Zhīfùbǎo, Alipay) règnent : cantine (\zh{食堂}), bus (\zh{公交}), supérette (\zh{便利店}), petite rue (\zh{路边摊}). L'argent est devenu « invisible » (\zh{隐形的钱}).
\textbf{Comparaison :} Dans les deux pays, la dématérialisation rend la dépense trop facile (\zh{太容易花钱了}). D'où l'urgence pédagogique de cette leçon : \zh{记账} n'est pas une corvée d'autrefois, c'est une compétence de survie financière moderne. Apprendre à \zh{省着花} (dépenser avec discernement) face à la facilité du paiement mobile est le vrai défi commun des jeunes de Yaoundé et de Pékin.
\end{savaistu}

\coursec{Production orale guidée et évaluation}

\courssub{Exercice sommatif : Débat organisé et évalué}

\begin{exercice}[Débat final : « Faut-il donner de l'argent de poche aux lycéens ? »]
\textbf{Consigne.} En groupes de \textbf{cinq} (1 Animateur, 2 Orateurs « Pour » \zh{甲方}, 2 Orateurs « Contre » \zh{乙方} — le Rapporteur \zh{报告员} est désigné parmi les 4 orateurs à la fin).
\textbf{Déroulement (6 min total) :}
\begin{enumerate}
\item \textbf{1 min} : Animateur lance le sujet (\zh{今天我们讨论...}).
\item \textbf{3 min} : Alternance arguments Pour / Contre. Chaque orateur doit utiliser \textbf{au moins une fois} \zh{因为...所以...} et \textbf{une formule de politesse/débat} (\zh{你说得对，可是...} / \zh{我不太同意} / \zh{有道理，但是...}).
\item \textbf{1 min} : Relances libres (Animateur : \zh{你呢？}, \zh{还有想法吗？}).
\item \textbf{1 min} : Rapporteur conclut : \zh{甲方认为... 乙方认为... 所以，我们认为...}.
\end{enumerate}
\textbf{Contraintes linguistiques obligatoires :} Vocabulaire budget (\zh{预算}, \zh{分配}, \zh{记账}, \zh{省着花}), Structures (\zh{应该/不应该}, \zh{因为...所以...}, \zh{觉得/认为}), Tons soignés (\zh{bú yīnggāi}, \zh{shéngzhe huā}, \zh{jí zhàng}).
\textbf{Rôle des observateurs :} Remplir la grille d'évaluation pour chaque orateur. Remettre à l'enseignant.
\end{exercice}

\begin{corrige}[Exercice — Débat final : Grille de réussite attendue]
Il n'y a pas de réponse unique ; l'évaluation porte sur la \textbf{méthode} et la \textbf{correction linguistique}. Une prestation réussie valide :
\begin{enumerate}
\item \textbf{Structure respectée :} Ouverture $\to$ Arguments Pour (min 2 args distincts) $\to$ Arguments Contre (min 2 args distincts) $\to$ Relances $\to$ Conclusion synthétique.
\item \textbf{Outils grammaticaux :} Chaque camp a utilisé \zh{因为...所以...} au moins une fois. \zh{应该/不应该} utilisé pour conseiller. \zh{觉得/认为} pour l'opinion.
\item \textbf{Interaction :} Les réponses commencent par \zh{你说得对，可是...} ou \zh{我不太同意，因为...} (pas de monologues juxtaposés). L'animateur a relancé au moins deux fois (\zh{你呢？}).
\item \textbf{Conclusion du rapporteur :} Commence par \zh{所以}. Résume les deux positions. Prend une position nuancée.
\item \textbf{Exemple de conclusion attendue (niveau Terminale A) :}
\zh{所以，我们认为应该给高中生零花钱，因为这是学习理财的机会。可是不能给太多，父母要教我们记账，省着花。} \\
\emph{Suǒyǐ, wǒmen rènwéi yīnggāi gěi gāozhōngshēng línghuāqián, yīnwèi zhè shì xuéxí lǐcái de jīhuì. Kěshì bù néng gěi tài duō, fùmǔ yào jiāo wǒmen jìzhàng, shéngzhe huā.} \\
« Donc, nous pensons qu'il faut donner de l'argent de poche aux lycéens, car c'est une occasion d'apprendre la finance. Mais il ne faut pas en donner trop, les parents doivent nous apprendre à tenir nos comptes et à dépenser avec sagesse. »
\end{enumerate}
\end{corrige}

\courssub{Activité de réinvestissement : « Mon budget idéal » (Expression suivie / Devoir)}

\begin{exercice}[Rédaction / Exposé oral enregistré : « Mon budget mensuel idéal »]
\textbf{Consigne.} Imaginez que vous avez \zh{2000元} (ou 150 000 FCFA) par mois. Rédigez un court texte (80-100 caractères) ou préparez un exposé oral de 1 minute présentant :
\begin{enumerate}
\item Vos revenus (\zh{收入} : \zh{零花钱}, \zh{兼职} job étudiant, \zh{红包}).
\item La répartition de votre budget (\zh{分配}) : \zh{吃饭}, \zh{交通}, \zh{学习用品}, \zh{娱乐}, \zh{存钱}. Utilisez \zh{给...} + montant.
\item Une règle de gestion personnelle avec \zh{应该} / \zh{不应该}.
\item Une conclusion avec \zh{因为...所以...}.
\end{enumerate}
\textbf{Modèle de départ :}
\zh{我的收入是两千元。我给吃饭八百元，交通三百元，学习用品两百元，娱乐两百元，存五百元。我觉得存钱很重要，因为以后可能急用钱，所以我每个月都要存。我不应该乱买零食，应该省着花。}
\emph{Wǒ de shōurù shì liǎng qiān yuán. Wǒ gěi chīfàn bā bǎi yuán, jiāotōng sān bǎi yuán, xuéxí yòngpǐn liǎng bǎi yuán, yúlè liǎng bǎi yuán, cún wǔ bǎi yuán. Wǒ juéde cún qián hěn zhòngyào, yīnwèi yǐhòu kěnéng jí yòng qián, suǒyǐ wǒ měi ge yuè dōu yào cún. Wǒ bú yīnggāi luàn mǎi língshí, yīnggāi shéngzhe huā.}
\end{exercice}

\begin{corrige}[Exercice — Mon budget idéal]
\textbf{Critères de réussite :}
\begin{itemize}
\item Vocabulaire du budget utilisé correctement (min. 6 items du tableau Sec. 1).
\item Classificateurs corrects (\zh{元/块}, \zh{笔}, \zh{份}).
\item Structures cibles présentes : \zh{给} + Poste + Montant ; \zh{应该/不应该} + V ; \zh{因为...所以...}.
\item Tons respectés à l'oral (enregistrement) / Caractères corrects à l'écrit (ordre des traits, pas de \zh{帐} pour \zh{账}).
\item Cohérence du propos (total des dépenses $\le$ revenu).
\end{itemize}
\end{corrige}

\coursec{Leçon 28 — Expression orale : Gestion du budget}

\courssub{Vocabulaire du budget : revenus, dépenses et gestion}

\begin{textebox}[Situation déclenchante : Discussion à la cantine]
\zh{— 小张，你这个月零花钱花完了吗？} \\
\zh{— 唉，别提了。上周我把钱都花在买新球鞋上了，现在连吃饭的钱都不够。} \\
\zh{— 你得学会做预算。我每个月先存一部分，剩下的才花。} \\
\zh{— 说得容易。下个月我一定省钱，不买零食了。}
\end{textebox}

\begin{center}
\begin{tabularx}{\linewidth}{|X|X|X|X|}
\hline
\rowcolor{popblueL} \textbf{Caractères} & \textbf{Pinyin} & \textbf{Nature} & \textbf{Traduction} \\
\hline
\zh{零花钱} & línghuāqián & nom & argent de poche \\
\hline
\zh{生活费} & shēnghuófèi & nom & frais de subsistance (logement, nourriture) \\
\hline
\zh{兼职} & jiānzhí & verbe/nom & job étudiant, travail à temps partiel \\
\hline
\zh{红包} & hóngbāo & nom & enveloppe rouge (cadeau d'argent) \\
\hline
\zh{收入} & shōurù & nom & revenu, rentrée d'argent \\
\hline
\zh{支出} & zhīchū & nom & dépense, sortie d'argent \\
\hline
\hline
\zh{交通} & jiāotōng & nom & transport \\
\hline
\zh{餐饮} & cānyǐn & nom & restauration, repas (formel) \\
\hline
\zh{吃饭} & chīfàn & verbe/nom & manger / repas (courant) \\
\hline
\zh{文具} & wénjù & nom & papeterie, fournitures scolaires \\
\hline
\zh{衣服} & yīfu & nom & vêtements \\
\hline
\zh{零食} & língshí & nom & snacks, grignotines \\
\hline
\zh{手机费} & shǒujīfèi & nom & forfait / crédit téléphone \\
\hline
\zh{网费} & wǎngfèi & nom & frais d'internet \\
\hline
\hline
\zh{花} & huā & verbe & dépenser (argent, temps) \\
\hline
\zh{省} & shěng & verbe & économiser, épargner \\
\hline
\zh{省钱} & shěngqián & verbe & économiser de l'argent \\
\hline
\zh{存} & cún & verbe & mettre de côté, déposer (banque) \\
\hline
\zh{存钱} & cúnqián & verbe & mettre de l'argent de côté \\
\hline
\zh{赚} & zhuàn & verbe & gagner (de l'argent) \\
\hline
\zh{够用} & gòu yòng & verbe & suffire, être suffisant \\
\hline
\zh{不够用} & bù gòu yòng & verbe & ne pas suffire, manquer \\
\hline
\hline
\zh{预算} & yùsuàn & nom/verbe & budget / budgétiser \\
\hline
\zh{计划} & jìhuà & nom/verbe & plan / planifier \\
\hline
\zh{块} & kuài & classif. & yuan (oral, familier) \\
\hline
\zh{元} & yuán & classif./nom & yuan (écrit, formel) \\
\hline
\zh{笔} & bǐ & classif. & trait / somme (d'argent), article (dépense) \\
\hline
\end{tabularx}
\end{center}

\begin{propriete}[Classificateurs monétaires : \zh{块} vs \zh{元} vs \zh{笔}]
\begin{itemize}
    \item \zh{块} : Unité orale courante (1 \zh{块} = 1 \zh{元} = 10 \zh{毛} = 100 \zh{分}). \emph{Ex:} \zh{两百块} (200 yuans).
    \item \zh{元} : Unité écrite, formelle, bancaire. \emph{Ex:} \zh{¥200元}.
    \item \zh{笔} : Classificateur pour une transaction, une somme dépensée ou reçue. \emph{Ex:} \zh{一笔收入} (une rentrée d'argent), \zh{三笔支出} (trois postes de dépenses).
\end{itemize}
\end{propriete}

\begin{popfigure}
\begin{tikzpicture}[
    mindmap, grow cyclic, every node/.style=concept, concept color=popblue!60,
    level 1/.append style={level distance=4.5cm, sibling angle=90, concept color=popblue!50},
    level 2/.append style={level distance=3.2cm, sibling angle=45, concept color=popgreen!50},
    level 3/.append style={level distance=2.8cm, sibling angle=30, concept color=poporange!50},
    text=popdark, font=\small
]
\node[align=center]{Budget\\ \zh{预算}}
    child [concept color=popblue!50] {node[align=center]{Revenus\\ \zh{收入}}
        child {node[align=center]{Argent de poche\\ \zh{零花钱}}}
        child {node[align=center]{Frais de vie\\ \zh{生活费}}}
        child {node[align=center]{Job étudiant\\ \zh{兼职}}}
        child {node[align=center]{Enveloppe rouge\\ \zh{红包}}}
    }
    child [concept color=popgreen!50] {node[align=center]{Dépenses\\ \zh{支出}}
        child {node[align=center]{Transport\\ \zh{交通}}}
        child {node[align=center]{Repas\\ \zh{餐饮/吃饭}}}
        child {node[align=center]{Fournitures\\ \zh{文具}}}
        child {node[align=center]{Vêtements\\ \zh{衣服}}}
        child {node[align=center]{Snacks\\ \zh{零食}}}
        child {node[align=center]{Téléphone/Net\\ \zh{手机费/网费}}}
    }
    child [concept color=poporange!50] {node[align=center]{Actions\\ \zh{动作}}
        child {node[align=center]{Dépenser\\ \zh{花}}}
        child {node[align=center]{Économiser\\ \zh{省/省钱}}}
        child {node[align=center]{Mettre de côté\\ \zh{存/存钱}}}
        child {node[align=center]{Gagner\\ \zh{赚}}}
    }
    child [concept color=poppurple!50] {node[align=center]{État\\ \zh{状况}}
        child {node[align=center]{Suffisant\\ \zh{够用}}}
        child {node[align=center]{Insuffisant\\ \zh{不够用}}}
        child {node[align=center]{Équilibré\\ \zh{收支平衡}}}
        child {node[align=center]{Déficit\\ \zh{入不敷出}}}
    };
\end{tikzpicture}
\legende{Carte mentale du vocabulaire budgétaire (HSK 3-4).}
\end{popfigure}

\courssub{Dialogues modèles : situations authentiques}

\begin{textebox}[Dialogue 1 : Comparer les budgets (Cameroon Context)]
\zh{— 阿明，你每个月有多少零花钱？} \\
\zh{— 我每个月有三千块，是父母给的生活费。你呢？} \\
\zh{— 我只有两千块。有时候我做家教，能赚五百块。} \\
\zh{— 那你一个月有两千五？不错！你把钱花在哪儿？} \\
\zh{— 主要花在交通和吃饭上。一笔交通费五百，吃饭一千多。} \\
\zh{— 我把钱主要花在买书和网费上。书太贵了，一笔就两百多。} \\
\zh{— 所以我们都得省钱。下学期我打算少买零食，存点钱买电脑。} \\
\zh{— 好主意！我们互相监督吧。}
\end{textebox}

\begin{center}
\begin{tabularx}{\linewidth}{|X|X|X|}
\hline
\rowcolor{popblueL} \textbf{Chinois} & \textbf{Pinyin} & \textbf{Traduction} \\
\hline
\zh{阿明，你每个月有多少零花钱？} & Āmíng, nǐ měi gè yuè yǒu duōshao línghuāqián? & Amine, combien as-tu d'argent de poche par mois ? \\
\hline
\zh{我每个月有三千块，是父母给的生活费。} & Wǒ měi gè yuè yǒu sān qiān kuài, shì fùmǔ gěi de shēnghuófèi. & J'ai 3000 yuans/mois, c'est les frais de vie que mes parents donnent. \\
\hline
\zh{我只有两千块。有时候我做家教，能赚五百块。} & Wǒ zhǐ yǒu liǎng qiān kuài. Yǒushíhou wǒ zuò jiājiào, néng zhuàn wǔ bǎi kuài. & Je n'ai que 2000. Parfois je fais du soutien scolaire, je peux gagner 500. \\
\hline
\zh{你把钱花在哪儿？} & Nǐ bǎ qián huā zài nǎr? & Où dépenses-tu ton argent ? (\zh{哪儿} = \zh{什么地方}) \\
\hline
\zh{一笔交通费五百，吃饭一千多。} & Yī bǐ jiāotōngfèi wǔ bǎi, chīfàn yī qiān duō. & Un poste transport 500, manger 1000+. \\
\hline
\zh{书太贵了，一笔就两百多。} & Shū tài guì le, yī bǐ jiù liǎng bǎi duō. & Les livres sont trop chers, un seul coûte 200+. \\
\hline
\zh{我们互相监督吧。} & Wǒmen hùxiāng jiāndū ba. & Surveillons-nous mutuellement / Motivons-nous ensemble. \\
\hline
\end{tabularx}
\end{center}

\begin{textebox}[Dialogue 2 : Négocier / Expliquer un dépassement]
\zh{— 妈妈，这个月钱不够用了，能不能再给我两百块？} \\
\zh{— 上个月不是给了你三百块吗？钱都花哪儿了？} \\
\zh{— 因为同学过生日，我买了礼物，又请客吃饭，所以超支了。} \\
\zh{— 下不为例。以后大额支出要跟我们商量。现在要学会记账。} \\
\zh{— 知道了。我以后每周记一次账，争取下个月收支平衡。}
\end{textebox}

\begin{center}
\begin{tabularx}{\linewidth}{|X|X|X|}
\hline
\rowcolor{popblueL} \textbf{Chinois} & \textbf{Pinyin} & \textbf{Traduction} \\
\hline
\zh{这个月钱不够用了，能不能再给我两百块？} & Zhège yuè qián bù gòu yòng le, néng bu néng zài gěi wǒ liǎng bǎi kuài? & Ce mois-ci l'argent ne suffit pas, peux-tu m'en redonner 200 ? \\
\hline
\zh{钱都花哪儿了？} & Qián dōu huā nǎr le? & L'argent a été dépensé où ? (\zh{了} = action passée/terminée) \\
\hline
\zh{因为同学过生日，我买了礼物，又请客吃饭，所以超支了。} & Yīnwèi tóngxué guò shēngrì, wǒ mǎi le lǐwù, yòu qǐngkè chīfàn, suǒyǐ chāozhī le. & Parce que c'était l'anniv d'un camarade, j'ai acheté un cadeau et invité à manger, donc dépassement de budget. \\
\hline
\zh{下不为例。} & Xià bù wéi lì. & Pas de prochaine fois / Ça ne se reproduira plus. (Chengyu oral) \\
\hline
\zh{大额支出要跟我们商量。} & Dà'é zhīchū yào gēn wǒmen shāngliáng. & Les grosses dépenses doivent être discutées avec nous. \\
\hline
\zh{学会记账} & xuéhuì jìzhàng & apprendre à tenir les comptes / noter ses dépenses. \\
\hline
\zh{收支平衡} & shōuzhī pínghéng & équilibre budgétaire (revenus = dépenses). \\
\hline
\end{tabularx}
\end{center}

\begin{aretenir}[Structures utiles des dialogues]
\begin{center}
\begin{tabularx}{\linewidth}{|X|X|X|}
\hline
\rowcolor{popgreenL} \textbf{Structure} & \textbf{Exemple (Caractères + Pinyin)} & \textbf{Traduction} \\
\hline
\zh{能不能} + V + \zh{？} & \zh{能不能再给我两百块？} (Néng bu néng zài gěi wǒ liǎng bǎi kuài?) & Peux-tu / Est-il possible de... ? (Politesse) \\
\hline
\zh{因为}..., \zh{所以}... & \zh{因为买礼物，所以超支了。} (Yīnwèi mǎi lǐwù, suǒyǐ chāozhī le.) & Parce que..., donc... (Corrélatif obligatoire) \\
\hline
\zh{又}... \zh{又}... & \zh{又买礼物，又请客吃饭。} (Yòu mǎi lǐwù, yòu qǐngkè chīfàn.) & À la fois... et... (Énumération d'actions passées) \\
\hline
\zh{下不为例} & \zh{下不为例。} (Xià bù wéi lì.) & Il n'y aura pas de prochaine fois (Engagement formel). \\
\hline
\zh{争取} + Objectif & \zh{争取下个月收支平衡。} (Zhēngqǔ xià gè yuè shōuzhī pínghéng.) & Viser / S'efforcer d'atteindre (objectif). \\
\hline
\end{tabularx}
\end{center}
\end{aretenir}

\courssub{Formules fonctionnelles pour l'oral : Présenter, Argumenter, Interagir}

\begin{definition}[Boîte à outils : Les formules par fonction communicative]
Ces formules sont les "briques" de votre exposé et de votre interaction. Mémorisez-les par blocs fonctionnels.
\end{definition}

\begin{center}
\begin{tabularx}{\linewidth}{|X|X|X|}
\hline
\rowcolor{popblueL} \textbf{Fonction} & \textbf{Formules chinoises (Caractères + Pinyin)} & \textbf{Français} \\
\hline
\textbf{Ouvrir / Annoncer le sujet} & \zh{大家好，我叫... 今天我要介绍一下我的零花钱情况。} (Dàjiā hǎo, wǒ jiào... Jīntiān wǒ yào jièshào yīxià wǒ de línghuāqián qíngkuàng.) & Bonjour, je m'appelle... Aujourd'hui je vais présenter ma situation d'argent de poche. \\
& \zh{首先，先说一下我的收入。} (Shǒuxiān, xiān shuō yīxià wǒ de shōurù.) & Premièrement, disons mes revenus. \\
\hline
\textbf{Structurer (Séquence)} & \zh{先}... \zh{，然后}... \zh{，接着}... \zh{，最后}... (Xiān... ránhòu... jiēzhe... zuìhòu...) & D'abord... Ensuite... Puis... Enfin... \\
\hline
\textbf{Structurer (Contraste)} & \zh{但是}... / \zh{不过}... / \zh{可是}... (Dànshì... / Bùguò... / Kěshì...) & Mais / Cependant / Toutefois. \\
\hline
\textbf{Justifier (Cause $\rightarrow$ Conséquence)} & \zh{因为}... \zh{，所以}... (Yīnwèi..., suǒyǐ...) & Parce que..., donc... (\textbf{Les deux sont obligatoires}). \\
\hline
\textbf{Exprimer l'opinion} & \zh{我觉得}... / \zh{我认为}... / \zh{在我看来}... (Wǒ juéde... / Wǒ rènwéi... / Wǒ kàn lái...) & Je pense que... / Je crois que... / À mon avis... \\
& \zh{应该} / \zh{得} / \zh{要} + V & Devoir / Falloir (Conseil, obligation morale). \\
\hline
\textbf{Exprimer l'intention / Projet} & \zh{我打算}... / \zh{我计划}... / \zh{我准备}... (Wǒ dǎsuàn... / Wǒ jìhuà... / Wǒ zhǔnbèi...) & Je prévois de... / Je compte... / Je me prépare à... \\
\hline
\textbf{Interagir (Questions)} & \zh{你呢？} / \zh{为什么？} / \zh{怎么做？} / \zh{多少钱？} (Nǐ ne? / Wèishéme? / Zěnme zuò? / Duōshao qián?) & Et toi ? / Pourquoi ? / Comment faire ? / Combien ? \\
\hline
\textbf{Interagir (Réactions)} & \zh{有道理。} / \zh{我同意。} / \zh{我不太同意，因为...} (Yǒu dàoli. / Wǒ tóngyì. / Wǒ bù tài tóngyì, yīnwèi...) & C'est logique. / Je suis d'accord. / Je ne suis pas tout à fait d'accord, car... \\
\hline
\textbf{Clore} & \zh{以上就是我的介绍，谢谢大家。} (Yǐshàng jiùshì wǒ de jièshào, xièxie dàjiā.) & Voilà mon exposé, merci à tous. \\
& \zh{我的发言完了，谢谢。} (Wǒ de fāyán wánle, xièxie.) & Mon intervention est terminée, merci. \\
\hline
\end{tabularx}
\end{center}

\begin{attention}[Pièges francophones : \zh{觉得} vs \zh{认为} vs \zh{以为}]
\begin{itemize}
    \item \zh{觉得} (juéde) : Sentiment, impression subjective, souvent oral. \emph{Ex:} \zh{我觉得这本书很好看。}
    \item \zh{认为} (rènwéi) : Opinion réfléchie, jugement, plus formel/écrit. \emph{Ex:} \zh{我认为学习中文很重要。}
    \item \zh{以为} (yǐwéi) : \textbf{Faux ami !} Signifie "Croire à tort que...", "Penser que... (mais la réalité est autre)". \emph{Ex:} \zh{我以为考试在明天，其实在后天。} (Je croyais que l'examen était demain, en fait il est après-demain). \textbf{N'utilisez jamais \zh{以为} pour donner votre avis.}
\end{itemize}
\end{attention}

\courssub{Les tons à travailler : points clés et exercices ciblés}

La maîtrise des tons est notée sur 6 points (Grille Bac). Travaillez en priorité les mots-outils structurels et le vocabulaire technique de la leçon.

\begin{propriete}[Rappel des 4 tons + Ton neutre]
\begin{itemize}
    \item \textbf{1er ton (─) :} Haut, plat, long. \emph{Ex:} \zh{花} (huā), \zh{钱} (qián), \zh{先} (xiān).
    \item \textbf{2e ton (/) :} Montant, interrogatif. \emph{Ex:} \zh{零} (líng), \zh{来} (lái), \zh{赚} (zhuàn).
    \item \textbf{3e ton (V) :} Bas, descendant-remontant (souvent "demi-ton" bas en milieu de phrase). \emph{Ex:} \zh{省} (shěng), \zh{打算} (dǎsuàn), \zh{把} (bǎ).
    \item \textbf{4e ton (\textbackslash{}) :} Descendant, bref, affirmatif. \emph{Ex:} \zh{花} (huà - différent de huā !), \zh{够} (gòu), \zh{不} (bù - isolé), \zh{块} (kuài).
    \item \textbf{Ton neutre} : Court, léger, sans marque. \emph{Ex:} \zh{钱} dans \zh{零花钱} (línghuāqian), \zh{了} (le), \zh{吧} (ba), \zh{呢} (ne).
\end{itemize}
\end{propriete}

\begin{propriete}[Changements de ton obligatoires (Sandhi) — \textbf{Notés à l'oral}]
\begin{enumerate}
    \item \textbf{\zh{一} (yī) :} Devant ton 4 $\rightarrow$ \textbf{2e ton (yí)}. Devant tons 1,2,3 $\rightarrow$ \textbf{4e ton (yì)}. Seul/chiffre/fin $\rightarrow$ 1er ton.
        \begin{itemize}
            \item \zh{一下} $\rightarrow$ \textbf{yíxià} (pas yīxià).
            \item \zh{一点} $\rightarrow$ \textbf{yìdiǎn}.
            \item \zh{一百} $\rightarrow$ \textbf{yìbǎi}.
        \end{itemize}
    \item \textbf{\zh{不} (bù) :} Devant ton 4 $\rightarrow$ \textbf{2e ton (bú)}. Devant tons 1,2,3 $\rightarrow$ 4e ton (bù).
        \begin{itemize}
            \item \zh{不够} $\rightarrow$ \textbf{búgòu} (pas bùgòu).
            \item \zh{不是} $\rightarrow$ \textbf{búshì}.
            \item \zh{不花} $\rightarrow$ \textbf{bùhuā}.
        \end{itemize}
    \item \textbf{Deux 3e tons consécutifs :} Le premier devient \textbf{2e ton}.
        \begin{itemize}
            \item \zh{省钱} (shěng + qián) $\rightarrow$ \textbf{shéngqián}.
            \item \zh{打算} (dǎ + suàn) $\rightarrow$ \textbf{dásuàn}.
            \item \zh{所以} (suǒ + yǐ) $\rightarrow$ \textbf{suóyǐ}.
        \end{itemize}
\end{enumerate}
\end{propriete}

\begin{experience}[Atelier tons : "Le budget tonal" (10 min)]
\begin{enumerate}
    \item \textbf{Écoute-Discrimination :} L'enseignant lit une liste (ex: \zh{xiān / xiàn}, \zh{ránhòu / ránhǒu}, \zh{búgòu / bùgòu}, \zh{shéngqián / shěngqián}). Élèves lèvent carte "1/2/3/4" ou "Juste/Faux".
    \item \textbf{Chaîne de tons :} En cercle. Élève 1 : \zh{先} (xiān). Élève 2 : \zh{然后} (ránhòu). Élève 3 : \zh{但是} (dànshì). Élève 4 : \zh{所以} (suǒyǐ). Accélérer le rythme. Sanction : recommencer si erreur de ton.
    \item \textbf{Phrases pièges :} Lire à voix haute en exagérant les sandhi :
    \zh{我先介绍一下。} (Wǒ xiān jièshào \textbf{yíxià}.)
    \zh{钱不够用。} (Qián \textbf{bú} gòu yòng.)
    \zh{我打算省钱。} (Wǒ \textbf{dá}suàn \textbf{shéng}qián.)
    \zh{因为所以。} (Yīnwèi \textbf{suó}yǐ.)
\end{enumerate}
\end{experience}

\begin{popfigure}
\begin{tikzpicture}[
    scale=0.9, transform shape,
    tone/.style={thick, draw=popblue, ->, >=Stealth},
    node/.style={circle, draw=popdark, fill=popblueL, minimum size=0.6cm, font=\small},
    label/.style={font=\tiny, text=popmuted, anchor=north}
]
% Axes temps (horizontal) vs hauteur (vertical)
\draw[->, popmuted] (0,0) -- (14,0) node[right, label] {Temps};
\draw[->, popmuted] (0,0) -- (0,4.5) node[above, label] {Hauteur};

% Tons isolés (points de repère)
\foreach \x/\name/\path in {
    1/1er/{(0.5,3.5) -- (1.5,3.5)},
    2/2e/{(2.5,1) -- (3.5,4)},
    3/3e/{(4.5,2) .. controls (5,0.5) and (6,0.5) .. (6.5,3)},
    4/4e/{(7.5,4) -- (8.5,1)}
} {
    \draw[tone, popblue!80] \path;
    \node[label] at (\x+0.5, 0) {\name\ ton};
}

% Sandhi 不 (bù) -> bú devant 4e ton (gòu)
\draw[tone, poporange, dashed] (9.5,1) -- (10.5,3.5) node[midway, above, font=\tiny, poporange] {bù $\rightarrow$ bú};
\node[label] at (10, 0) {Sandhi \zh{不}};

% Sandhi 一 (yī) -> yí devant 4e ton (xià)
\draw[tone, popgreen, dashed] (11.5,1) -- (12.5,3.5) node[midway, above, font=\tiny, popgreen] {yī $\rightarrow$ yí};
\node[label] at (12, 0) {Sandhi \zh{一}};

% 3e ton + 3e ton -> 2e + 3e
\draw[tone, poppurple, dashed] (13.5,2) -- (14.5,3.5) node[midway, above, font=\tiny, poppurple] {3+3 $\rightarrow$ 2+3};
\node[label] at (14, 0) {Sandhi 3+3};
\end{tikzpicture}
\legende{Repères visuels des tons et sandhi principaux (flèches pleines = tons canoniques, flèches pointillées = transformations obligatoires).}
\end{popfigure}

\courssub{Jeux de rôle et débat guidé : Mise en situation}

\begin{methode}[Organisation des jeux de rôle (Binômes / Groupes de 3)]
\begin{enumerate}
    \item \textbf{Distribution des cartes-rôles} (5 min prep) : Chaque élève reçoit une fiche avec son profil (revenus, personnalité, contrainte).
    \item \textbf{Jeu} (3-4 min) : Interaction spontanée. L'enseignant circule, note les structures cibles (\zh{把}, \zh{因为/所以}, connecteurs).
    \item \textbf{Débriefing collectif} (5 min) : 1 binôme rejoue devant la classe. Correction bienveillante (tons, structure, vocabulaire).
\end{enumerate}
\end{methode}

\begin{experience}[Jeu de rôle 1 : « L'ami dépensier » (Binôme)]
\begin{itemize}
    \item \textbf{Rôle A (Conseiller) :} Tu es économe. Tu as 2000元/月. Tu veux aider ton ami à tenir son budget. Utilise : \zh{你应该...}, \zh{建议你...}, \zh{把钱花在...上}.
    \item \textbf{Rôle B (Dépensier) :} Tu as 2000元/月 mais tu es toujours à sec. Tu aimes : jeux vidéo (\zh{游戏}), snacks (\zh{零食}), vêtements (\zh{衣服}), sortir (\zh{出去玩}). Tu te plains : \zh{钱不够用}, \zh{我想买...}.
    \item \textbf{Objectif A :} Obtenir un engagement concret de B (\zh{我打算下个月少买...}).
    \item \textbf{Objectif B :} Justifier tes dépenses avec \zh{因为...所以...}.
\end{itemize}
\end{experience}

\begin{experience}[Jeu de rôle 2 : « Négociation parents/ados » (Binôme)]
\begin{itemize}
    \item \textbf{Rôle A (Parent) :} Revenu fixe. Exige : \zh{记账}, \zh{收支平衡}, \zh{大额支出商量}. Refuse le "trou" mensuel. Formules : \zh{下不为例}, \zh{你得学会...}.
    \item \textbf{Rôle B (Lycéen) :} Besoin d'argent pour : \zh{买参考书} (livres de révision), \zh{报补习班} (cours de soutien), \zh{买球鞋} (baskets). Doit convaincre avec un budget prévisionnel : \zh{我计划...}, \zh{这个月特殊因为...}.
    \item \textbf{Contrainte linguistique :} Minimum 3 \zh{因为...所以...} et 2 \zh{把...花在...上} dans l'échange.
\end{itemize}
\end{experience}

\begin{experience}[Débat guidé : « Les lycéens devraient-ils faire un job étudiant ? » (Groupes de 4 : 2 Pour, 2 Contre + 1 Modérateur)]
\end{experience}

\begin{methode}[Structure du débat (15 min total)]
\begin{enumerate}
    \item \textbf{Préparation arguments (3 min) :} Chaque camp note 3 arguments en mots-clés chinois sur fiche.
    \item \textbf{Tour 1 - Affirmation (4 min) :} Pour 1 argumente ($\approx$ 45s), Contre 1 répond/contre-argumente ($\approx$ 45s). Alternance.
    \item \textbf{Tour 2 - Approfondissement (4 min) :} Pour 2 / Contre 2. Nouveaux arguments ou développement.
    \item \textbf{Synthèse modérateur (2 min) :} \zh{总的来说，赞成方认为... 反方认为... 我个人觉得...} (Globalement, les pour pensent... les contre pensent... moi je pense...).
    \item \textbf{Vote classe} : Main levée \zh{赞成} / \zh{反对} / \zh{弃权}.
\end{enumerate}
\end{methode}

\begin{aretenir}[Vocabulaire débat : \zh{兼职} (Job étudiant)]
\begin{center}
\begin{tabularx}{\linewidth}{|X|X|X|X|}
\hline
\rowcolor{poppurpleL} \textbf{Caractères} & \textbf{Pinyin} & \textbf{Nature} & \textbf{Traduction} \\
\hline
\zh{赞成} & zànchéng & verbe & être pour, approuver \\
\hline
\zh{反对} & fǎnduì & verbe & être contre, s'opposer \\
\hline
\zh{锻炼} & duànliàn & verbe & exercer, se forger (caractère, capacité) \\
\hline
\zh{社会经验} & shèhuì jīngyàn & nom & expérience sociale \\
\hline
\zh{影响学习} & yǐngxiǎng xuéxí & v+o & affecter les études \\
\hline
\zh{时间管理} & shíjiān guǎnlǐ & nom & gestion du temps \\
\hline
\zh{经济独立} & jīngjì dúlì & adj/nom & indépendance financière \\
\hline
\zh{安全隐患} & ānquán yǐnhuàn & nom & risque pour la sécurité \\
\hline
\end{tabularx}
\end{center}
\end{aretenir}

\courssub{Production orale guidée et évaluation}

Après avoir pratiqué les dialogues types et les jeux de rôle de la section précédente, vous êtes maintenant prêts à passer à l'exercice de synthèse : la production orale autonome. Cette étape finale vise à vérifier votre capacité à mobiliser le vocabulaire budgétaire, les structures grammaticales (notamment les connecteurs logiques et la structure \zh{把}) et les formules fonctionnelles dans un exposé structuré, puis à interagir spontanément avec l'auditoire.

\courssub{La tâche finale : « Mon budget et moi »}

La consigne officielle impose un exposé court (8 à 10 phrases) sur votre argent de poche (\zh{零花钱} \emph{línghuāqián}). Le plan est \textbf{imposé} en quatre étapes logiques, chacune introduite par un connecteur obligatoire. Cela permet à l'examinateur de suivre votre raisonnement sans effort.

\begin{definition}[Connecteurs logiques obligatoires]
Les quatre mots-outils suivants structurent obligatoirement votre exposé. Ils placent les événements ou les idées dans une séquence temporelle ou logique.
\begin{center}
\begin{tabularx}{\linewidth}{|X|X|X|X|}
\hline
\rowcolor{popblueL} \textbf{Connecteur} & \textbf{Pinyin} & \textbf{Fonction} & \textbf{Équivalent français} \\
\hline
\zh{先} & xiān & Étape 1 : Introduction du sujet, montant perçu & « D'abord », « Premièrement » \\
\hline
\zh{然后} & ránhòu & Étape 2 : Énumération des dépenses & « Ensuite », « Puis » \\
\hline
\zh{但是} & dànshì & Étape 3 : Difficulté, contraste, imprévu & « Mais », « Cependant » \\
\hline
\zh{所以} & suǒyǐ & Étape 4 : Conclusion, avis, conséquence & « Donc », « C'est pourquoi » \\
\hline
\end{tabularx}
\end{center}
\end{definition}

\begin{propriete}[Position des connecteurs en début d'énoncé]
En chinois, ces connecteurs se placent \textbf{au tout début de la phrase} (sujet + verbe...), souvent suivis d'une virgule (，) à l'oral marquée par une pause.
\emph{Schéma :} \textbf{Connecteur + Sujet + Verbe + Compléments.}
\emph{Exemple correct :} \zh{先，我介绍一下我的零花钱。} (D'abord, je présente mon argent de poche.)
\emph{Exemple correct :} \zh{然后，我把钱花在交通上。}
\textbf{Attention :} Ne pas dire \zh{我先有...} (sauf si « 先 » modifie le verbe « faire d'abord » comme dans \zh{我先吃饭} « Je mange d'abord »). Pour structurer le discours, le connecteur est en position initiale.
\end{propriete}

\courssub{Modèle d'exposé et analyse}

Voici un modèle conforme aux attendus (8 phrases). Observez comment chaque paragraphe correspond à une étape du plan. Notez l'absence de répétition du connecteur \zh{先} en début de phrase 2.

\begin{textebox}[Modèle d'exposé : « Mon budget et moi »]
\zh{先，我介绍一下我的零花钱。} \\
\zh{我每个月有两百块钱。} \\
\zh{然后，我把钱花在交通和吃饭上。} \\
\zh{有时候，我还买书和文具。} \\
\zh{但是，有时候钱不够用。} \\
\zh{所以，我觉得学生应该学会省钱。} \\
\zh{我打算下个月少买零食。} \\
\zh{我的介绍完了，谢谢大家！}
\end{textebox}

\begin{center}
\begin{tabularx}{\linewidth}{|X|X|X|}
\hline
\rowcolor{popblueL} \textbf{Phrase (Caractères)} & \textbf{Pinyin} & \textbf{Traduction} \\
\hline
\zh{先，我介绍一下我的零花钱。} & Xiān, wǒ jièshào yíxià wǒ de línghuāqián. & D'abord, je vais présenter brièvement mon argent de poche. (\zh{yíxià} sandhi) \\
\hline
\zh{我每个月有两百块钱。} & Wǒ měi gè yuè yǒu liǎng bǎi kuài qián. & J'ai 200 yuans chaque mois. (\zh{liǎng} pour 200) \\
\hline
\zh{然后，我把钱花在交通和吃饭上。} & Ránhòu, wǒ bǎ qián huā zài jiāotōng hé chīfàn shàng. & Ensuite, je dépense l'argent dans les transports et la nourriture. \\
\hline
\zh{有时候，我还买书和文具。} & Yǒushíhou, wǒ hái mǎi shū hé wénjù. & Parfois, j'achète aussi des livres et de la papeterie. \\
\hline
\zh{但是，有时候钱不够用。} & Dànshì, yǒushíhou qián \textbf{bú} gòu yòng. & Mais parfois l'argent ne suffit pas. (\zh{búgòu} sandhi) \\
\hline
\zh{所以，我觉得学生应该学会省钱。} & Suǒyǐ, wǒ juéde xuéshēng yīnggāi xuéhuì \textbf{shéng}qián. & Donc, je pense que les élèves devraient apprendre à économiser. (\zh{shéngqián} sandhi) \\
\hline
\zh{我打算下个月少买零食。} & Wǒ \textbf{dá}suàn xià gè yuè shǎo mǎi língshí. & Je prévois d'acheter moins de snacks le mois prochain. (\zh{dá}suàn sandhi) \\
\hline
\zh{我的介绍完了，谢谢大家！} & Wǒ de jièshào wánle, xièxie dàjiā! & Ma présentation est finie, merci à tous ! \\
\hline
\end{tabularx}
\end{center}

\begin{aretenir}[Vocabulaire clé de l'exposé + Points de vigilance]
\begin{center}
\begin{tabularx}{\linewidth}{|X|X|X|X|}
\hline
\rowcolor{popblueL} Caractères & Pinyin & Nature & Traduction / Note \\
\hline
\zh{零花钱} & línghuāqián & nom & argent de poche \\
\hline
\zh{块 (钱)} & kuài (qián) & classif./nom & yuan (oral) / argent \\
\hline
\zh{把...花在...上} & bǎ... huā zài... shàng & structure & dépenser (qqch) pour (qqch) — \textbf{Objet défini requis} \\
\hline
\zh{交通} & jiāotōng & nom & transports \\
\hline
\zh{文具} & wénjù & nom & papeterie, fournitures scolaires \\
\hline
\zh{够用} & gòu yòng & verbe & suffire, être assez \\
\hline
\zh{省钱} & shěngqián (\textbf{shéngqián}) & verbe & économiser de l'argent — \textbf{Sandhi 3+3 !} \\
\hline
\zh{打算} & dǎsuàn (\textbf{dásuàn}) & verbe & prévoir, avoir l'intention de — \textbf{Sandhi 3+3 !} \\
\hline
\zh{零食} & língshí & nom & snacks, grignotines \\
\hline
\zh{介绍完了} & jièshào wánle & V+Compl & avoir fini de présenter (\zh{完} = achèvement) \\
\hline
\end{tabularx}
\end{center}
\end{aretenir}

\courssub{Focus grammatical : La structure \zh{把} (bǎ) pour les dépenses}

Le modèle utilise \zh{我把钱花在交通和吃饭上}. C'est la structure standard pour dire « dépenser [objet défini] dans [domaine] ».

\begin{propriete}[Structure \zh{把} avec \zh{花} (dépenser) — Règle de l'objet défini]
\emph{Sujet + 把 + Objet (défini/connu) + 花在 + Lieu/Domaine + 上.}
\begin{itemize}
    \item \textbf{Condition :} L'objet (\zh{钱}, \zh{零花钱}, \zh{工资}) doit être \textbf{connu, spécifique, ou quantifié}. On ne dit pas \zh{我把钱花...} si on parle d'argent en général indéfini (on utiliserait \zh{花钱在...上} sans \zh{把}).
    \item \zh{在...上} indique la destination de la dépense (domaine, activité, objet d'achat).
    \item \zh{了} en fin de phrase = action terminée / situation nouvelle (ex: \zh{钱花完了}).
\end{itemize}
\end{propriete}

\begin{exemplebox}[Variantes et pièges]
\zh{我把零花钱都花在吃饭上了。} (Wǒ bǎ línghuāqián dōu huā zài chīfàn shàng \textbf{le}.) — J'ai dépensé \textbf{tout} mon argent de poche en nourriture. (\zh{都} + \zh{le} = totalité terminée).
\zh{他把工资花在房租上。} (Tā bǎ gōngzī huā zài fángzū shàng.) — Il dépense son salaire pour le loyer.
\textbf{Erreur fréquente :} \zh{我把钱花在买书。} (Incorrect : \zh{买书} est une action, pas un lieu/domaine). \rightarrow \textbf{Correct :} \zh{我把钱花在书上。} ou \zh{我把钱用来买书。} (\zh{用来} = utiliser pour).
\end{exemplebox}


\courssub{Préparation écrite (5 minutes) et déroulement oral}

L'examen ne teste pas votre mémoire, mais votre capacité à \textbf{structurer} et \textbf{parler}. Vous avez 5 minutes pour rédiger un \textbf{plan uniquement} (mots-clés, chiffres, connecteurs), \textbf{pas un texte complet}.

\begin{methode}[Déroulement de l'épreuve orale (10 min par binôme + classe)]
\begin{enumerate}
    \item \textbf{Préparation individuelle (5 min) :} Remplissez une fiche « Plan » (voir tableau ci-dessous). Interdit d'écrire des phrases complètes en caractères. Pinyin ou mots-clés seulement.
    \item \textbf{Passage en binôme (3-4 min) :} L'élève A présente son exposé (1 min 30). L'élève B écoute activement et pose \textbf{2 questions} (\zh{你呢？} / \zh{为什么？} / \zh{怎么省钱？}). Inversion des rôles.
    \item \textbf{Devant la classe (volontaires ou tirage au sort) :} Exposé face à la classe (1 min 30). Questions de 2 camarades (interaction notée sur 2 pts).
\end{enumerate}
\end{methode}

\begin{center}
\begin{tabularx}{\linewidth}{|X|X|X|}
\hline
\rowcolor{popgreenL} \textbf{Étape du plan} & \textbf{Connecteur} & \textbf{Mots-clés à noter (Exemple élève)} \\
\hline
1. Montant perçu / Source & \zh{先} & \zh{每月} \zh{300块} \zh{父母给} / \zh{兼职赚} \\
\hline
2. Dépenses principales (3 postes) & \zh{然后} & \zh{把钱花在} \zh{交通} \zh{、} \zh{吃饭} \zh{、} \zh{文具} \zh{上} \\
\hline
3. Difficulté / Épargne actuelle & \zh{但是} & \zh{钱不够用} / \zh{存} \zh{50块} / \zh{省钱} \\
\hline
4. Avis / Résolution future & \zh{所以} & \zh{觉得 学生 应该 省钱} / \zh{打算 下月 少买 零食} \\
\hline
Formules ouverture / fin & — & \zh{先介绍一下...} / \zh{介绍完了，谢谢大家} \\
\hline
\end{tabularx}
\end{center}

\courssub{Gestion du stress et attitude professionnelle}

Parler en public dans une langue étrangère est anxiogène. La grille d'évaluation valorise la fluidité et l'interaction, qui s'effondrent si vous lisez.

\begin{methode}[Gérer son stress à l'oral : 4 réflexes]
\begin{enumerate}
    \item \textbf{Respirer :} Inspirez profondément par le nez (4 temps), expirez par la bouche (6 temps) avant de commencer. Cela calme le rythme cardiaque et pose la voix.
    \item \textbf{Parler lentement :} L'examinateur préfère une phrase lente et tonale qu'une phrase rapide et incompréhensible. Articulez les tons (surtout le 3e ton bas \zh{shěng}, \zh{dǎ} et le 4e ton descendant \zh{kuài}, \zh{gòu}, \zh{suǒyǐ}).
    \item \textbf{Regarder l'auditoire :} Ne fixez pas votre fiche. Jetez un œil sur le mot-clé \rightarrow levez la tête \rightarrow dites la phrase. Contact visuel = points d'interaction.
    \item \textbf{Notes-plan uniquement :} Votre fiche ne contient que \zh{先 300块} \zh{然后 交通 吃饭} \zh{但是 不够} \zh{所以 省钱}. Si vous séchez, dites : \zh{等一下} (Attendez un peu) ou \zh{怎么说？} (Comment dire ?) — cela montre une stratégie de communication, pas un échec.
\end{enumerate}
\end{methode}

\begin{attention}[Piège fatal : Lire son texte]
\textbf{Un exposé lu = 0 point en Fluidité et 0 point en Interaction (soit -6 pts sur 20).}
\begin{itemize}
    \item La lecture se détecte : regard fixe sur la feuille, débit monotone, absence de pauses naturelles, pas de contact visuel.
    \item Si vous avez écrit des phrases complètes (brouillon), \textbf{ne les lisez pas}. Utilisez-les comme filet de sécurité mental, mais parlez à partir des mots-clés.
    \item \emph{Astuce :} Écrivez votre plan au dos de la convocation ou sur un petit bristol (format carte de visite) : cela vous force à être synthétique et à lever la tête.
\end{itemize}
\end{attention}

\courssub{Grille d'évaluation détaillée (20 points)}

Cette grille est alignée sur les critères du Baccalauréat (LV2). Chaque descripteur est observable.

\begin{center}
\begin{tabularx}{\linewidth}{|X|X|X|}
\hline
\rowcolor{poporangeL} \textbf{Critère} & \textbf{Bareme} & \textbf{Indicateurs de réussite (Niveau A / B)} \\
\hline
\textbf{Contenu \& Structure} (4 pts) & 0 -- 4 & \textbf{4 :} Plan respecté (4 étapes), 8-10 phrases, infos précises (montant, 3 postes). \textbf{2 :} Plan incomplet, phrases < 6, infos vagues. \textbf{0 :} Hors sujet. \\
\hline
\textbf{Formules fonctionnelles} (4 pts) & 0 -- 4 & \textbf{4 :} 4 connecteurs utilisés correctement + formules d'ouverture/fin (\zh{介绍一下}, \zh{完了谢谢}) + vocabulaire budget précis (\zh{把...花在...上}, \zh{省钱}, \zh{打算}). \textbf{2 :} 2-3 connecteurs, vocabulaire pauvre. \\
\hline
\textbf{Tons \& Prononciation} (6 pts) & 0 -- 6 & \textbf{6 :} Tons justes sur mots-clés (\zh{liǎng, bǎi, huā, gòu, shěng, dǎsuàn}), distinction j/q/x, zh/ch/sh, en/eng, an/ang, sandhi \zh{不/一/3+3} respectés. \textbf{4 :} Erreurs sur tons isolés mais compréhensible. \textbf{2 :} Tons aléatoires, incompréhensible par moments. \\
\hline
\textbf{Fluidité} (4 pts) & 0 -- 4 & \textbf{4 :} Débit naturel, pauses logiques (après connecteurs), auto-corrections rares, pas de lecture. \textbf{2 :} Hésitations longues, phrases hachées, lecture partielle. \textbf{0 :} Lecture totale, silences > 10s. \\
\hline
\textbf{Interaction} (2 pts) & 0 -- 2 & \textbf{2 :} Répond aux questions (\zh{你呢？}, \zh{为什么？}) par une phrase complète (\zh{因为...所以...}), regarde l'interlocuteur. \textbf{1 :} Réponse monosyllabique (\zh{对}, \zh{是}) ou en français. \textbf{0 :} Ne répond pas, ne pose pas de question en binôme. \\
\hline
\end{tabularx}
\end{center}

\courssub{Visualisation : La frise logique de l'exposé}

Cette figure résume l'enchaînement obligatoire des idées. Mémorisez cette flèche : \textbf{Montant $\rightarrow$ Dépenses $\rightarrow$ Problème $\rightarrow$ Solution/Avis}.

\begin{popfigure}
\begin{tikzpicture}[
    node distance=1.2cm and 0.5cm,
    box/.style={rectangle, rounded corners, draw=popblue, thick, fill=popblueL, text width=2.8cm, align=center, minimum height=1.8cm, font=\small},
    arrow/.style={-Stealth, thick, draw=poporange, shorten >=2pt, shorten <=2pt},
    labelbox/.style={rectangle, draw=poporange, fill=poporangeL, rounded corners, font=\tiny\bfseries, text width=1.5cm, align=center, minimum height=0.7cm}
]
% Nœuds principaux
\node[box, align=center] (etape1){\textbf{ÉTAPE 1}\\\zh{先} (Xiān)\\\emph{D'abord}\\Montant perçu\\\zh{我每个月有...块}};
\node[box, right=of etape1, align=center] (etape2){\textbf{ÉTAPE 2}\\\zh{然后} (Ránhòu)\\\emph{Ensuite}\\Dépenses\\\zh{我把钱花在...上}};
\node[box, right=of etape2, align=center] (etape3){\textbf{ÉTAPE 3}\\\zh{但是} (Dànshì)\\\emph{Mais}\\Difficulté / Épargne\\\zh{钱不够用 / 存钱}};
\node[box, right=of etape3, align=center] (etape4){\textbf{ÉTAPE 4}\\\zh{所以} (Suǒyǐ)\\\emph{Donc}\\Avis / Résolution\\\zh{我觉得...应该...}};

% Flèches principales
\draw[arrow] (etape1.east) -- (etape2.west) node[midway, above, labelbox] {Séquence};
\draw[arrow] (etape2.east) -- (etape3.west) node[midway, above, labelbox] {Énumération};
\draw[arrow] (etape3.east) -- (etape4.west) node[midway, above, labelbox] {Contraste $\rightarrow$ Conséquence};

% Étiquettes de connecteurs (grandes)
\node[above=0.2cm of etape1.north, font=\large\bfseries, text=popdark] {\zh{先}};
\node[above=0.2cm of etape2.north, font=\large\bfseries, text=popdark] {\zh{然后}};
\node[above=0.2cm of etape3.north, font=\large\bfseries, text=popdark] {\zh{但是}};
\node[above=0.2cm of etape4.north, font=\large\bfseries, text=popdark] {\zh{所以}};
\end{tikzpicture}
\legende{Frise logique de l'exposé « Mon budget et moi » : les 4 connecteurs obligatoires structurent le récit du constat à la résolution.}
\end{popfigure}

\courssub{Exemples supplémentaires pour l'interaction (Questions/Réponses)}

En binôme ou face à la classe, vous devez comprendre et produire ces échanges types.

\begin{exemplebox}[Échanges types : Questions des camarades]
\begin{itemize}
    \item \textbf{Q :} \zh{你每个月有多少零花钱？} (Nǐ měi gè yuè yǒu duōshao línghuāqián ?) — Combien as-tu par mois ?
    \item \textbf{R :} \zh{我有三百块。} (Wǒ yǒu sān bǎi kuài.) — J'ai 300 yuans.
    \item \textbf{Q :} \zh{你把钱花在什么地方？} (Nǐ bǎ qián huā zài shénme dìfang ?) — Où dépenses-tu ton argent ?
    \item \textbf{R :} \zh{主要花在吃饭和交通上。} (Zhǔyào huā zài chīfàn hé jiāotōng shàng.) — Principalement nourriture et transports.
    \item \textbf{Q :} \zh{为什么有时候钱不够用？} (Wèishéme yǒushíhou qián bù gòu yòng ?) — Pourquoi parfois pas assez ?
    \item \textbf{R :} \zh{因为我想买新衣服，所以钱不够。} (Yīnwèi wǒ xiǎng mǎi xīn yīfu, suǒyǐ qián bù gòu.) — Parce que je veux acheter des vêtements neufs, donc pas assez.
    \item \textbf{Q :} \zh{你怎么省钱？} (Nǐ zěnme shěngqián ?) — Comment économises-tu ?
    \item \textbf{R :} \zh{我少买零食，走路不坐车。} (Wǒ shǎo mǎi língshí, zǒulù bù zuò chē.) — J'achète moins de snacks, je marche au lieu de prendre le bus.
\end{itemize}
\end{exemplebox}

\begin{propriete}[Structure \zh{因为}... \zh{所以}... pour justifier (Réponse à \zh{为什么})]
C'est la paire corrélative standard pour « Parce que... donc... ». Les deux conjoints sont \textbf{obligatoires} en chinois standard (contrairement au français où « donc » est souvent omis).
\emph{Schéma :} \zh{因为} + Cause + \zh{，所以} + Conséquence.
\emph{Exemple :} \zh{因为交通费贵，所以我打算走路。} (Yīnwèi jiāotōngfèi guì, suǒyǐ wǒ dǎsuàn zǒulù.) — Parce que le transport est cher, je prévois de marcher.
\end{propriete}

\courssub{Exercice résolu : Simulation d'évaluation}

\begin{exoresolu}[Évaluation d'un exposé élève]
\textbf{Énoncé :} Écoutez (ou lisez) l'exposé simulé d'un élève de Terminale. Notez-le sur 20 avec la grille ci-dessus et justifiez chaque note.

\emph{Exposé de l'élève « Paul » :}
\zh{先，我有零花钱。每个月两百块。然后，我花钱吃饭、坐车。但是，钱不够。所以，我要省钱。下个月我不买零食。完了，谢谢。}

\tcblower
\textbf{Solution détaillée :}

\begin{center}
\begin{tabularx}{\linewidth}{|X|X|X|}
\hline
\rowcolor{popblueL} \textbf{Critère} & \textbf{Note / Max} & \textbf{Justification} \\
\hline
Contenu \& Structure & \textbf{3 / 4} & Plan respecté (4 étapes), mais phrases très courtes (7 phrases seulement, limite basse). Manque de détails (pas de structure \zh{把}, pas de \zh{文具} etc.). Phrase 1 incomplète (\zh{我有零花钱} sans quantificateur précis immédiat). \\
\hline
Formules fonctionnelles & \textbf{3 / 4} & 4 connecteurs présents et bien placés. Formule de fin (\zh{完了，谢谢}) correcte. Manque formule d'ouverture polie (\zh{介绍一下}). Vocabulaire minimaliste (\zh{花钱} au lieu de \zh{把钱花在...上}, \zh{坐车} au lieu de \zh{交通}). \\
\hline
Tons \& Prononciation & \textbf{4 / 6} & (Hypothèse orale) Si tons corrects sur \zh{liǎng, bǎi, huā, gòu, shěng} $\rightarrow$ 5-6. Si erreurs sur \zh{xiān/xiàn, ránhòu, dànshì, suǒyǐ} ou sandhi \zh{búgòu, shéngqián} $\rightarrow$ 3-4. Notons 4 (standard « compréhensible avec erreurs »). \\
\hline
Fluidité & \textbf{3 / 4} & Débit haché (phrases nominales \zh{每个月两百块} sans verbe \zh{有}). Pas de lecture détectée. Pauses correctes après connecteurs. \\
\hline
Interaction & \textbf{1 / 2} & (Hypothèse) Répond \zh{对} ou \zh{是} aux questions \zh{为什么？} au lieu de phrase complète avec \zh{因为}...\zh{所以}... \\
\hline
\textbf{TOTAL} & \textbf{14 / 20} & \textbf{Niveau : Assez Bien.} Progression possible : allonger les phrases (SVO complet), utiliser structure \zh{把}, développer l'interaction par phrases complètes. \\
\hline
\end{tabularx}
\end{center}
\end{exoresolu}

\courssub{Exercice final : À vous de jouer !}

\begin{exercice}[Préparez votre exposé noté]
Vous avez \textbf{5 minutes} pour remplir votre fiche de préparation (recto seulement). Utilisez \textbf{vos vrais chiffres} (argent de poche réel au Cameroun : 5 000 -- 20 000 FCFA $\approx$ 50 -- 200 CNY) ou des chiffres inventés plausibles.

\begin{center}
\begin{tabularx}{\linewidth}{|X|X|}
\hline
\rowcolor{popgreenL} \textbf{Votre Fiche Plan (Recto)} & \textbf{Mots-clés (Chinois + Chiffres)} \\
\hline
\zh{先} : Montant / Source & \zh{每月 \_\_\_\_ 块 / 父母给 / 兼职} \\
\hline
\zh{然后} : 3 postes de dépenses & \zh{把钱花在 \_\_\_\_ 、 \_\_\_\_ 、 \_\_\_\_ 上} \\
\hline
\zh{但是} : Problème / Épargne & \zh{钱不够 / 存 \_\_\_\_ 块 / 省钱} \\
\hline
\zh{所以} : Avis / Action future & \zh{觉得 学生 应该 \_\_\_\_ / 打算 下月 \_\_\_\_} \\
\hline
Formules ouverture / fin & \zh{先介绍一下... / 介绍完了，谢谢大家} \\
\hline
\end{tabularx}
\end{center}

\textbf{Consignes :}
\begin{enumerate}
    \item Remplissez les tirets en caractères ou pinyin.
    \item Entraînez-vous à voix basse 2 minutes (chronomètre : 1 min 30 visé).
    \item Présentez à votre binôme. Lui note les connecteurs entendus (coche sur fiche).
    \item Échangez les rôles.
    \item 2 binômes passent devant la classe pour évaluation formelle (grille 20 pts).
\end{enumerate}
\end{exercice}

\begin{corrige}[Exercice : Préparation de l'exposé]
Pas de correction unique (production personnelle). \textbf{Grille d'auto-évaluation avant passage :}
\begin{itemize}
    \item [ ] Mes 4 connecteurs (\zh{先, 然后, 但是, 所以}) sont écrits sur ma fiche.
    \item [ ] J'ai au moins 1 chiffre précis (montant) et 3 noms de dépenses.
    \item [ ] J'utilise la structure \zh{把钱花在...上} à l'étape 2.
    \item [ ] J'ai une phrase avec \zh{觉得...应该...} ou \zh{打算...} à l'étape 4.
    \item [ ] Je connais la réponse type à \zh{为什么？} (\zh{因为...所以...}).
    \item [ ] Mon exposé dure entre 1 min 15 et 1 min 45 à vitesse normale.
\end{itemize}
Si toutes les cases sont cochées $\rightarrow$ Vous visez 16+/20. Bonne chance !
\end{corrige}

\begin{savaistu}[Argent de poche : Cameroun vs Chine]
Au Cameroun, l'argent de poche (\zh{零花钱}) est souvent hebdomadaire ou mensuel, donné par les parents pour le transport (taxi/moto), le \zh{午餐} (déjeuner) et les photocopies. En Chine, les lycéens reçoivent souvent un \zh{生活费} (subsides de vie) mensuel plus conséquent (500--1500 CNY) car ils logent parfois à l'internat (\zh{住校}) et gèrent eux-mêmes cantine, lessive, fournitures. Le paiement mobile (\zh{微信支付}, \zh{支付宝}) est roi en Chine (même pour 1 yuan) ; au Cameroun, le cash (FCFA) et \zh{Mobile Money} (Orange Money, MTN MoMo) cohabitent. Dire \zh{我用手机付钱} (Je paie avec mon téléphone) est universel aujourd'hui !
\end{savaistu}

\newpage

\coursec{Activité d'intégration}

\courssub{Objectifs de l'activité}

À la fin de cette activité d'intégration, l'élève sera capable de :
\begin{itemize}
\item prendre la parole en chinois dans un débat organisé sur un sujet de la vie quotidienne : la gestion du budget et l'argent de poche ;
\item mobiliser le vocabulaire du budget (\zh{零花钱}, \zh{钱}, \zh{花}, \zh{存}, \zh{买}, \zh{需要}…) dans des échanges réels ;
\item utiliser les formules fonctionnelles de la leçon : donner son avis (\zh{我觉得…}), marquer l'accord (\zh{我同意}), le désaccord poli (\zh{我不太同意}), relancer la parole (\zh{你呢？你觉得怎么样？}) et conclure (\zh{所以，我们认为…}) ;
\item construire une argumentation simple avec la structure \zh{因为…所以…} ;
\item prononcer correctement les tons, en particulier dans les mots-clés du débat (\zh{零花钱}, \zh{应该}, \zh{同意}, \zh{觉得}, \zh{存钱}) ;
\item écouter les autres, réagir à leurs propos et voter pour l'équipe la plus convaincante.
\end{itemize}

\courssub{Situation de communication}

C'est la semaine du club chinois du lycée de Yaoundé. Le thème de l'année est « \zh{钱和生活} » (\emph{Qián hé shēnghuó}, « l'argent et la vie »). Pour la séance de clôture, la présidente du club a proposé un grand débat en chinois, devant toute la classe et le professeur. Le sujet choisi concerne directement les élèves : à Yaoundé comme à Douala ou à Buea, certains lycéens reçoivent de l'argent de poche pour le transport, la cantine ou le téléphone ; d'autres non. Faut-il en donner ? La question est posée en chinois sur l'affiche du club :

\begin{textebox}[Affiche du club chinois — Lycée de Yaoundé]
\textbf{大辩论：学生应该有零花钱吗？}\\
\emph{Dà biànlùn: Xuésheng yīnggāi yǒu língqiánhuā ma?}\\
« Grand débat : les lycéens devraient-ils recevoir de l'argent de poche ? »\\[4pt]
甲方：学生应该有零花钱。\\
\emph{Jiǎfāng: Xuésheng yīnggāi yǒu língqiánhuā.}\\
« Équipe A : les lycéens doivent recevoir de l'argent de poche. »\\
乙方：学生不应该有零花钱。\\
\emph{Yǐfāng: Xuésheng bù yīnggāi yǒu língqiánhuā.}\\
« Équipe B : les lycéens ne doivent pas recevoir d'argent de poche. »\\[4pt]
时间：星期五下午三点。地点：中文教室。欢迎大家来听！\\
\emph{Shíjiān: Xīngqīwǔ xiàwǔ sān diǎn. Dìdiǎn: Zhōngwén jiàoshì. Huānyíng dàjiā lái tīng!}\\
« Date : vendredi à 15 h. Lieu : salle de chinois. Tout le monde est bienvenu ! »
\end{textebox}

Le professeur divise la classe : une équipe défend la position \zh{甲方}, l'autre la position \zh{乙方}. Chaque équipe prépare trois arguments, désigne un animateur, des orateurs et un rapporteur. À la fin, la classe vote.

\courssub{Consignes détaillées}

\textbf{Tâche 1 — Comprendre le sujet (10 min).} Lis l'affiche du club à voix haute avec ton groupe. Souligne les mots du budget que tu connais déjà (\zh{零花钱}, \zh{应该}). Vérifie dans le tableau de vocabulaire le sens et les tons de \zh{辩论} (biànlùn, 4\up{e} + 4\up{e} ton) et de \zh{甲方} / \zh{乙方}. Puis, en français, reformule la question du débat en une phrase pour être sûr de l'avoir comprise.

\textbf{Tâche 2 — Préparer les arguments (20 min).} Dans ton équipe, rédige par écrit \textbf{trois arguments} en chinois, chacun avec la structure \zh{因为…所以…}. Chaque argument doit contenir au moins un mot du lexique du budget. Exemple de départ : \zh{因为学生要坐车去学校，所以他们需要零花钱。} Répartissez les rôles : un animateur, trois orateurs (un argument chacun), un rapporteur.

\textbf{Tâche 3 — Travailler les tons (10 min).} Répétez à voix haute les mots-clés du débat en insistant sur les tons : \zh{零花钱} (língqiánhuā : 2-2-1), \zh{应该} (yīnggāi : 1-1), \zh{同意} (tóngyì : 2-4), \zh{觉得} (juéde : 2 - ton neutre), \zh{存钱} (cúnqián : 2-2). Le professeur écoute chaque équipe et corrige.

\textbf{Tâche 4 — Le débat (25 min).} L'animateur de l'équipe A ouvre le débat avec \zh{我觉得学生应该有零花钱。} Chaque orateur présente son argument. L'animateur relance l'autre équipe avec \zh{你呢？你觉得怎么样？} L'équipe B répond avec \zh{我不太同意。因为…所以…} Chaque équipe a droit à une réplique par argument.

\textbf{Tâche 5 — Conclure et voter (10 min).} Le rapporteur de chaque équipe conclut avec \zh{所以，我们认为…}. La classe vote à main levée pour l'équipe la plus convaincante, puis le professeur annonce les résultats de la grille d'évaluation.

\courssub{Vocabulaire du budget}

\begin{definition}[Lexique essentiel — Gestion du budget]
Le tableau ci-dessous regroupe les mots et expressions indispensables pour parler d'argent de poche, de dépenses et d'épargne. Mémorise les tons (marques sur les voyelles) et la nature grammaticale.
\end{definition}

\begin{center}
\begin{tabularx}{\linewidth}{|X|X|X|X|}
\hline
\rowcolor{popblueL} Caractères & Pinyin & Nature & Traduction \\
\hline
零花钱 & língqiánhuā & nom & argent de poche \\
\hline
钱 & qián & nom & argent, monnaie \\
\hline
花 & huā & verbe & dépenser (argent, temps) \\
\hline
存钱 & cúnqián & locution verbale & économiser de l'argent, mettre de côté \\
\hline
买 & mǎi & verbe & acheter \\
\hline
需要 & xūyào & verbe & avoir besoin de \\
\hline
应该 & yīnggāi & verbe modal & devoir, il faut (obligation morale / conseil) \\
\hline
辩论 & biànlùn & nom / verbe & débat, débattre \\
\hline
甲方 / 乙方 & jiǎfāng / yǐfāng & nom & partie A (pour) / partie B (contre) \\
\hline
同意 & tóngyì & verbe & être d'accord \\
\hline
觉得 & juéde & verbe & penser, trouver (avis subjectif) \\
\hline
有用 & yǒuyòng & adjectif & utile \\
\hline
没用 & méiyòng & adjectif & inutile \\
\hline
浪费 & làngfèi & verbe / adjectif & gaspiller / gaspillage \\
\hline
太 & tài & adverbe & trop (excès) \\
\hline
但是 & dànshì & conjonction & mais, cependant \\
\hline
所以 & suǒyǐ & conjonction & donc, par conséquent \\
\hline
因为 & yīnwèi & conjonction & parce que \\
\hline
\end{tabularx}
\end{center}

\begin{aretenir}[Points de vigilance lexicaux]
\begin{itemize}
\item \zh{零花钱} : composé de \zh{零} (éparpillé, menu), \zh{花} (dépenser) et \zh{钱} (argent). Litt. « argent pour menues dépenses ». Tons : 2-2-1.
\item \zh{存钱} : \zh{存} (cún, 2\up{e} ton) = déposer, conserver + \zh{钱}. Ne pas confondre avec \zh{春} (chūn, printemps).
\item \zh{觉得} : le second caractère \zh{得} se prononce ton neutre (léger) \emph{de} quand le verbe signifie « penser/trouver ». Ex: \zh{我觉得好} (juéde hǎo).
\item \zh{应该} : exprime un devoir moral ou un conseil avisé. Plus fort que \zh{可以} (pouvoir), moins contraignant que \zh{必须} (devoir absolu).
\end{itemize}
\end{aretenir}

\courssub{Dialogues modèles}

\begin{exemplebox}[Dialogue 1 : Demander de l'argent de poche (contexte familial)]
\textbf{Contexte :} Un élève demande à sa mère de l'argent pour la semaine.
\zh{儿：妈妈，这周的零花钱给我吧。}\\
\emph{Ér: Māma, zhè zhōu de língqiánhuā gěi wǒ ba.}\\
« Maman, donne-moi l'argent de poche de cette semaine. »\\[4pt]
\zh{妈：多少？上周不是给了你五千法郎吗？}\\
\emph{Mā: Duōshao? Shàng zhōu bùshì gěi le nǐ wǔqiān Fǎláng ma?}\\
« Combien ? La semaine dernière, je ne t'ai pas donné 5000 francs ? »\\[4pt]
\zh{儿：我存了一点，但不够买本子和坐车。}\\
\emph{Ér: Wǒ cún le yìdiǎn, dàn bùgòu mǎi běnzi hé zuò chē.}\\
« J'en ai économisé un peu, mais ça ne suffit pas pour acheter des cahiers et le transport. »\\[4pt]
\zh{妈：好吧，这是两千。记得存钱，别乱花。}\\
\emph{Mā: Hǎoba, zhè shì liǎngqiān. Jìde cúnqián, bié luàn huā.}\\
« Bon, voilà 2000. N'oublie pas d'économiser, ne dépense pas n'importe comment. »
\end{exemplebox}

\begin{exemplebox}[Dialogue 2 : Échanger entre camarades (contexte cour de récréation)]
\textbf{Contexte :} Deux amis comparent leur gestion de budget.
\zh{甲：你每个月有多少零花钱？}\\
\emph{Jiǎ: Nǐ měi gè yuè yǒu duōshao língqiánhuā?}\\
« Toi, tu as combien d'argent de poche par mois ? »\\[4pt]
\zh{乙：三千法郎。我通常存一半，花一半。}\\
\emph{Yǐ: Sānqiān Fǎláng. Wǒ tōngcháng cún yībàn, huā yībàn.}\\
« 3000 francs. D'habitude, j'économise la moitié, je dépense la moitié. »\\[4pt]
\zh{甲：我全花了，买零食、充话费。下个月我要学着存钱。}\\
\emph{Jiǎ: Wǒ quán huā le, mǎi língshí, chōng huàfèi. Xià gè yuè wǒ yào xuézhe cúnqián.}\\
« Moi, j'ai tout dépensé : snacks, recharge téléphone. Le mois prochain, je vais apprendre à économiser. »
\end{exemplebox}

\courssub{Formules fonctionnelles pour le débat}

\begin{propriete}[Structures syntaxiques du débat oral]
\begin{center}
\begin{tabularx}{\linewidth}{|X|X|X|}
\hline
\rowcolor{popblueL} Fonction & Structure type & Exemple (Caractères + Pinyin + Traduction) \\
\hline
Ouvrir / Donner son avis & \cle{我觉得 + [phrase]} & \zh{我觉得学生应该有零花钱。} (Wǒ juéde xuésheng yīnggāi yǒu língqiánhuā.) — Je pense que les élèves doivent avoir de l'argent de poche. \\
\hline
Argumenter (Cause $\to$ Conséquence) & \cle{因为 + cause，所以 + conséquence} & \zh{因为要坐车，所以需要钱。} (Yīnwèi yào zuò chē, suǒyǐ xūyào qián.) — Parce qu'il faut prendre le transport, il faut de l'argent. \\
\hline
Exprimer l'accord & \cle{我同意 (+ [ton avis])} & \zh{我同意，零花钱很有用。} (Wǒ tóngyì, língqiánhuā hěn yǒuyòng.) — Je suis d'accord, l'argent de poche est très utile. \\
\hline
Exprimer le désaccord (poli) & \cle{我不太同意 (+ [argument])} & \zh{我不太同意。因为会浪费钱。} (Wǒ bù tài tóngyì. Yīnwèi huì làngfèi qián.) — Je ne suis pas tout à fait d'accord. Parce que ça gaspille l'argent. \\
\hline
Relancer / Demander l'avis & \cle{你呢？} / \cle{你觉得怎么样？} & \zh{乙方，你呢？你觉得怎么样？} (Yǐfāng, nǐ ne? Nǐ juéde zěnmeyàng?) — Équipe B, et vous ? Qu'en pensez-vous ? \\
\hline
Contredire / Nuancer (Réplique) & \cle{可是 / 不过，因为…所以…} & \zh{可是，因为有零花钱可以学存钱，所以很有用。} (Kěshì, yīnwèi yǒu língqiánhuā kěyǐ xué cúnqián, suǒyǐ hěn yǒuyòng.) — Mais, parce qu'avec l'argent de poche on apprend à économiser, c'est utile. \\
\hline
Conclure (au nom du groupe) & \cle{所以，我们认为 + [synthèse]} & \zh{所以，我们认为应该有零花钱，但不多。} (Suǒyǐ, wǒmen rènwéi yīnggāi yǒu língqiánhuā, dàn bù duō.) — Donc, nous pensons qu'il faut de l'argent de poche, mais pas beaucoup. \\
\hline
\end{tabularx}
\end{center}
\end{propriete}

\begin{methode}[Comment construire un argument \zh{因为…所以…} en 3 étapes]
\begin{enumerate}
\item \textbf{Identifier la cause} (fait, besoin, contrainte) : \zh{学生要买书} (Les élèves doivent acheter des livres).
\item \textbf{Identifier la conséquence} (action, nécessité, résultat) : \zh{他们需要零花钱} (Ils ont besoin d'argent de poche).
\item \textbf{Assembler} : \zh{因为学生要买书，所以他们需要零花钱。} \emph{(Yīnwèi xuésheng yào mǎi shū, suǒyǐ tāmen xūyào língqiánhuā.)}
\end{enumerate}
\emph{Astuce :} Garde l'ordre Sujet + Verbe dans les deux parties. Ne mets pas le sujet après \zh{所以} s'il est le même que celui de la cause (on peut l'omettre ou le répéter pour la clarté).
\end{methode}

\courssub{Les tons à travailler : entraînement ciblé}

\begin{propriete}[Règles de ton et pièges pour cette leçon]
\begin{itemize}
\item \textbf{\zh{钱} (qián) : 2\up{e} ton montant.} Beaucoup de francophones le prononcent plat (1\up{e} ton) ou descendant. \textbf{Exercice :} \zh{零花钱} (líng-\textbf{qián}-huā), \zh{存钱} (cún-\textbf{qián}), \zh{没钱} (méi-\textbf{qián}).
\item \textbf{\zh{应该} (yīnggāi) : 1\up{e} + 1\up{e} ton (hauts et plats).} Ne pas faire monter \zh{gāi}.
\item \textbf{\zh{同意} (tóngyì) : 2\up{e} + 4\up{e} ton.} \zh{同} monte, \zh{意} tombe sèchement.
\item \textbf{\zh{觉得} (juéde) : 2\up{e} ton + ton neutre.} \zh{得} perd son ton propre (4\up{e} ton isolé) et devient léger \emph{de}.
\item \textbf{\zh{因为} (yīnwèi) : 1\up{e} + 4\up{e} ton.} \zh{为} ton descendant.
\item \textbf{\zh{所以} (suǒyǐ) : 3\up{e} + 3\up{e} ton.} Règle du double 3\up{e} ton : le premier \zh{所} se prononce 2\up{e} ton (\textbf{suó}) à l'oral ! \rightarrow \emph{suóyǐ}.
\end{itemize}
\end{propriete}

\begin{experience}[Atelier tons : « Le marché aux arguments »]
\begin{enumerate}
\item Chaque élève pioche une carte « Argument » (ex: \zh{因为要买水} / \zh{因为想存钱}).
\item Il doit le lire à voix haute en respectant \textbf{scrupuleusement} les tons de \zh{因为} (suóyǐ) et du vocabulaire budget.
\item La classe valide (pouce en l'air) ou corrige (pouce en bas + modèle du prof).
\item On enchaîne avec la conséquence : \zh{所以…} (suǒyǐ…).
\end{enumerate}
\end{experience}

\courssub{Jeu de rôle et débat guidé}

\begin{experience}[Mise en situation : « Le Grand Débat du Club Chinois »]
\textbf{Organisation (5 min) :}
\begin{itemize}
\item Classe divisée en deux grands groupes : \zh{甲方} (Pour) / \zh{乙方} (Contre).
\item Rôles par groupe : 1 \textbf{Animateur} (ouvre, relance, gère le temps), 3 \textbf{Orateurs} (1 argument chacun), 1 \textbf{Rapporteur} (conclut), les autres = \textbf{Observateurs} (grille d'évaluation).
\end{itemize}

\textbf{Préparation écrite (15 min) :}
Chaque orateur rédige son argument complet sur une fiche :
\zh{论点 : 因为……，所以……。}
\emph{Exemple Orateur A1 :} \zh{因为学生每天要坐车、买水、买笔，所以他们需要零花钱。}
\emph{Exemple Orateur B1 :} \zh{因为有的学生拿零花钱买游戏、吃零食，所以零花钱会被浪费。}

\textbf{Déroulement du débat (25 min) :}
\begin{enumerate}
\item \textbf{Animateur A :} \zh{大家好！今天我们辩论：学生应该有零花钱吗？甲方认为：应该有。我觉得学生应该有零花钱。}
\item \textbf{Orateur A1, A2, A3 :} Présentent leurs 3 arguments (\zh{因为…所以…}).
\item \textbf{Animateur A :} \zh{乙方，你呢？你觉得怎么样？}
\item \textbf{Animateur B :} \zh{我不太同意。乙方认为：不应该有。}
\item \textbf{Orateur B1, B2, B3 :} Présentent leurs 3 arguments.
\item \textbf{Phase de réplique (libre, 5 min) :} Les animateurs relancent : \zh{甲方/乙方，你们怎么看对方的观点？} Les orateurs répondent avec \zh{可是/不过，因为…所以…}.
\item \textbf{Rapporteur A :} \zh{所以，我们认为学生应该有零花钱，但是不应该太多，要学会存钱。}
\item \textbf{Rapporteur B :} \zh{所以，我们认为学生不应该经常有零花钱，父母直接买需要的东西比较好。}
\end{enumerate}
\end{experience}

\courssub{Production orale guidée et évaluation}

\begin{exoresolu}[Débat modèle complet — Extraits commentés]
Reconstitue un débat complet à partir des éléments préparés : ouverture, arguments, relance, réplique, conclusions.
\tcblower
\textbf{Animateur A :} 大家好！今天我们辩论：学生应该有零花钱吗？甲方认为：应该有。我觉得学生应该有零花钱。\\
\emph{Dàjiā hǎo! Jīntiān wǒmen biànlùn: Xuésheng yīnggāi yǒu língqiánhuā ma? Jiǎfāng rènwéi: yīnggāi yǒu. Wǒ juéde xuésheng yīnggāi yǒu língqiánhuā.}\\
« Bonjour à tous ! Aujourd'hui nous débattons : les lycéens devraient-ils avoir de l'argent de poche ? L'équipe A pense : oui. Je pense que les lycéens doivent recevoir de l'argent de poche. »

\textbf{Orateur A1 :} 因为学生每天要坐车去学校、买水喝、买笔记本，所以他们需要零花钱。\\
\emph{Yīnwèi xuésheng měitiān yào zuò chē qù xuéxiào, mǎi shuǐ hē, mǎi bǐjìběn, suǒyǐ tāmen xūyào língqiánhuā.}\\
« Parce que les élèves doivent chaque jour prendre le transport pour l'école, acheter de l'eau à boire, acheter des cahiers, ils ont besoin d'argent de poche. »

\textbf{Animateur A :} 乙方，你呢？你觉得怎么样？\\
\emph{Yǐfāng, nǐ ne? Nǐ juéde zěnmeyàng?}\\
« Équipe B, et vous ? Qu'en pensez-vous ? »

\textbf{Animateur B :} 我不太同意。乙方认为：学生不应该有零花钱。\\
\emph{Wǒ bù tài tóngyì. Yǐfāng rènwéi: Xuésheng bù yīnggāi yǒu língqiánhuā.}\\
« Je ne suis pas tout à fait d'accord. L'équipe B pense : les élèves ne devraient pas avoir d'argent de poche. »

\textbf{Orateur B1 :} 因为有的学生拿零花钱买游戏、买零食，不买有用的东西，所以零花钱会被浪费。\\
\emph{Yīnwèi yǒu de xuésheng ná língqiánhuā mǎi yóuxì, mǎi língshí, bù mǎi yǒuyòng de dōngxi, suǒyǐ língqiánhuā huì bèi làngfèi.}\\
« Parce que certains élèves prennent l'argent de poche pour acheter des jeux, des snacks, n'achètent pas de choses utiles, l'argent de poche sera gaspillé. »

\textbf{Orateur A2 (Réplique) :} 可是，因为有零花钱，学生可以学习怎么存钱、怎么管钱，所以零花钱很有用。\\
\emph{Kěshì, yīnwèi yǒu língqiánhuā, xuésheng kěyǐ xuéxí zěnme cúnqián, zěnme guǎn qián, suǒyǐ língqiánhuā hěn yǒuyòng.}\\
« Mais, parce qu'avec l'argent de poche les élèves peuvent apprendre comment économiser, comment gérer l'argent, l'argent de poche est très utile. »

\textbf{Rapporteur A :} 所以，我们认为学生应该有零花钱，但是不应该太多，父母要教他们存钱。\\
\emph{Suǒyǐ, wǒmen rènwéi xuésheng yīnggāi yǒu língqiánhuā, dànshì bù yīnggāi tài duō, fùmǔ yào jiāo tāmen cúnqián.}\\
« Donc, nous pensons que les élèves doivent avoir de l'argent de poche, mais pas trop, les parents doivent leur apprendre à économiser. »

\textbf{Rapporteur B :} 所以，我们认为学生不应该经常有零花钱，父母可以直接给学生买需要的东西，这样不浪费。\\
\emph{Suǒyǐ, wǒmen rènwéi xuésheng bù yīnggāi jīngcháng yǒu língqiánhuā, fùmǔ kěyǐ zhíjiē gěi xuésheng mǎi xūyào de dōngxi, zhèyàng bù làngfèi.}\\
« Donc, nous pensons que les élèves ne devraient pas souvent avoir d'argent de poche, les parents peuvent directement acheter aux élèves les choses nécessaires, comme ça pas de gaspillage. »

\textbf{Commentaire linguistique :}
\begin{itemize}
\item \textbf{Registre :} L'usage de \zh{大家好} et \zh{今天我们辩论…} instaure un cadre formel adapté à l'exposé oral.
\item \textbf{Politesse du désaccord :} \zh{我不太同意} (litt. « je ne suis pas trop d'accord ») est la norme en contexte scolaire/chinois pour exprimer une opposition sans heurter. \zh{我不同意} est perçu comme brutal.
\item \textbf{Cohésion :} Les connecteurs \zh{可是} (réplique), \zh{但是} (nuance) structurent l'échange.
\item \textbf{Passage Je/Nous :} Les orateurs utilisent \zh{我觉得} (je), les rapporteurs \zh{我们认为} (nous), marquant la synthèse collective.
\item \textbf{Vocabulaire ciblé :} \zh{坐车}, \zh{买水}, \zh{买零食}, \zh{存钱}, \zh{管钱}, \zh{浪费}, \zh{有用} sont tous réinvestis.
\end{itemize}
\end{exoresolu}

\begin{methode}[Grille d'évaluation orale — Débat « Gestion du budget » (sur 20)]
Utilise cette grille pour t'auto-évaluer ou évaluer un camarade (observateur).
\begin{center}
\begin{tabularx}{\linewidth}{|X|c|X|}
\hline
\rowcolor{popblueL} \textbf{Critère} & \textbf{Points} & \textbf{Indicateurs de réussite} \\
\hline
\textbf{Tons \& Prononciation} & 5 / 5 & Tons corrects sur \zh{零花钱, 应该, 同意, 觉得, 因为, 所以} (attention \zh{所以} $\to$ \emph{suóyǐ}). Pas de confusion \zh{qian/qián}. \\
\hline
\textbf{Formules fonctionnelles} & 5 / 5 & Utilisation correcte d'au moins \textbf{4 formules} distinctes : \zh{我觉得}, \zh{因为…所以…}, \zh{我不太同意}, \zh{你呢/你觉得怎么样}, \zh{可是/不过}, \zh{我们认为}. \\
\hline
\textbf{Fluidité \& Correction} & 5 / 5 & Phrases complètes (SVO), ordre des mots respecté (compléments de lieu/temps avant verbe), peu d'hésitations, auto-correction possible. \\
\hline
\textbf{Interaction \& Écoute} & 5 / 5 & Répond au contenu de l'argument adverse (pas un monologue), relance le partenaire, respecte le tour de parole, langage corporel engagé. \\
\hline
\end{tabularx}
\end{center}
\end{methode}

\begin{exercice}[Entraînement individuel — À faire à la maison ou en classe]
\begin{enumerate}
\item \textbf{Écriture :} Écris en caractères (avec pinyin) 3 arguments \textbf{POUR} et 3 arguments \textbf{CONTRE} l'argent de poche, en utilisant \zh{因为…所以…} et au moins 3 mots du vocabulaire par argument.
\item \textbf{Oral :} Enregistre-toi (téléphone) en simulant l'animateur A qui ouvre le débat + l'orateur A1 + la relance vers l'équipe B. Écoute-toi : les tons de \zh{应该, 零花钱, 所以} sont-ils stables ?
\item \textbf{Traduction (Thème) :} Traduis en chinois :
    \begin{enumerate}
    \item « Je pense qu'économiser de l'argent est très utile. »
    \item « Parce qu'il veut acheter un téléphone, il a besoin de beaucoup d'argent. »
    \item « Nous ne sommes pas d'accord : l'argent de poche n'est pas gaspillé. »
    \item « Et toi ? Qu'est-ce que tu en penses ? »
    \end{enumerate}
\end{enumerate}
\end{exercice}

\begin{corrige}[Exercice — Corrigé commenté]
\begin{enumerate}
\item \textbf{Arguments POUR :}
    \begin{itemize}
    \item \zh{因为学生要坐车上学，所以需要零花钱。} (Transport)
    \item \zh{因为零花钱可以让学生学习存钱，所以很有用。} (Éducation financière)
    \item \zh{因为买文具、买水都要钱，所以学生应该有零花钱。} (Besoins quotidiens)
    \end{itemize}
\textbf{Arguments CONTRE :}
    \begin{itemize}
    \item \zh{因为学生会买零食、玩游戏，所以零花钱会被浪费。} (Gaspillage)
    \item \zh{因为父母可以直接买需要的东西，所以不需要给零花钱。} (Alternative parentale)
    \item \zh{因为有零花钱学生不学习存钱，所以不好。} (Risque éducatif)
    \end{itemize}
\item \textbf{Oral :} Vérifie particulièrement : \zh{yīnggāi} (1-1), \zh{língqiánhuā} (2-2-1), \zh{suǒyǐ} (lu 2-3 \rightarrow \emph{suóyǐ}).
\item \textbf{Traduction :}
    \begin{enumerate}
    \item \zh{我觉得存钱很有用。} (Wǒ juéde cúnqián hěn yǒuyòng.)
    \item \zh{因为他想买手机，所以需要很多钱。} (Yīnwèi tā xiǎng mǎi shǒujī, suǒyǐ xūyào hěn duō qián.)
    \item \zh{我们不同意：零花钱没有被浪费。} (Wǒmen bù tóngyì: Língqiánhuā méiyǒu bèi làngfèi.) \emph{Note : forme passive \zh{被} pour « être gaspillé »}.
    \item \zh{你呢？你觉得怎么样？} (Nǐ ne? Nǐ juéde zěnmeyàng?)
    \end{enumerate}
\end{enumerate}
\end{corrige}

\courssub{Bilan et prolongements}

Cette activité a permis de réunir, dans une situation de communication authentique (débat contradictoire), toutes les compétences de la leçon :
\begin{itemize}
\item \textbf{Lexique :} Maîtrise active du vocabulaire budget (\zh{零花钱, 存钱, 浪费, 需要, 买}…).
\item \textbf{Grammaire :} Structure argumentative \zh{因为…所以…}, verbes modaux (\zh{应该}), adverbes de degré (\zh{太, 很}), connecteurs logiques (\zh{可是, 但是, 所以}).
\item \textbf{Fonctionnel :} Formules codifiées pour l'interaction orale (ouvrir, argumenter, accorder, désaccorder poliment, relancer, conclure).
\item \textbf{Phonologie :} Travail ciblé sur les tons instables pour francophones (\zh{qián, suǒyǐ, juéde, yīnggāi}).
\item \textbf{Socioculturel :} Ancrage dans la réalité camerounaise (transport, cantine, francs CFA, club de Yaoundé) et découverte de la culture du débat structuré \zh{甲方/乙方} très prisée en Chine (compétitions \zh{辩论赛}).
\end{itemize}

\begin{savaistu}[Culture : L'argent de poche en Chine et au Cameroun]
\begin{itemize}
\item \textbf{Cameroun :} L'argent de poche (\zh{零花钱}) sert souvent aux frais quotidiens : transport (taxi-moto, car), cantine (\zh{食堂} / \zh{小卖部}), communication (\zh{充话费}). Montants variables selon les familles (1000–5000 FCFA/semaine).
\item \textbf{Chine :} L'argent de poche existe (\zh{零花钱}), mais la culture de l'\textbf{enveloppe rouge} (\zh{红包} \emph{hóngbāo}) au Nouvel An, mariages, anniversaires est majeure. Les enfants reçoivent des sommes importantes qu'ils remettent souvent aux parents pour « garder » (\zh{帮你保管}) ou mettre à la banque. L'éducation financière (\zh{理财教育}) commence tôt : tirelires, comptes bancaires pour mineurs, applications de gestion.
\item \textbf{Point commun :} La préoccupation parentale : « Apprendre à gérer l'argent sans le gaspiller » (\zh{学会理财，不乱花钱}).
\end{itemize}
\end{savaistu}

\begin{experience}[Prolongement : Débat « Sujet libre »]
Change de sujet pour réinvestir les structures. Propositions :
\begin{itemize}
\item \zh{学生应该带手机上学吗？} (Les élèves doivent-ils apporter leur téléphone à l'école ?)
\item \zh{周末应该补课吗？} (Faut-il faire des cours le week-end ?)
\item \zh{穿校服好不好？} (L'uniforme scolaire : bien ou pas ?)
\end{itemize}
Même organisation, même grille. L'objectif : automatiser les formules et stabiliser les tons dans de nouveaux contenus lexicaux.
\end{experience}

\newpage

\coursec{Méthodes et exercices résolus}

\coursec{Leçon 28 — Expression orale : gestion du budget}

\courssub{Vocabulaire clé et structures pour présenter un budget personnel}

\begin{methode}[Présenter oralement son budget mensuel en 5 étapes]
\textbf{Objectif :} Construire un exposé court (1--2 min) clair et structuré sur ses revenus et dépenses.
\begin{enumerate}
    \item \textbf{Saluer et annoncer le thème :} Utiliser une formule d'ouverture standard (\zh{大家好，今天我要谈谈我的预算。} -- Dàjiā hǎo, jīntiān wǒ yào tántan wǒ de yùsuàn. / « Bonjour à tous, aujourd'hui je vais parler de mon budget. »).
    \item \textbf{Présenter les revenus (收入 shōurù) :} Citer la source (argent de poche, job d'été, bourse) et le montant approximatif. Structure : \cle{Sujet + 收入 + 是 + Montant + 元/块}. Exemple : \zh{我的收入主要是零花钱，每月大概五百元。} (Wǒ de shōurù zhǔyào shì línghuáqián, měiyuē dàgài wǔbǎi yuán.)
    \item \textbf{Détailer les dépenses principales (支出 zhīchū) :} Utiliser les classificateurs adaptés (\cle{笔 bǐ} pour une somme/dépense, \cle{项 xiàng} pour un poste). Structure : \cle{Sujet + 在 + Domaine + 上 + 动词 + 金额}. Exemple : \zh{我在餐饮上花了两百块。} (Wǒ zài cānyǐn shàng huā le liǎngbǎi kuài.)
    \item \textbf{Exprimer l'épargne ou le déficit (结余 jiéyú / 赤字 chìzì) :} Structure : \cle{收入 - 支出 = 结余/赤字}. Exemple : \zh{最后剩下一百元，我存进银行。} (Zuìhòu shèng xià yībǎi yuán, wǒ cún jìn yínháng.)
    \item \textbf{Conclure et ouvrir au débat :} \zh{这就是我的预算情况，谢谢大家。你们呢？} (Zhè jiùshì wǒ de yùsuàn qíngkuàng, xièxie dàjiā. Nǐmen ne?)
\end{enumerate}
\cle{Réflexe :} Toujours placer le complément de lieu (\zh{在...上}) \textbf{avant} le verbe (\zh{花}, \zh{用}) et le complément de quantité \textbf{après} le verbe (\zh{花了两百块}).
\end{methode}

\begin{textebox}[Modèle d'exposé : Mon budget de lycéen à Yaoundé]
大家好，今天我介绍我的月预算。\\
Dàjiā hǎo, jīntiān wǒ jièshào wǒ de yuè yùsuàn.\\
(Bonjour à tous, aujourd'hui je présente mon budget mensuel.)

我的收入是父母给的零花钱，大约一百一十元。\\
Wǒ de shōurù shì fùmǔ gěi de línghuáqián, dàyuē yībǎi yīshí yuán.\\
(Mon revenu est l'argent de poche de mes parents, environ 110 yuans.)

我的支出有三项。第一，交通费，坐摩托车或公交车，花三十元。\\
Wǒ de zhīchū yǒu sān xiàng. Dì yī, jiāotōngfèi, zuò mótuōchē huò gōngjiāochē, huā sānshí yuán.\\
(Mes dépenses ont trois postes. Premier, le transport, moto ou bus, dépense 30 yuans.)

第二，餐饮费，中午吃米饭或恩多莱，花二十五元。\\
Dì èr, cānyǐnfèi, zhōngwǔ chī mǐfàn huò ēnduōlái, huā èrshíwǔ yuán.\\
(Deuxième, les repas, midi riz ou ndolé, dépense 25 yuans.)

第三，话费，买流量和通话时间，花五元。\\
Dì sān, huàfèi, mǎi liúliàng hé tōnghuà shíjiān, huā wǔ yuán.\\
(Troisième, le crédit, acheter data et temps d'appel, dépense 5 yuans.)

总共花六十元。收入减去支出，结余五十元。\\
Zǒnggòng huā liùshí yuán. Shōurù jiǎnqù zhīchū, jiéyú wǔshí yuán.\\
(Total dépensé 60 yuans. Revenu moins dépenses, solde 50 yuans.)

我把结余存进手机银行备用。这就是我的预算，谢谢大家！\\
Wǒ bǎ jiéyú cún jìn shǒujī yínháng bèiyòng. Zhè jiùshì wǒ de yùsuàn, xièxie dàjiā!\\
(Je mets le solde sur ma banque mobile pour secours. Voilà mon budget, merci tous !)
\end{textebox}

\begin{exoresolu}[Exposé type : Rédaction guidée à partir de données]
\textbf{Énoncé :} À partir des données ci-dessous, rédigez en chinois (caractères + pinyin) un exposé oral de 8--10 phrases pour le présenter devant la classe. Utilisez le vocabulaire du tableau.
\begin{itemize}
    \item Revenus : Argent de poche des parents (父母给的零花钱) : 10 000 FCFA / mois $\approx$ 110 CNY.
    \item Dépenses : Transport (moto-taxi/bus) : 30 000 FCFA / mois ; Repas midi (ndolé/riz) : 25 000 FCFA / mois ; Crédit téléphone : 5 000 FCFA / mois.
    \item Épargne : Reste $\approx$ 0 FCFA (équilibre serré).
\end{itemize}
\textbf{Vocabulaire fourni :}
\begin{center}\begin{tabularx}{\linewidth}{|X|X|X|X|}\hline \rowcolor{popblueL} Caractères & Pinyin & Nature & Traduction \\ \hline
收入 & shōurù & nom/verbe & revenu / gagner de l'argent \\ \hline
支出 & zhīchū & nom/verbe & dépense / dépenser \\ \hline
零花钱 & línghuáqián & nom & argent de poche \\ \hline
交通 & jiāotōng & nom & transport \\ \hline
餐饮 & cānyǐn & nom & restauration / repas \\ \hline
话费 & huàfèi & nom & frais de téléphone / crédit \\ \hline
结余 & jiéyú & verbe/nom & solde positif / épargner le reste \\ \hline
赤字 & chìzì & nom & déficit / solde négatif \\ \hline
预算 & yùsuàn & nom/verbe & budget / budgéter \\ \hline
大约 & dàyuē & adv. & environ / à peu près \\ \hline
\end{tabularx}\end{center}

\tcblower
\textbf{Solution détaillée :}
\textbf{1. Analyse grammaticale :} Le texte suit la structure \textit{Sujet + Temps + Lieu + Verbe + Objet/Complément}. Les montants sont donnés en Yuan (元 yuán / 块 kuài) pour l'exercice chinois standard, mais l'élève peut mentionner FCFA (非洲法郎 Fēizhōu Fǎláng) en précisant la devise.
\textbf{2. Rédaction de l'exposé (Caractères + Pinyin + Traduction) :} (Voir le modèle ci-dessus dans la \texttt{textebox} « Modèle d'exposé »).
\textbf{3. Points de vigilance (Tons \& Ordre) :}
\begin{itemize}
    \item \zh{大约} (dàyuē) : 4e ton + 1er ton (ne pas dire *dàyuě).
    \item \zh{摩托车} (mótuōchē) : 2e + 1er + 1er ton.
    \item \zh{恩多莱} (Ēnduōlái) : Transcription phonétique pour \textbf{Ndolé} (1er + 1er + 2e ton).
    \item Structure \zh{花 + 金额} (verbe + quantité) : \zh{花三十元} et non *\zh{三十元花}.
    \item \zh{把} (bǎ) pour la manipulation de l'objet \zh{结余} avant le verbe \zh{存}.
\end{itemize}
\textbf{Conclusion :} L'exposé est fluide, utilise les connecteurs logiques (\zh{第一/第二/第三}), respecte l'ordre des mots SVO et la place des compléments.
\end{exoresolu}

\courssub{Formules fonctionnelles pour interagir et débattre}

\begin{methode}[Participer à un débat sur la gestion d'argent : 4 formules types]
\textbf{Objectif :} Réagir, donner son avis, nuancer, relancer l'interlocuteur.
\begin{enumerate}
    \item \textbf{Exprimer son accord / désaccord :}
    \begin{itemize}
        \item Accord : \zh{我赞同你的看法。} (Wǒ zàntóng nǐ de kànfǎ.) / \zh{说得很对。} (Shuō de hěn duì.)
        \item Désaccord poli : \zh{我有不同的看法。} (Wǒ yǒu bùtóng de kànfǎ.) / \zh{我不太同意，因为……} (Wǒ bù tài tóngyì, yīnwèi…)
    \end{itemize}
    \item \textbf{Donner son avis / Argumenter :} Structure \cle{我认为 / 我觉得 + Sujet + Prédicat}. \zh{我认为存钱很重要，以备不时之需。} (Wǒ rènwéi cúnqián hěn zhòngyào, yǐbèi bùshí zhī xū. -- Je pense qu'épargner est important, pour prévoir les imprévus.)
    \item \textbf{Demander des précisions / Relancer :} \zh{你能具体说说吗？} (Nǐ néng jùtǐ shuōshuo ma ?) / \zh{比如说呢？} (Bǐrú shuō ne ?) / \zh{那你怎么处理……？} (Nà nǐ zěnme chǔlǐ… ?)
    \item \textbf{Résumer / Conclure :} \zh{总之，量入为出最关键。} (Zǒngzhī, liàngrùwéichū zuì guānjiàn. -- En résumé, vivre selon ses moyens est la clé.)
\end{enumerate}
\cle{Structure clé :} \zh{量入为出} (liàng rù wéi chū) -- Chengyu (成语) 4 caractères : « Mesurer les revenus pour décider les dépenses » = Vivre selon ses moyens.
\end{methode}

\begin{exoresolu}[Transformation de phrases \& Complétion de dialogue]
\textbf{Énoncé :}
\textbf{A.} Transformez les phrases françaises en chinois en utilisant les formules fonctionnelles ci-dessus.
1. Je ne suis pas d'accord, car dépenser tout son argent est dangereux. \\
2. Pourriez-vous donner un exemple concret d'épargne ? \\
3. En résumé, il faut noter ses dépenses chaque jour.

\textbf{B.} Complétez le dialogue suivant (choisissez parmi : \zh{量入为出} / \zh{以备不时之需} / \zh{记账} / \zh{月光族}).
甲：你每个月都花光工资吗？(Tu dépenses tout ton salaire chaque mois ?)
乙：以前是，我是典型的 \underline{\hspace{3cm}}。
甲：现在呢？你怎么改变的？
乙：我开始 \underline{\hspace{3cm}}，坚持 \underline{\hspace{3cm}}，把结余 \underline{\hspace{3cm}} 存起来。

\tcblower
\textbf{Solution détaillée :}
\textbf{A. Transformation (Thème) :}
1. \zh{我不太同意，因为把钱花光很危险。} (Wǒ bù tài tóngyì, yīnwèi bǎ qián huāguāng hěn wēixiǎn.) \textit{Note : \zh{花光} = dépenser entièrement (complément de résultat).}
2. \zh{你能举个具体的存钱例子吗？} (Nǐ néng jǔ gè jùtǐ de cúnqián lìzi ma ?) \textit{Note : \zh{举个例子} (jǔ gè lìzi) = donner un exemple.}
3. \zh{总之，每天记账最关键。} (Zǒngzhī, měitiān jìzhàng zuì guānjiàn.) \textit{Note : \zh{记账} (jìzhàng) = tenir ses comptes / noter dépenses.}

\textbf{B. Complétion (Vocabulaire en contexte) :}
甲：你每个月都花光工资吗？
乙：以前是，我是典型的 \textbf{\zh{月光族}} (yuèguāngzú -- « clan lune claire » = ceux qui dépensent tout leur salaire avant la fin du mois).
甲：现在呢？你怎么改变的？
乙：我开始 \textbf{\zh{记账}} (jìzhàng), 坚持 \textbf{\zh{量入为出}} (liàngrùwéichū), 把结余 \textbf{\zh{以备不时之需}} (yǐbèi bùshí zhī xū) 存起来。
\textit{Analyse :} \zh{月光族} est un terme sociologique courant (HSK 4+). \zh{记账} est l'action concrète. \zh{量入为出} est le principe. \zh{以备不时之需} (pour parer aux besoins imprévus) justifie l'épargne.
\end{exoresolu}

\courssub{Les tons à travailler : paires minimales et phrases pièges}

\begin{methode}[Maîtriser les tons du vocabulaire budgétaire : Méthode « 3R »]
\textbf{Objectif :} Éviter les confusions tonales qui changent le sens (ex: \zh{买} acheter vs \zh{卖} vendre).
\begin{enumerate}
    \item \textbf{Repérer (Identifier) :} Isoler les paires minimales du lexique budget. Les écrire en pinyin avec traits de tons. Les classer par ton (1, 2, 3, 4, neutre).
    \item \textbf{Répéter (Drill) :} Méthode « Main sur la gorge » pour sentir la vibration (tons bas/3e) vs « Main levée » (ton montant/2e). Lire les tableaux de paires minimales 3 fois : lent $\rightarrow$ normal $\rightarrow$ rapide.
    \item \textbf{Réutiliser en phrase (Contextualiser) :} Insérer le mot dans une phrase type \zh{我要……} ou \zh{这个……很贵}. Vérifier les changements de ton (sandhi) pour \zh{一} (yī) et \zh{不} (bù).
\end{enumerate}
\cle{Focus Tons :} \zh{收} (shōu, 1er) vs \zh{手} (shǒu, 3e) ; \zh{支} (zhī, 1er) vs \zh{只} (zhǐ, 3e) ; \zh{预} (yù, 4e) vs \zh{雨} (yǔ, 3e) ; \zh{算} (suàn, 4e) vs \zh{蒜} (suàn, 4e -- homophone parfait, contexte différent).
\end{methode}

\begin{exoresolu}[Discrimination auditive et lecture tonale]
\textbf{Énoncé :}
\textbf{1. Discrimination (Simulation orale) :} L'enseignant lit une phrase. L'élève coche le mot entendu (A ou B).
\begin{itemize}
    \item Phrase 1 : \zh{我的 \textbf{收/手} 入增加了。} (Mon revenu / Ma main a augmenté.)
    \item Phrase 2 : \zh{请 \textbf{算/蒜} 一下总数。} (Calcule / L'ail le total.)
    \item Phrase 3 : \zh{这个 \textbf{预/雨} 算很准确。} (Ce budget / Cette pluie est précis.)
\end{itemize}
\textbf{2. Lecture à haute voix (Tons composés \& \zh{不} / \zh{一}) :} Lisez le paragraphe suivant en marquant les tons modifiés (sandhi) pour \zh{不} (bù $\rightarrow$ bú devant 4e ton) et \zh{一} (yī $\rightarrow$ yí/yì).
\zh{我不买贵的东西。一共五十元。不贵。}

\tcblower
\textbf{Solution détaillée :}
\textbf{1. Réponses attendues :}
\begin{itemize}
    \item Phrase 1 : \textbf{\zh{收}} (shōu, 1er ton). \textit{Contexte :} \zh{收入} (revenu). \zh{手入} n'existe pas.
    \item Phrase 2 : \textbf{\zh{算}} (suàn, 4e ton). \textit{Contexte :} \zh{算一下} (calculer un coup). \zh{蒜一下} n'a pas de sens (sauf « éplucher l'ail » \zh{剥蒜}).
    \item Phrase 3 : \textbf{\zh{预}} (yù, 4e ton). \textit{Contexte :} \zh{预算} (budget). \zh{雨算} n'existe pas.
\end{itemize}
\textbf{2. Lecture corrigée avec sandhi (Pinyin annoté) :}
\zh{我 \textbf{bù} 买 贵 的 东西。} (Wǒ \textbf{bù} mǎi guì de dōngxi.) -- \zh{不} (4e) + \zh{买} (3e) $\rightarrow$ \zh{不} garde son ton 4 \textbf{bù} (règle : \zh{不} devient 2e ton \textbf{uniquement} devant un 4e ton). \textit{(Piège classique)}.
\zh{一共 五十 元。} (Yìgòng wǔshí yuán.) -- \zh{一} yī (1er) + \zh{共} gòng (4e) $\rightarrow$ \zh{一} devient \textbf{yì} (4e ton) devant ton 4.
\zh{\textbf{Bú} 贵。} (\textbf{Bú} guì.) -- \zh{不} + \zh{贵} (4e ton) $\rightarrow$ \zh{不} devient \textbf{bú} (2e ton). \checkmark

\textbf{Tableau récapitulatif des tons du vocabulaire clé :}
\begin{center}\begin{tabularx}{\linewidth}{|X|X|X|X|}\hline \rowcolor{popblueL} Caractères & Pinyin & Ton & Sens \\ \hline
收入 & shōu rù & 1 - 4 & Revenu \\ \hline
支出 & zhī chū & 1 - 1 & Dépense \\ \hline
预算 & yù suàn & 4 - 4 & Budget \\ \hline
结余 & jié yú & 2 - 2 & Solde / Épargne \\ \hline
赤字 & chì zì & 4 - 4 & Déficit \\ \hline
月光族 & yuè guāng zú & 4 - 1 - 2 & « Clan lune claire » (dépensier) \\ \hline
记账 & jì zhàng & 4 - 4 & Tenir les comptes \\ \hline
量入为出 & liàng rù wéi chū & 4 - 4 - 2 - 1 & Vivre selon ses moyens \\ \hline
\end{tabularx}\end{center}
\end{exoresolu}

\courssub{Jeu de rôle guidé : Négociation au marché / Gestion imprévu}

\begin{methode}[Réussir un jeu de rôle « Achat / Imprévu budgétaire » en 4 temps]
\textbf{Contexte :} Scénario camerounais : Marché Mokolo (Yaoundé) ou Marché Central (Douala). Rôle A : Étudiant (budget serré). Rôle B : Vendeur / Parent / Ami.
\begin{enumerate}
    \item \textbf{Analyser la fiche de rôle :} Identifier \textit{Objectif} (acheter \zh{ndolé} + riz pour 2000 FCFA max), \textit{Contrainte} (prix affiché 2500 FCFA), \textit{Attitude} (poli, ferme).
    \item \textbf{Préparer les « blocs linguistiques » (Chunks) :}
    \begin{itemize}
        \item Demander prix : \zh{这个多少钱？} (Zhège duōshao qián ?) / \zh{能便宜点儿吗？} (Néng piányi diǎnr ma ?)
        \item Argumenter budget : \zh{我只有两千块。} (Wǒ zhǐ yǒu liǎngqiān kuài.) / \zh{超预算了。} (Chāo yùsuàn le.)
        \item Proposer compromis : \zh{少给点菜/米行吗？} (Shǎo gěi diǎn cài/mǐ xíng ma ?)
        \item Clôturer : \zh{成交 / 不买了。} (Chéngjiāo / Bù mǎi le.)
    \end{itemize}
    \item \textbf{Répéter avec partenaire :} Échanger les rôles. Focus sur l'intonation interrogative (montante pour \zh{吗}, descendante pour questions ouvertes).
    \item \textbf{Auto-évaluation :} Ai-je utilisé \zh{只} (seulement) ? Ai-je placé \zh{在} (lieu) avant le verbe ? Les tons de \zh{便宜} (piányi -- 2e + 2e/neutre) sont-ils justes ?
\end{enumerate}
\end{methode}

\begin{textebox}[Modèle de dialogue : Imprévu médical \& Réajustement budgétaire]
甲：哎，你看起来很担心。怎么了？\\
Āi, nǐ kàn qǐlái hěn dānxīn. Zěnme le ?\\
(Hé, tu as l'air inquiet. Qu'est-ce qu'il y a ?)

乙：我生病了，买药要五千法郎。可是我的零花钱不够了。\\
Wǒ shēngbìng le, mǎi yào yào wǔqiān Fǎláng. Kěshì wǒ de línghuáqián bù gòu le.\\
(Je suis malade, les médicaments coûtent 5000 FCFA. Mais mon argent de poche ne suffit plus.)

甲：那你原本的「娱乐费」怎么办？\\
Nà nǐ yuánběn de « yúlèfèi » zěnme bàn ?\\
(Et ton « budget loisirs » initial, tu en fais quoi ?)

乙：只能取消这个周末看电影、吃零食的计划了。\\
Zhǐ néng qǔxiāo zhège zhōumò kàndiànyǐng, chī língshí de jìhuà le.\\
(Je ne peux qu'annuler le projet ciné/grignotage de ce week-end.)

甲：如果实在不够，我可以借你两千，下次你还我。\\
Rúguǒ shízài bù gòu, wǒ kěyǐ jiè nǐ liǎngqiān, xiàcì nǐ huán wǒ.\\
(Si vraiment pas assez, je peux te prêter 2000, la prochaine fois tu me rembourses.)

乙：太感谢了！我下周发零花钱第一时间还你。\\
Tài gǎnxiè le! Wǒ xià zhōu fā línghuáqián dì yī shíjiān huán nǐ.\\
(Merci infiniment ! La semaine prochaine je reçois mon argent, je te rembourse dès que possible.)
\end{textebox}

\begin{exoresolu}[Scénario : Imprévu médical -- Production guidée]
\textbf{Énoncé :} Vous êtes l'élève A. Votre camarade (B) vous annonce qu'il doit payer 5000 FCFA pour une ordonnance (看病开药 kànbìng kāiyào) non prévue. Votre budget « Loisirs » (看电影 kàndiànyǐng / 吃零食 chī língshí) est de 5000 FCFA.
\textbf{Tâche :} Rédigez les 6 répliques du dialogue (A puis B, alterné) en chinois (caractères + pinyin + traduction). Utilisez : \zh{怎么办} (zěnme bàn -- que faire), \zh{借} (jiè -- prêter/emprunter), \zh{下次} (xiàcì -- la prochaine fois).

\tcblower
\textbf{Solution détaillée :}
\textbf{Dialogue type (Niveau HSK 3/4) :} (Voir le modèle ci-dessus dans la \texttt{textebox} « Modèle de dialogue »).
\textbf{Analyse grammaticale :}
\begin{itemize}
    \item \zh{不够了} (bù gòu le) : \zh{够} (gòu, suffire) + \zh{了} (changement d'état / nouveau statut « plus assez »).
    \item \zh{只能} (zhǐ néng) : « Ne pouvoir que... » (restriction). Place avant le verbe modal \zh{能}.
    \item \zh{取消...计划} (qǔxiāo... jìhuà) : Structure V + O standard.
    \item \zh{借} (jiè) : Ambiguïté \zh{借给} (jiè gěi -- prêter À qqn) vs \zh{跟...借} (gēn... jiè -- emprunter À qqn). Ici \zh{借你} (te prêter) contexte clair.
    \item \zh{第一时间} (dì yī shíjiān) : « Dès que possible / à la première occasion » (locution temporelle courante).
\end{itemize}
\end{exoresolu}

\courssub{Production orale finale : Exposé structuré \& Débat évalué}

\begin{methode}[Structurer un exposé oral noté (Bac) : Plan « 3 parties + 3 connecteurs »]
\textbf{Objectif :} Produire une intervention continue de 2--3 minutes (niveau B1/B2 oral) sur « La gestion du budget chez les jeunes Camerounais vs Chinois ».
\begin{enumerate}
    \item \textbf{Introduction (30s) :} Accroche + Sujet + Plan annoncé.
        \zh{大家好。今天的话题是「青年预算管理」。我将从三个方面谈谈：收入来源、主要支出、和建议。} (Dàjiā hǎo. Jīntiān de huàtí shì « Qīngnián yùsuàn guǎnlǐ ». Wǒ jiāng cóng sān gè fāngmiàn tántan: shōurù láiyuán, zhǔyào zhīchū, hé jiànyì.)
    \item \textbf{Développement (1min 30s) -- 3 arguments = 3 paragraphes oraux :}
        \begin{itemize}
            \item \textbf{Arg 1 (Revenus) :} \zh{首先，收入方面……} (Shǒuxiān, shōurù fāngmiàn…) -- Comparer : Parents (Cameroun) vs Bourses/Jobs étudiants (Chine).
            \item \textbf{Arg 2 (Dépenses) :} \zh{其次，支出结构不同……} (Qícì, zhīchū jiégòu bùtóng…) -- Transport/Alimentation (CM) vs Logement/Transport/Shopping (CN). Utiliser \zh{相比之下} (xiāngbǐ zhīxià -- en comparaison).
            \item \textbf{Arg 3 (Outils/Épargne) :} \zh{最后，工具和习惯……} (Zuìhòu, gōngjù hé xíguàn…) -- \zh{记账软件} (jìzhàng ruǎnjiàn -- apps budget) Alipay/WeChat Pay (CN) vs Cahier/Excel/Mobile Money (CM). \zh{量入为出} universel.
        \end{itemize}
    \item \textbf{Conclusion (30s) :} Synthèse + Ouverture débat.
        \zh{总之，无论在喀麦隆还是中国，记账和储蓄都是好习惯。建议同学们从今天开始记账。谢谢！} (Zǒngzhī, wúlùn zài Kāmàilóng háishì Zhōngguó, jìzhàng hé chǔxù dōu shì hǎo xíguàn. Jiànyì tóngxuémen cóng jīntiān kāishǐ jìzhàng. Xièxie!)
    \item \textbf{Phase Débat (Interaction) :} Écouter la question $\rightarrow$ Reformuler $\rightarrow$ Répondre avec \zh{我认为/因为/所以}.
\end{enumerate}
\cle{Critères d'évaluation (Grille Bac) :} 1. Prononciation/Tons (25\%) 2. Cohérence/Connecteurs (25\%) 3. Vocabulaire/Structures (25\%) 4. Interaction/Réactivité (25\%).
\end{methode}

\begin{textebox}[Modèle d'introduction -- Exposé Bac « Moonlight Clan »]
各位评委、同学们好。\\
Gèwèi píngwěi, tóngxuémen hǎo.\\
(Mesdames/Messieurs les jurés, chers camarades, bonjour.)

今天我要谈论的现象是「月光族」。\\
Jīntiān wǒ yào tánlùn de xiànxiàng shì « Yuèguāngzú ».\\
(Le phénomène dont je vais parler aujourd'hui est le « Moonlight Clan ».)

指那些月收入刚好等于月支出、甚至入不敷出、毫无积蓄的年轻人。\\
Zhǐ nàxiē yuè shōurù gānghǎo děngyú yuè zhīchū, shènzhì rù bù fú chū, háowú jīxù de niánqīng rén.\\
(Cela désigne ceux dont le revenu mensuel juste égal aux dépenses, voire insuffisant, sans aucune épargne, jeunes gens.)

这一现象在喀麦隆和中国的大城市都很普遍。\\
Zhè yī xiànxiàng zài Kāmàilóng hé Zhōngguó de dà chéngshì dōu hěn pǔbiàn.\\
(Ce phénomène est très répandu dans les grandes villes du Cameroun et de Chine.)

我认为，这既反映了消费升级的愿望，也暴露了理财教育的缺失。\\
Wǒ rènwéi, zhè jì fǎnyìng le xiāofèi shēngjí de yuànwàng, yě bàolù le lǐcái jiàoyù de quēshī.\\
(Je pense que cela reflète à la fois le souhait de montée en gamme de la consommation, et expose le manque d'éducation financière.)

下面，我将从成因、危害和对策三个方面展开。\\
Xiàmian, wǒ jiāng cóng chéngyīn, wèihài hé duìcè sān gè fāngmiàn zhǎnkāi.\\
(Ensuite, je vais développer sous trois aspects : causes, méfaits, et contre-mesures.)
\end{textebox}

\begin{exoresolu}[Sujet d'examen type : Rédaction du plan détaillé + Introduction écrite]
\textbf{Sujet :} « \zh{谈谈你对「月光族」现象的看法，并提出建议。} » (Parlez du phénomène « Moonlight Clan » et proposez des solutions). Temps de prép : 10 min. Temps de parole : 2 min.
\textbf{Travail demandé :} Rédigez le \textbf{plan détaillé} (mots-clés en chinois + pinyin) et l'\textbf{introduction complète} (caractères + pinyin + traduction).

\tcblower
\textbf{Solution détaillée :}
\textbf{1. Plan détaillé (提纲 Tígāng) :}
\begin{itemize}
    \item \textbf{开场} : \zh{定义「月光族」} (Dìngyì « Yuèguāngzú » -- Définir : revenus = dépenses, épargne = 0).
    \item \textbf{原因分析} (Yuányīn fēnxī -- Analyse causes) -- Connecteur : \zh{主要原因有二} (Zhǔyào yuányīn yǒu èr -- Deux causes principales) :
        \begin{enumerate}
            \item \zh{消费观念变化，冲动消费多} (Xiāofèi guānniàn biànhuà, chōngdòng xiāofèi duō -- Mentalité conso change, achats impulsifs).
            \item \zh{收入低，生活成本高（房租、交通）} (Shōurù dī, shēnghuó chéngběn gāo -- Revenus bas, coût vie élevé).
        \end{enumerate}
    \item \textbf{后果与危害} (Hòuguǒ yǔ wèihài -- Conséquences) : \zh{无储蓄，抗风险能力弱} (Wú chǔxù, kàng fēngxiǎn nénglì ruò -- Pas d'épargne, résilience faible).
    \item \textbf{建议对策} (Jiànyì duìcè -- Solutions) -- Connecteur : \zh{建议从三方面入手} (Jiànyì cóng sān fāngmiàn rùshǒu) :
        \begin{enumerate}
            \item \zh{强制记账，明细收支} (Qiángzhì jìzhàng, míngxì shōuzhī -- Tenir comptes obligatoirement, détailler recettes/dépenses).
            \item \zh{设立「应急基金」，先存后花} (Shèlì « Yìngjí jījīn », xiān cún hòu huā -- Créer fonds urgence, épargner d'abord).
            \item \zh{学习理财知识，避免网贷} (Xuéxí lǐcái zhīshì, bìmiǎn wǎngdài -- Apprendre finance, éviter prêts en ligne).
        \end{enumerate}
    \item \textbf{结语} : \zh{拒做「月光族」，从今天做起。} (Jù zuò « Yuèguāngzú », cóng jīntiān zuò qǐ -- Refuser d'être « Moonlight », commencer aujourd'hui.)
\end{itemize}

\textbf{2. Introduction complète (开场白 Kāichǎngbái) :} (Voir le modèle ci-dessus dans la \texttt{textebox} « Modèle d'introduction »).
\textbf{Points de niveau HSK 4+ validés :} \zh{入不敷出} (rù bù fú chū -- revenus ne couvrent pas dépenses, chengyu), \zh{暴露} (bàolù -- révéler/mettre à jour), \zh{理财教育} (lǐcái jiàoyù -- éducation financière), \zh{消费升级} (xiāofèi shēngjí -- montée en gamme conso), \zh{缺失} (quēshī -- manque/absence).
\end{exoresolu}

\begin{attention}[Les 5 erreurs qui coûtent des points au Bac (Leçon 28)]
\begin{enumerate}
    \item \textbf{Tons : Confusion \zh{收} (shōu, 1er) / \zh{手} (shǒu, 3e) et \zh{支} (zhī, 1er) / \zh{只} (zhǐ, 3e).} $\rightarrow$ \textit{Ex: *\zh{我手入不够} au lieu de \zh{我收入不够}.} \textbf{Drill :} Paires minimales quotidiennes.
    \item \textbf{Ordre des mots : Place du complément de lieu \zh{在...上} / \zh{在...里}.} $\rightarrow$ \textit{Erreur : *\zh{我花钱在吃饭上}. Correct : \zh{我在吃饭上花钱} (Sujet + \zh{在} Lieu + Verbe + Objet/Quantité).}
    \item \textbf{Mots-outils \zh{了} (le) : Double emploi ou oubli.} $\rightarrow$ \textit{Erreur : *\zh{我花了两百块了} (double \zh{le} aspectuel + modal). Correct : \zh{我花了两百块} (action terminée) OU \zh{我花两百块了} (nouvel état « j'en suis à 200 dépensés »).}
    \item \textbf{Classificateurs : Oubli ou mauvais choix pour argent/dépenses.} $\rightarrow$ \textit{Erreur : *\zh{三个钱}, *\zh{一张预算}. Correct : \zh{一笔收入/支出} (poste), \zh{一项预算} (projet), \zh{十块钱} (montant).}
    \item \textbf{Calque du français : « Je veux économiser de l'argent » $\neq$ *\zh{我想省钱} (je veux être économe/radin).} $\rightarrow$ \textit{Correct : \zh{我想存钱} (cúnqián -- mettre de côté/épargner) ou \zh{我想节约} (jiéyuē -- faire des économies/réduire dépenses). \zh{省} shěng = économiser (verbe action) ; \zh{存} cún = mettre en réserve/banque.}
\end{enumerate}
\cle{Astuce orale :} Si vous séchez sur un mot (ex: « overdraft »), paraphrasez : \zh{银行让我花超过我的钱} (La banque me laisse dépenser plus que mon argent) = \zh{透支} (tòuzhī). Montrez votre compétence stratégique !
\end{attention}

\newpage

\coursec{Exercices}

\courssub{Je vérifie mes connaissances}

\begin{exercice}[QCM : le vocabulaire du budget]
Pour chaque question, entoure la bonne réponse.
\begin{enumerate}
\item \zh{预算} (yùsuàn) signifie :
\begin{itemize}
\item[A.] le salaire
\item[B.] le budget
\item[C.] la dette
\end{itemize}
\item Pour dire « économiser de l'argent », on dit :
\begin{itemize}
\item[A.] \zh{花钱} (huā qián)
\item[B.] \zh{省钱} (shěng qián)
\item[C.] \zh{借钱} (jiè qián)
\end{itemize}
\item \zh{零花钱} (línghuāqián) désigne :
\begin{itemize}
\item[A.] l'argent de poche
\item[B.] le prix du transport
\item[C.] les frais de scolarité
\end{itemize}
\item Quelle phrase signifie « Je ne suis pas d'accord » ?
\begin{itemize}
\item[A.] \zh{我同意。} (Wǒ tóngyì.)
\item[B.] \zh{我觉得很好。} (Wǒ juéde hěn hǎo.)
\item[C.] \zh{我不同意。} (Wǒ bù tóngyì.)
\end{itemize}
\end{enumerate}
\end{exercice}

\begin{exercice}[Vrai ou faux, en justifiant]
Dis si chaque affirmation est vraie (V) ou fausse (F), puis justifie ta réponse en une phrase en français.
\begin{enumerate}
\item \zh{存} (cún) signifie « dépenser ».
\item Dans la phrase \zh{我每个月存一点儿钱。}, le mot \zh{一点儿} signifie « un peu ».
\item Pour relancer la parole dans un débat, on peut dire \zh{你呢？} (Nǐ ne ?).
\item \zh{太贵了！} (Tài guì le !) signifie « C'est trop bon marché ! ».
\end{enumerate}
\end{exercice}

\begin{exercice}[Vocabulaire : complète le tableau]
Complète le tableau suivant en écrivant le caractère chinois ou la traduction française qui manque.
\begin{center}
\begin{tabularx}{\linewidth}{|X|X|X|}
\hline
\rowcolor{popblueL} Caractères & Pinyin & Traduction \\
\hline
钱 & qián & \ldots \\
\hline
\ldots & huā & dépenser \\
\hline
存 & cún & \ldots \\
\hline
\ldots & piányi & bon marché \\
\hline
贵 & guì & \ldots \\
\hline
\ldots & gōngzī & salaire \\
\hline
\end{tabularx}
\end{center}
\end{exercice}

\begin{exercice}[Pinyin : associe les caractères à leur prononciation]
Associe chaque mot de la colonne A à son pinyin dans la colonne B. Attention aux tons !
\begin{center}
\begin{tabularx}{\linewidth}{|X|X|}
\hline
\rowcolor{popblueL} Colonne A (caractères) & Colonne B (pinyin) \\
\hline
1. 买 & a. mài \\
\hline
2. 卖 & b. mǎi \\
\hline
3. 省 & c. shēng \\
\hline
4. 生 & d. shěng \\
\hline
5. 预算 & e. jiàqián \\
\hline
6. 价钱 & f. yùsuàn \\
\hline
\end{tabularx}
\end{center}
\end{exercice}

\courssub{Je m'entraîne}

\begin{exercice}[Traduction : du français au chinois]
Traduis les phrases suivantes en chinois (caractères obligatoires ; ajoute le pinyin).
\begin{enumerate}
\item Je dépense mon argent pour le transport.
\item Tu devrais économiser un peu.
\item Je ne suis pas d'accord.
\item Chaque mois, je mets un peu d'argent de côté.
\item Ce sac est trop cher, je ne l'achète pas.
\end{enumerate}
\end{exercice}

\begin{exercice}[Dialogue à trous : les formules fonctionnelles]
Complète le dialogue suivant avec les formules de la boîte : \zh{我觉得} (wǒ juéde), \zh{我同意} (wǒ tóngyì), \zh{你呢？} (Nǐ ne ?), \zh{所以} (suǒyǐ). Puis joue ce dialogue en binôme en soignant les tons.
\begin{textebox}[Dialogue à compléter : 零花钱的问题]
A：你觉得中学生应该有零花钱吗？\\
\emph{Nǐ juéde zhōngxuéshēng yīnggāi yǒu línghuāqián ma ?}\\
« Penses-tu que les collégiens et lycéens devraient avoir de l'argent de poche ? »\\
B：\ldots\ldots 应该有，\ldots\ldots 不能太多。\\
« \ldots\ldots\ il devrait y en avoir, \ldots\ldots\ pas trop. »\\
A：\ldots\ldots ！太多钱对学生不好。\\
« \ldots\ldots\ ! Trop d'argent n'est pas bon pour les élèves. »\\
B：\ldots\ldots ？你每个月有多少零花钱？\\
« Et \ldots\ldots\ ? Combien d'argent de poche as-tu chaque mois ? »
\end{textebox}
\end{exercice}

\begin{exercice}[Remets les mots dans l'ordre]
Remets les mots dans le bon ordre pour former des phrases correctes, puis traduis-les en français.
\begin{enumerate}
\item \zh{钱 / 我 / 一点儿 / 每个月 / 存}
\item \zh{应该 / 你 / 省 / 一点儿 / 钱}
\item \zh{因为 / 太贵 / 所以 / 我 / 不买 / 了}
\item \zh{先 / 我 / 买 / 本子 / 然后 / 买 / 笔}
\end{enumerate}
\end{exercice}

\begin{exercice}[Texte à trous : les mots-outils]
Complète le texte suivant avec les mots de la boîte : \zh{的} (de), \zh{了} (le), \zh{因为} (yīnwèi), \zh{所以} (suǒyǐ), \zh{先} (xiān), \zh{然后} (ránhòu).
\begin{textebox}[Texte à compléter : 我的预算]
我每个月有五千法郎\ldots\ldots 零花钱。我\ldots\ldots 买车票，\ldots\ldots 买学习用品。\ldots\ldots 我想买一双新鞋，\ldots\ldots 我这个月存\ldots\ldots 两千法郎。\\
\emph{Wǒ měi ge yuè yǒu wǔqiān fǎláng \ldots\ldots línghuāqián. Wǒ \ldots\ldots mǎi chēpiào, \ldots\ldots mǎi xuéxí yòngpǐn. \ldots\ldots wǒ xiǎng mǎi yì shuāng xīn xié, \ldots\ldots wǒ zhè ge yuè cún \ldots\ldots liǎngqiān fǎláng.}\\
« Chaque mois, j'ai cinq mille francs \ldots\ldots argent de poche. \ldots\ldots j'achète des tickets de transport, \ldots\ldots des fournitures scolaires. \ldots\ldots je veux acheter une nouvelle paire de chaussures, \ldots\ldots ce mois-ci j'ai mis \ldots\ldots deux mille francs de côté. »
\end{textebox}
\end{exercice}

\courssub{Je me prépare au Bac}

\begin{exercice}[Type Bac — Compréhension de l'écrit et questions de langue]
Lis le texte suivant, puis réponds aux questions.
\begin{textebox}[Texte : 小王的零花钱]
小王是雅温得一所中学的学生。他每个月有八千法郎的零花钱。以前，他一个星期就把钱花完了：买零食、买饮料、请朋友吃饭。到了月底，他连买车票的钱都没有了。\\
后来，老师告诉他：「你应该先做一个预算，然后按计划花钱。」现在，小王先买学习用品和车票，然后存两千法郎，剩下的钱才买零食。\\
小王说：「因为我会管理我的钱，所以我再也不担心月底没有钱了。」\\
\emph{Xiǎo Wáng shì Yǎwēndé yì suǒ zhōngxué de xuésheng. Tā měi ge yuè yǒu bāqiān fǎláng de línghuāqián. Yǐqián, tā yí ge xīngqī jiù bǎ qián huā wán le : mǎi língshí, mǎi yǐnliào, qǐng péngyou chīfàn. Dào le yuèdǐ, tā lián mǎi chēpiào de qián dōu méiyǒu le.}\\
\emph{Hòulái, lǎoshī gàosu tā : « Nǐ yīnggāi xiān zuò yí ge yùsuàn, ránhòu àn jìhuà huā qián. » Xiànzài, Xiǎo Wáng xiān mǎi xuéxí yòngpǐn hé chēpiào, ránhòu cún liǎngqiān fǎláng, shèngxià de qián cái mǎi língshí.}\\
\emph{Xiǎo Wáng shuō : « Yīnwèi wǒ huì guǎnlǐ wǒ de qián, suǒyǐ wǒ zài yě bù dānxīn yuèdǐ méiyǒu qián le. »}\\
« Petit Wang est élève dans un lycée de Yaoundé. Chaque mois, il a huit mille francs d'argent de poche. Avant, il dépensait tout en une semaine : snacks, boissons, repas offerts aux amis. À la fin du mois, il n'avait même plus de quoi acheter un ticket de transport. Plus tard, son professeur lui a dit : "Tu devrais d'abord faire un budget, puis dépenser selon ton plan." Maintenant, Petit Wang achète d'abord les fournitures scolaires et les tickets, puis met deux mille francs de côté ; ce n'est qu'avec le reste qu'il achète des snacks. Petit Wang dit : "Parce que je sais gérer mon argent, je ne m'inquiète plus de ne plus avoir d'argent à la fin du mois." »
\end{textebox}
\textbf{Questions de compréhension} (réponds en français) :
\begin{enumerate}
\item Combien d'argent de poche Petit Wang reçoit-il chaque mois ?
\item Comment dépensait-il son argent avant ? Cite deux exemples du texte.
\item Quel conseil le professeur lui a-t-il donné ?
\item Que fait Petit Wang maintenant avec son argent ? Décris les trois étapes.
\end{enumerate}
\textbf{Questions de langue} :
\begin{enumerate}
\item Relève dans le texte la structure \zh{先……然后……} et explique son emploi.
\item Trouve dans le texte une phrase avec \zh{因为……所以……} et traduis-la.
\item Explique l'emploi de \zh{了} dans la phrase \zh{他连买车票的钱都没有了。}
\end{enumerate}
\end{exercice}

\begin{exercice}[Type Bac — Thème et version]
\textbf{Thème.} Traduis les phrases suivantes en chinois (caractères obligatoires) :
\begin{enumerate}
\item Tu ne devrais pas dépenser tout ton argent de poche en une semaine.
\item Parce que le transport est cher, je fais d'abord un budget chaque mois.
\item Je pense que les lycéens devraient apprendre à gérer leur argent.
\end{enumerate}
\textbf{Version.} Traduis le passage suivant en français :
\begin{textebox}[Texte à traduire]
我姐姐在大学学习。她的钱不多，所以她每个月都做一个预算。她先付房租，然后买吃的，最后才买衣服。\\
\emph{Wǒ jiějie zài dàxué xuéxí. Tā de qián bù duō, suǒyǐ tā měi ge yuè dōu zuò yí ge yùsuàn. Tā xiān fù fángzū, ránhòu mǎi chī de, zuìhòu cái mǎi yīfu.}
\end{textebox}
\end{exercice}

\begin{exercice}[Type Bac — Expression écrite et préparation du débat oral]
\textbf{Sujet.} Tu prépares un débat en classe sur la question : \zh{中学生应该有零花钱吗？} (Les lycéens devraient-ils avoir de l'argent de poche ?)
\begin{enumerate}
\item \textbf{À l'écrit :} rédige un texte de \textbf{80 à 100 caractères chinois} dans lequel tu donnes ton avis. Tu utiliseras obligatoirement :
\begin{itemize}
\item au moins deux formules fonctionnelles parmi : \zh{我觉得……}, \zh{我同意}, \zh{我不同意}, \zh{我认为……} ;
\item au moins une fois la structure \zh{因为……所以……} ;
\item au moins une fois les connecteurs \zh{先……然后……}.
\end{itemize}
\item \textbf{À l'oral :} présente ensuite ton point de vue devant la classe en \textbf{dix phrases}, en prononçant correctement les tons, et réponds aux questions de deux camarades en utilisant \zh{你呢？} pour relancer la parole.
\end{enumerate}
\end{exercice}

\newpage

\coursec{Corrigés des exercices}

\begin{corrige}[Exercice 1 — QCM : le vocabulaire du budget]
\begin{enumerate}
\item \textbf{B} — \zh{预算} (yùsuàn) signifie « le budget ». Le salaire se dit \zh{工资} (gōngzī) et la dette \zh{债} (zhài).
\item \textbf{B} — « économiser de l'argent » se dit \zh{省钱} (shěng qián). \zh{花钱} (huā qián) signifie « dépenser de l'argent » et \zh{借钱} (jiè qián) « emprunter de l'argent ».
\item \textbf{A} — \zh{零花钱} (línghuāqián) désigne « l'argent de poche ». Littéralement : « argent pour les petites dépenses » (\zh{零} = petit, fractionné ; \zh{花} = dépenser).
\item \textbf{C} — \zh{我不同意。} (Wǒ bù tóngyì.) signifie « Je ne suis pas d'accord ». La négation se fait avec \zh{不} placé devant le verbe \zh{同意} (tóngyì, être d'accord).
\end{enumerate}
\end{corrige}

\begin{corrige}[Exercice 2 — Vrai ou faux, en justifiant]
\begin{enumerate}
\item \textbf{Faux.} \zh{存} (cún) signifie « mettre de côté, épargner, déposer (à la banque) ». « Dépenser » se dit \zh{花} (huā).
\item \textbf{Vrai.} \zh{一点儿} (yìdiǎnr) signifie « un peu ». Dans \zh{我每个月存一点儿钱。} (Wǒ měi ge yuè cún yìdiǎnr qián.), la phrase signifie « Chaque mois, je mets un peu d'argent de côté ».
\item \textbf{Vrai.} \zh{你呢？} (Nǐ ne ?) est une formule de relance qui signifie « Et toi ? ». La particule \zh{呢} reprend le sujet du tour précédent pour interroger l'interlocuteur.
\item \textbf{Faux.} \zh{太贵了！} (Tài guì le !) signifie « C'est trop cher ! ». « Trop bon marché » se dirait \zh{太便宜了！} (Tài piányi le !). Attention à ne pas confondre \zh{贵} (guì, cher) et \zh{便宜} (piányi, bon marché).
\end{enumerate}
\end{corrige}

\begin{corrige}[Exercice 3 — Vocabulaire : complète le tableau]
\begin{center}
\begin{tabularx}{\linewidth}{|X|X|X|}
\hline
\rowcolor{popblueL} Caractères & Pinyin & Traduction \\
\hline
钱 & qián & l'argent \\
\hline
花 & huā & dépenser \\
\hline
存 & cún & mettre de côté, épargner \\
\hline
便宜 & piányi & bon marché \\
\hline
贵 & guì & cher \\
\hline
工资 & gōngzī & salaire \\
\hline
\end{tabularx}
\end{center}
\textbf{Points de vigilance :} ne pas confondre \zh{花} (huā, dépenser ; aussi « fleur ») et \zh{华} (huá, comme dans \zh{中华}) ; \zh{便宜} se prononce piányi (et non biànyí) dans le sens de « bon marché » — la lecture biànyí existe dans \zh{方便}… non, attention : \zh{方便} se lit fāngbiàn ; la lecture biànyí de \zh{便宜} est littéraire et ne s'emploie pas au niveau du lycée. Retenir : « bon marché » = piányi.
\end{corrige}

\begin{corrige}[Exercice 4 — Pinyin : associe les caractères à leur prononciation]
\begin{center}
\begin{tabularx}{\linewidth}{|X|X|}
\hline
\rowcolor{popblueL} Colonne A (caractères) & Colonne B (pinyin) \\
\hline
1. 买 & b. mǎi (3\textsuperscript{e} ton) — acheter \\
\hline
2. 卖 & a. mài (4\textsuperscript{e} ton) — vendre \\
\hline
3. 省 & d. shěng (3\textsuperscript{e} ton) — économiser \\
\hline
4. 生 & c. shēng (1\textsuperscript{er} ton) — naître, brut \\
\hline
5. 预算 & f. yùsuàn — budget \\
\hline
6. 价钱 & e. jiàqián — prix \\
\hline
\end{tabularx}
\end{center}
\textbf{Points de vigilance :} la paire \zh{买} (mǎi, acheter) / \zh{卖} (mài, vendre) est un piège classique pour les francophones : seuls le ton et un trait supplémentaire sur \zh{卖} les distinguent. De même, \zh{省} (shěng) et \zh{生} (shēng) ne diffèrent que par le ton. À l'oral, exagérer le contraste : 3\textsuperscript{e} ton qui descend puis remonte pour \zh{买}, 4\textsuperscript{e} ton sec et descendant pour \zh{卖}.
\end{corrige}

\begin{corrige}[Exercice 5 — Traduction : du français au chinois]
\begin{enumerate}
\item \zh{我把钱花在交通上。} ou plus simplement \zh{我花钱买车票。}\\
\emph{Wǒ bǎ qián huā zài jiāotōng shàng. / Wǒ huā qián mǎi chēpiào.}\\
« Je dépense mon argent pour le transport. » (Dans la première version, la construction en \zh{把} met l'objet \zh{钱} avant le verbe ; la structure est \zh{把钱花在……上}.)
\item \zh{你应该省一点儿钱。}\\
\emph{Nǐ yīnggāi shěng yìdiǎnr qián.}\\
« Tu devrais économiser un peu. » (Le modal \zh{应该} se place avant le verbe.)
\item \zh{我不同意。}\\
\emph{Wǒ bù tóngyì.}\\
« Je ne suis pas d'accord. »
\item \zh{我每个月存一点儿钱。}\\
\emph{Wǒ měi ge yuè cún yìdiǎnr qián.}\\
« Chaque mois, je mets un peu d'argent de côté. » (Le complément de temps \zh{每个月} se place avant le verbe.)
\item \zh{这个包太贵了，我不买。}\\
\emph{Zhè ge bāo tài guì le, wǒ bù mǎi.}\\
« Ce sac est trop cher, je ne l'achète pas. » (Classificateur \zh{个} devant \zh{包} ; structure d'exclamation \zh{太……了}.)
\end{enumerate}
\end{corrige}

\begin{corrige}[Exercice 6 — Dialogue à trous : les formules fonctionnelles]
Dialogue complété :
\begin{itemize}
\item B：\zh{我觉得应该有，但是不能太多。}\\
\emph{Wǒ juéde yīnggāi yǒu, dànshì bù néng tài duō.}\\
« Je pense qu'il devrait y en avoir, mais pas trop. »\\
Première case : \zh{我觉得} (wǒ juéde, « je pense que »), formule d'opinion. Deuxième case : le connecteur d'opposition \zh{但是} (dànshì, « mais »), car la seconde partie limite la première (« il en faut, \emph{mais} pas trop »).
\item A：\zh{我同意！太多钱对学生不好。}\\
\emph{Wǒ tóngyì ! Tài duō qián duì xuésheng bù hǎo.}\\
« Je suis d'accord ! Trop d'argent n'est pas bon pour les élèves. » (Formule d'accord : \zh{我同意}.)
\item B：\zh{你呢？你每个月有多少零花钱？}\\
\emph{Nǐ ne ? Nǐ měi ge yuè yǒu duōshao línghuāqián ?}\\
« Et toi ? Combien d'argent de poche as-tu chaque mois ? » (Formule de relance : \zh{你呢？}.)
\end{itemize}
\textbf{Remarque :} le mot \zh{所以} (suǒyǐ, « donc ») de la boîte ne convient pas à la deuxième case : \zh{应该有，所以不能太多} (« il en faut, donc pas trop ») n'est pas logique. La bonne solution est bien \zh{但是}. Pour la prestation orale : soigner le 3\textsuperscript{e} ton de \zh{我} (wǒ) et le ton montant de \zh{同意} (tóngyì, 2\textsuperscript{e} + 4\textsuperscript{e} ton).
\end{corrige}

\begin{corrige}[Exercice 7 — Remets les mots dans l'ordre]
\begin{enumerate}
\item \zh{我每个月存一点儿钱。}\\
\emph{Wǒ měi ge yuè cún yìdiǎnr qián.}\\
« Chaque mois, je mets un peu d'argent de côté. » (Ordre : sujet + complément de temps + verbe + objet.)
\item \zh{你应该省一点儿钱。}\\
\emph{Nǐ yīnggāi shěng yìdiǎnr qián.}\\
« Tu devrais économiser un peu d'argent. » (Le modal \zh{应该} précède le verbe \zh{省}.)
\item \zh{因为太贵了，所以我不买。}\\
\emph{Yīnwèi tài guì le, suǒyǐ wǒ bù mǎi.}\\
« Parce que c'est trop cher, je ne l'achète pas. » (En chinois, \zh{因为} et \zh{所以} peuvent se combiner dans la même phrase, contrairement au français où « parce que » et « donc » s'excluent.)
\item \zh{我先买本子，然后买笔。}\\
\emph{Wǒ xiān mǎi běnzi, ránhòu mǎi bǐ.}\\
« J'achète d'abord les cahiers, puis les stylos. » (\zh{先} se place avant le verbe ; \zh{然后} introduit la deuxième action.)
\end{enumerate}
\end{corrige}

\begin{corrige}[Exercice 8 — Texte à trous : les mots-outils]
Texte complété :

我每个月有五千法郎\zh{的}零花钱。我\zh{先}买车票，\zh{然后}买学习用品。\zh{因为}我想买一双新鞋，\zh{所以}我这个月存\zh{了}两千法郎。

\emph{Wǒ měi ge yuè yǒu wǔqiān fǎláng de línghuāqián. Wǒ xiān mǎi chēpiào, ránhòu mǎi xuéxí yòngpǐn. Yīnwèi wǒ xiǎng mǎi yì shuāng xīn xié, suǒyǐ wǒ zhè ge yuè cún le liǎngqiān fǎláng.}

« Chaque mois, j'ai cinq mille francs d'argent de poche. J'achète d'abord les tickets de transport, puis les fournitures scolaires. Parce que je veux acheter une nouvelle paire de chaussures, ce mois-ci j'ai mis deux mille francs de côté. »

\textbf{Justifications :}
\begin{itemize}
\item \zh{的} : marque du complément du nom — « cinq mille francs \emph{d'}argent de poche ».
\item \zh{先……然后……} : chronologie des actions (« d'abord… puis… »).
\item \zh{因为……所以……} : cause et conséquence.
\item \zh{了} : aspect accompli — l'action d'épargner est réalisée (« j'ai mis de côté »).
\end{itemize}
\textbf{Point de vigilance :} devant \zh{两}, on dit \zh{两千} (liǎngqiān) et non \zh{二千} : pour les quantités, on emploie \zh{两} devant les classificateurs et les grands nombres (\zh{两千}, \zh{两万}).
\end{corrige}

\begin{corrige}[Exercice 9 — Type Bac — Compréhension de l'écrit et questions de langue]
\textbf{Barème indicatif :} compréhension 8 points (2 par question) ; questions de langue 6 points (2 par question).

\textbf{Questions de compréhension :}
\begin{enumerate}
\item Petit Wang reçoit \textbf{huit mille francs} d'argent de poche chaque mois (\zh{八千法郎}).
\item Avant, il dépensait tout son argent en une semaine, par exemple en achetant des \textbf{snacks} (\zh{买零食}) et des \textbf{boissons} (\zh{买饮料}), ou en \textbf{offrant des repas à ses amis} (\zh{请朋友吃饭}).
\item Le professeur lui a conseillé de \textbf{faire d'abord un budget, puis de dépenser selon son plan} (\zh{先做一个预算，然后按计划花钱}).
\item Maintenant, Petit Wang : (1) achète d'abord les fournitures scolaires et les tickets de transport ; (2) met ensuite deux mille francs de côté ; (3) n'achète des snacks qu'avec l'argent restant.
\end{enumerate}
\textbf{Questions de langue :}
\begin{enumerate}
\item Structure relevée : \zh{小王先买学习用品和车票，然后存两千法郎}. La structure \zh{先……然后……} (xiān… ránhòu…) exprime la chronologie : « d'abord… ensuite… ». \zh{先} se place avant le premier verbe, \zh{然后} introduit la seconde proposition.
\item Phrase relevée : \zh{因为我会管理我的钱，所以我再也不担心月底没有钱了。}\\
\emph{Yīnwèi wǒ huì guǎnlǐ wǒ de qián, suǒyǐ wǒ zài yě bù dānxīn yuèdǐ méiyǒu qián le.}\\
« Parce que je sais gérer mon argent, je ne m'inquiète plus de ne plus avoir d'argent à la fin du mois. » (\zh{再也不} = « ne… plus jamais ».)
\item Dans \zh{他连买车票的钱都没有了。}, \zh{了} marque un \textbf{changement de situation} : auparavant il avait de l'argent, désormais il n'en a plus. C'est le \zh{了} de nouvelle situation, fréquent avec \zh{没有} (\zh{没有了} = « ne plus avoir »).
\end{enumerate}
\textbf{Points de vigilance :} ne pas traduire \zh{连……都……} mot à mot ; cette structure insiste (« même… ») : « il n'avait même plus l'argent du ticket ».
\end{corrige}

\begin{corrige}[Exercice 10 — Type Bac — Thème et version]
\textbf{Barème indicatif :} thème 9 points (3 par phrase) ; version 6 points.

\textbf{Thème :}
\begin{enumerate}
\item \zh{你不应该一个星期就把零花钱花完。}\\
\emph{Nǐ bù yīnggāi yí ge xīngqī jiù bǎ línghuāqián huā wán.}\\
« Tu ne devrais pas dépenser tout ton argent de poche en une semaine. » (Négation du modal : \zh{不应该} ; \zh{就} insiste sur la rapidité ; \zh{花完} = « finir de dépenser », avec le complément de résultat \zh{完}.)
\item \zh{因为交通很贵，所以我每个月先做一个预算。}\\
\emph{Yīnwèi jiāotōng hěn guì, suǒyǐ wǒ měi ge yuè xiān zuò yí ge yùsuàn.}\\
« Parce que le transport est cher, je fais d'abord un budget chaque mois. »
\item \zh{我觉得中学生应该学会管理自己的钱。}\\
\emph{Wǒ juéde zhōngxuéshēng yīnggāi xuéhuì guǎnlǐ zìjǐ de qián.}\\
« Je pense que les lycéens devraient apprendre à gérer leur argent. » (\zh{学会} = « apprendre à, acquérir la compétence » ; \zh{自己的} = « son propre ».)
\end{enumerate}
\textbf{Version :}

« Ma grande sœur étudie à l'université. Elle n'a pas beaucoup d'argent, alors elle fait un budget chaque mois. Elle paie d'abord le loyer, puis achète de la nourriture, et ce n'est qu'enfin qu'elle achète des vêtements. »

\textbf{Points de vigilance :} \zh{最后才} (zuìhòu cái) exprime une restriction temporelle : « ce n'est qu'à la fin que… » ; \zh{买吃的} signifie littéralement « acheter de quoi manger » ; ne pas oublier que \zh{都} dans \zh{每个月都做} marque l'habitude (« chaque mois sans exception »).
\end{corrige}

\begin{corrige}[Exercice 11 — Type Bac — Expression écrite et préparation du débat oral]
\textbf{Barème indicatif :} respect des contraintes (formules, \zh{因为……所以……}, \zh{先……然后……}) 6 points ; correction grammaticale et caractères 6 points ; richesse du vocabulaire 4 points ; prestation orale (tons, fluidité, relance avec \zh{你呢？}) 4 points.

\textbf{Proposition de rédaction modèle (environ 90 caractères) :}

\zh{我觉得中学生应该有零花钱，但是不能太多。因为我们要买车票和学习用品，所以每个月有一点儿钱很方便。我认为我们应该先做一个预算，然后按计划花钱。我不同意随便花钱：会管理钱的学生，以后才会管理自己的生活。}

\emph{Wǒ juéde zhōngxuéshēng yīnggāi yǒu línghuāqián, dànshì bù néng tài duō. Yīnwèi wǒmen yào mǎi chēpiào hé xuéxí yòngpǐn, suǒyǐ měi ge yuè yǒu yìdiǎnr qián hěn fāngbiàn. Wǒ rènwéi wǒmen yīnggāi xiān zuò yí ge yùsuàn, ránhòu àn jìhuà huā qián. Wǒ bù tóngyì suíbiàn huā qián : huì guǎnlǐ qián de xuésheng, yǐhòu cái huì guǎnlǐ zìjǐ de shēnghuó.}

« Je pense que les lycéens devraient avoir de l'argent de poche, mais pas trop. Parce que nous devons acheter les tickets de transport et les fournitures scolaires, avoir un peu d'argent chaque mois est très pratique. J'estime que nous devrions d'abord faire un budget, puis dépenser selon le plan. Je ne suis pas d'accord pour dépenser n'importe comment : un élève qui sait gérer son argent saura plus tard gérer sa propre vie. »

\textbf{Vérification des contraintes :} formules fonctionnelles \zh{我觉得}, \zh{我认为}, \zh{我不同意} (trois, minimum deux) ; \zh{因为……所以……} présent ; \zh{先……然后……} présent.

\textbf{Conseils pour l'oral :}
\begin{itemize}
\item Soigner le 3\textsuperscript{e} ton de \zh{我} (wǒ) et de \zh{有} (yǒu) ; bien détacher les deux 4\textsuperscript{es} tons de \zh{预算} (yùsuàn).
\item Marquer une courte pause après la proposition en \zh{因为}… avant la conséquence introduite par \zh{所以}.
\item Pour relancer : \zh{你呢？你觉得零花钱应该有多少？} (Nǐ ne ? Nǐ juéde línghuāqián yīnggāi yǒu duōshao ?) — « Et toi ? Combien d'argent de poche devrait-on avoir, selon toi ? »
\end{itemize}
\textbf{Points de vigilance :} ne pas écrire \zh{因为} sans \zh{所以} dans une longue phrase écrite ; ne pas placer \zh{先} après le verbe ; compter les caractères pour rester entre 80 et 100.
\end{corrige}

\newpage

\coursec{Fiche bilan}

\begin{popfigure}
\begin{tikzpicture}[mindmap, grow cyclic, every node/.style={concept, font=\small},
  root concept/.append style={concept color=popblue, font=\bfseries},
  level 1/.append style={level distance=3.4cm, sibling angle=60},
  level 1 concept/.append style={font=\footnotesize}]
\node[root concept]{Gestion du budget \zh{管理预算}}
  child[concept color=poporange]{ node {Lexique \zh{预算 收入 支出 节省}} }
  child[concept color=popgreen]{ node {Dialogues modèles \zh{对话}} }
  child[concept color=poppurple]{ node {Formules : présenter, opiner, relancer} }
  child[concept color=poppink]{ node {Tons à travailler} }
  child[concept color=popgold]{ node {Jeux de rôle et débat} }
  child[concept color=popblue!60]{ node {Exposé simple et évaluation} };
\end{tikzpicture}
\legende{Carte mentale de la leçon : parler de la gestion du budget en chinois.}
\end{popfigure}

\begin{bilan}
\textbf{1. Lexique clé du budget}
\begin{center}
\begin{tabularx}{\linewidth}{|X|X|X|X|}\hline
\rowcolor{popblueL} Caractères & Pinyin & Nature & Traduction \\ \hline
预算 & yùsuàn & n. & budget \\ \hline
收入 & shōurù & n. & revenu \\ \hline
支出 & zhīchū & n. / v. & dépense ; dépenser \\ \hline
节省 & jiéshěng & v. & économiser \\ \hline
存钱 & cún qián & v. & mettre de l'argent de côté \\ \hline
花钱 & huā qián & v. & dépenser de l'argent \\ \hline
够 & gòu & adj. / v. & suffire \\ \hline
浪费 & làngfèi & v. / adj. & gaspiller \\ \hline
计划 & jìhuà & n. / v. & plan ; planifier \\ \hline
需要 & xūyào & v. & avoir besoin de \\ \hline
\end{tabularx}
\end{center}

\textbf{2. Structures à connaître par cœur}
\begin{center}
\begin{tabularx}{\linewidth}{|X|X|X|}\hline
\rowcolor{popblueL} Structure & Exemple & Traduction \\ \hline
Sujet + 觉得 + phrase & 我觉得应该先存钱。Wǒ juéde yīnggāi xiān cún qián. & Je pense qu'il faut d'abord économiser. \\ \hline
应该 / 不应该 + verbe & 我们不应该浪费钱。Wǒmen bù yīnggāi làngfèi qián. & Nous ne devrions pas gaspiller l'argent. \\ \hline
先 … 然后 … & 先算收入，然后算支出。Xiān suàn shōurù, ránhòu suàn zhīchū. & D'abord calculer les revenus, puis les dépenses. \\ \hline
A + 比 + B + adj. & 存钱比花钱难。Cún qián bǐ huā qián nán. & Économiser est plus difficile que dépenser. \\ \hline
如果 … 就 … & 如果钱不够，就要节省。Rúguǒ qián bù gòu, jiù yào jiéshěng. & Si l'argent ne suffit pas, il faut économiser. \\ \hline
越 … 越 … & 越早计划越好。Yuè zǎo jìhuà yuè hǎo. & Plus on planifie tôt, mieux c'est. \\ \hline
\end{tabularx}
\end{center}

\textbf{3. Formules fonctionnelles pour le débat}
\begin{itemize}
\item Présenter : \zh{我今天要谈的是……} (Wǒ jīntiān yào tán de shì…) « Aujourd'hui, je vais parler de… »
\item Donner son avis : \zh{我觉得…… / 在我看来……} (Wǒ juéde… / Zài wǒ kànlái…) « Je pense que… / À mon avis… »
\item Être d'accord : \zh{我同意你的看法。} (Wǒ tóngyì nǐ de kànfǎ.) « Je suis d'accord avec toi. »
\item Ne pas être d'accord : \zh{我不太同意，因为……} (Wǒ bú tài tóngyì, yīnwèi…) « Je ne suis pas tout à fait d'accord, parce que… »
\item Relancer : \zh{你怎么看？/ 你呢？} (Nǐ zěnme kàn ? / Nǐ ne ?) « Qu'en penses-tu ? / Et toi ? »
\item Conclure : \zh{总之，……} (Zǒngzhī,…) « Bref, … »
\end{itemize}

\textbf{4. Méthode express pour l'exposé oral}
\begin{enumerate}
\item Annoncer le sujet (\zh{我今天要谈的是家庭预算。}).
\item Donner deux ou trois idées avec \zh{先……然后……最后……}.
\item Illustrer avec un exemple concret (marché de Yaoundé, argent de poche).
\item Donner son avis avec \zh{我觉得……}.
\item Conclure avec \zh{总之……} et remercier : \zh{谢谢大家！} (Xièxie dàjiā !).
\end{enumerate}

\textbf{5. Tons à travailler}
\begin{itemize}
\item \zh{预算} yùsuàn : 4\textsuperscript{e} + 4\textsuperscript{e} ton, deux tons descendants nets.
\item \zh{节省} jiéshěng : 2\textsuperscript{e} + 3\textsuperscript{e} ton, monter puis descendre-remonter.
\item \zh{收入} shōurù : 1\textsuperscript{er} + 4\textsuperscript{e} ton, haut et plat puis descendant.
\item Attention à \zh{买} mǎi (acheter, 3\textsuperscript{e} ton) et \zh{卖} mài (vendre, 4\textsuperscript{e} ton) : ne pas confondre !
\end{itemize}
\end{bilan}

\begin{aretenir}[Checklist avant le Bac]
\begin{itemize}
\item[$\square$] Je connais par cœur les dix mots clés du budget (caractères, pinyin, tons).
\item[$\square$] Je sais présenter un sujet avec \zh{我今天要谈的是……}.
\item[$\square$] Je sais donner mon avis avec \zh{我觉得……} et \zh{在我看来……}.
\item[$\square$] Je sais exprimer l'accord et le désaccord poliment.
\item[$\square$] Je sais relancer le débat avec \zh{你怎么看？}.
\item[$\square$] Je maîtrise la comparaison avec \zh{比} (A + 比 + B + adjectif).
\item[$\square$] Je sais utiliser \zh{如果……就……} pour exprimer une condition.
\item[$\square$] Je sais organiser mon exposé avec \zh{先……然后……最后……}.
\item[$\square$] Je prononce correctement les tons de \zh{预算}, \zh{节省}, \zh{收入}.
\item[$\square$] Je ne confonds pas \zh{买} mǎi (acheter) et \zh{卖} mài (vendre).
\item[$\square$] Je sais conclure avec \zh{总之……} et remercier l'auditoire.
\item[$\square$] Je peux parler deux minutes sans lire mes notes.
\end{itemize}
\end{aretenir}

\findecours

\end{document}
